\documentclass[12pt]{report}

\usepackage{amssymb}
\usepackage{tikz}
\usepackage{graphicx}
\usepackage{multirow}
\usepackage{pgfplots}
\usepackage[dvipsnames, table, x11names, svgnames]{xcolor}
\usepackage{caption}
\usepackage{anyfontsize}
\usepackage{etoolbox}
\usepackage[italian]{babel}
\usepackage[pass]{geometry}
\usepackage{tcolorbox}
\usepackage{ragged2e}
\usepackage{listings}
\usepackage{ltablex}
\usepackage{environ}
\usepackage{titlesec}
\usepackage{tabularx}
\usetikzlibrary{positioning}
\usepackage[hidelinks]{hyperref}

\tcbuselibrary{theorems, listings, breakable, skins}

% Togliamo la fastidiosissima cosa che appena vado a capo mi mette lo skip
% come se fosse un paragrafo. Modifichiamo anche lo skip verticale.
\setlength{\parindent}{0pt}
\setlength{\parskip}{1ex}

% Modifichiamo lo stile del generico capitolo
\titleformat{\chapter}[display]
{\normalfont\huge\bfseries\linespread{1.3}\selectfont\raggedright}
{\chaptertitlename\ \thechapter}
{5pt}
{\Huge}

\titlespacing*{\chapter}{0pt}{0pt}{2pt}
\titlespacing*{\section}{0pt}{5pt}{3pt}
\titlespacing*{\subsection}{0pt}{5pt}{3pt}
\titlespacing*{\subsubsection}{0pt}{5pt}{3pt}

\newcommand{\generaFrontespizio}[4]{
    \begin{titlepage}
		
		\begin{center}
			
			\hbox{}
			\vfill
			
			{
				\color{BrickRed}
				\bfseries
				\fontsize{40pt}{40pt}\selectfont
				#1
				\par
			}
			
			\vspace{0.5cm}
			
			\vfill
			\hbox{}
			
			\begin{minipage}{1\textwidth}
				\fontsize{12pt}{14pt}\selectfont \textit{I seguenti appunti sono a cura di \textbf{Catalin Ceban} e si basano sul corso di "#1" erogato nell'A.A. #2 
				\ifstrequal{#4}{m}{ dal professore}{\ifstrequal{#4}{p}{ dai professori}{ dalla professoressa}} #3 per il C.d.L Informatica (33503)}
			\end{minipage}
			
			\vspace{0.5cm}
			{\large \today}
			
		\end{center}
		
	
	\end{titlepage}

	\newgeometry{
		includefoot,
		left=3.5cm,
		right=3.5cm,
		top=2.9cm,
		bottom=2cm
	}

	\tableofcontents
	\clearpage
	
}

\colorlet{coloreTeoremi}{DarkGoldenrod2}
\colorlet{coloreDefinizioni}{Coral3}
\colorlet{coloreProposizioni}{LightSalmon3}
\colorlet{coloreCorollari}{MediumOrchid3}
\colorlet{coloreOsservazioni}{Burlywood4}
\colorlet{coloreSfondoBox}{gray!5}

\newlength{\boxDistanzaVerticale}
\setlength{\boxDistanzaVerticale}{7pt}
\newlength{\boxDistanzaOrizzontale}
\setlength{\boxDistanzaOrizzontale}{10pt}
\newlength{\boxSpessoreRiga}
\setlength{\boxSpessoreRiga}{3pt}

\tcbset{
	stileBox/.style={
		blanker,
		borderline west={\boxSpessoreRiga}{0pt}{#1},
		left=\boxDistanzaOrizzontale,
		toptitle=\boxDistanzaVerticale, 
		bottom=\boxDistanzaVerticale,
		right=\boxDistanzaOrizzontale,
		pad after break=2mm,
		colback=coloreSfondoBox,
		colbacktitle=tcbcolback,
		coltitle=#1,
		fonttitle=\bfseries,
		enhanced,
		breakable
	}
}

\newtcbtheorem[number within=section]{teorema}{Teorema}{stileBox=coloreTeoremi}{th}

\newtcbtheorem[number within=section]{definizione}{Definizione}{stileBox=coloreDefinizioni}{def}

\newtcbtheorem[number within=section]{proposizione}{Proposizione}{stileBox=coloreProposizioni}{prop}

\newtcbtheorem[number within=section]{corollario}{Corollario}{stileBox=coloreCorollari}{cor}

\newtcbtheorem[number within=section]{osservazione}{Osservazione}{stileBox=coloreOsservazioni}{oss}

\newtcbtheorem[number within=section]{dimostrazione}{Dimostrazione}{stileBox=coloreOsservazioni}{oss}

\newtcbtheorem[number within=section]{algoritmo}{Spiegazione algoritmo}{stileBox=coloreProposizioni}{alg}

\newtcbtheorem[number within=section]{esercizio}{Esercizio}{stileBox=coloreTeoremi}{es}

\lstdefinestyle{stile-listings}{
	basicstyle=\ttfamily\small,
	inputencoding=utf8,
	extendedchars=true,
	literate=
	{à}{{\`a}}1
	{è}{{\`e}}1
	{é}{{\'e}}1
	{ì}{{\`i}}1
	{ò}{{\`ò}}1
	{ù}{{\`u}}1,
	keywordstyle=\color{blue!55!black}\bfseries,
	stringstyle=\color{Coral!55!black},
	numbers=left,
	numberstyle=\tiny\color{black!45},
	numbersep=8pt,
	backgroundcolor=\color{gray!15},
	frame=single,
	rulecolor=\color{gray!70!black},
	framesep=6pt,
	framextopmargin=4pt,
	framexbottommargin=4pt,
	xleftmargin=7pt,
	xrightmargin=7pt,
	breaklines=true,
	breakatwhitespace=false,
	showstringspaces=false,
	tabsize=2,
	captionpos=b,
	aboveskip=15pt,
	belowskip=8pt
}

\newcolumntype{C}{>{\Centering\arraybackslash}X}

\NewEnviron{tabella}[2][]{%
	\par\vspace{0.5em}%
	\begingroup
	\centering
	\renewcommand{\arraystretch}{1.5}%
	\begin{tabularx}{\textwidth}{#2}
		\hline
		\rowcolor{BrickRed!60!white}
		\BODY
		\hline
	\end{tabularx}%
	\if\relax\detokenize{#1}\relax
	% Nessuna caption fornita
	\else
	\vspace{0.5em}%
	\captionof{table}{#1}%
	\fi
	\par
	\endgroup
	\vspace{0.5em}%
}




\newcommand{\N}{\mathbb{N}}
\newcommand{\Z}{\mathbb{Z}}
\newcommand{\Q}{\mathbb{Q}}
\newcommand{\R}{\mathbb{R}}
\newcommand{\C}{\mathbb{C}}
\newcommand{\numeroromano}[1]
{\MakeUppercase{\romannumeral #1}}

\begin{document}
	
	\generaFrontespizio{Analisi \numeroromano{1}}{2025/2026}{Valeriano Aiello, Andrea Terracina e Flavia Lanzara}{p}
	
	\chapter{Insiemi numerici}
	Per introdurre lo studente allo studio della materia, si deve partire dalle basi. Ciò implica la definizione degli insiemi numerici di base e delle operazioni possibili in questi.
	
	\begin{description}
		
		\item [Numeri naturali] $\N$ = $\{0, 1, 2, 3, \dots\}$, su cui sono definite le seguenti operazioni:
		\begin{description}
			
			\item [Somma (+):] Presi $n, m \in \N$, viene definita la somma come la conosciamo dalle elementari, cioè $n + m$.
			
			\item [Prodotto ($\cdot$):] Presi $n, m \in \N$, definiamo il prodotto tra questi due numeri come $n \cdot m := \underbrace{n + n + \dots + n}_{m\text{ volte}}$.
			
		\end{description}
		
		Enunciamo di seguito alcune proprietà, singole o comuni, di queste operazioni:
		\begin{enumerate}
			\item \textbf{Proprietà commutativa}: per la somma equivale a dire che $n + m = m + n$, mentre per il prodotto $n \cdot m = m \cdot n$.
			\item \textbf{Elemento neutro}, cioè l'elemento che dato un numero $n$ restituisce $n$ stesso quando applicata l'operazione. Per la somma l'elemento neutro è $0$, dunque $n + 0 = n$, mentre per il prodotto è $1$, dunque $n \cdot 1 = n$.
			\item \textbf{Associatività}: per la somma vale che $n + (m + q) = (n + m) + q$, e analogamente per il prodotto $n \cdot (m \cdot q) = (n \cdot m) \cdot q$.
			\item \textbf{Distributività}: esclusiva per il prodotto, per questa proprietà vale che $n \cdot (m + q) = n\cdot m + n\cdot q$
		\end{enumerate}
	
		\item [Numeri interi] $\Z = \{\dots, -2, -1, 0, 1, 2, 3 \dots\}$, vale che $\N \subset \Z$, cioè $\N$ è un \textbf{sottoinsieme proprio} di $\Z$, il che vuol dire che $\Z$ contiene tutti gli elementi di $\N$ e anche di più. Simmetricamente, vuole anche dire che tutti gli elementi di $\N$ sono contenuti in $\Z$, ma ci sono elementi di $\Z$ che non sono contenuti in $\N$.\\
		All'insieme dei numeri interi si aggiunge una proprietà fondamentale, che permette di definire una nuova operazione. Tale proprietà è quella dell'\textbf{inverso additivo}, la quale dice che $\forall n\in \Z, \, \exists m \text{ t.c. } n + m = 0$, e tale elemento $m$ viene appunto chiamato \textbf{inverso additivo di $n$} e si indica comunemente come $-n$.\\
		Tale proprietà permette la definizione dell'operazione di \textbf{sottrazione}, che non è altro che una comune somma tra un numero positivo e un numero negativo (che non è detto essere il suo inverso additivo, ma è sicuramente l'inverso additivo di un altro numero).
		
		\item [Numeri razionali] $\Q = \{\frac{p}{q},\, p \in \Z,\, q \in \N, q \neq 0\}$, contiene dunque tutti quei numeri che possono essere scritti come rapporto tra un numero intero e un numero naturale. Importante la proibizione del valore nullo $0$ come denominatore della frazione, in quanto è \textbf{illegale} dividere per $0$. \\
		Vale qui che $\Z \subset \Q$, in quanto ogni numero $p \in \Z$ può essere riscritto come frazione tra il numero stesso e il valore $1$, valore neutro per il prodotto e ora anche per il rapporto.\\
		Una conseguenza della definizione di questo insieme è l'esistenza dell'\textbf{inverso moltiplicativo} di un elemento. Formalmente $\forall a\in \Q, \, \exists b \text{ t.c. } a \cdot b = 1$. Tale elemento $b$ è appunto detto \textit{inverso moltiplicativo di $a$}.
	
	\end{description}
	
	I numeri razionali della forma $\frac{p}{q}$ si possono tutti rappresentare in due modi:
	\begin{itemize}
		
		\item \textbf{Rappresentazione decimale limitata}, dove ogni numero decimale è costituito da un numero finito di cifre a sinistra della virgola (\textbf{parte intera}) e un numero finito di cifre a destra della virgola (\textbf{parte decimale}). Ad esempio $\frac{1}{2} = 0.5$.
		
		\item \textbf{Rappresentazione decimale periodica}, dove ognu numero decimale ha una parte intera finita e una parte decimale con una o più cifre che si ripetono costantemente. Ad esempio $\frac{1}{3} = 0.333333\dots = 0.\overline{3}$.
		
	\end{itemize}
	
	Definendo queste rappresentazioni, già gli antichi greci si accorsero di un problema, di un limite a questo insieme numerico. Prendendo ad esempio un quadrato di lato $l = 1$, come troviamo la diagonale?\\
	Semplice: conoscendo già il teorema di Pitagora, si resero conto che semplicemente la diagonale $d$ si poteva ricavare con $\sqrt{1^2 + 1^2} = \sqrt{2}$. Tracciando poi un arco, si resero conto anche che questo valore era un valore maggiore di $1$ ma sicuramente minore di $2$.\\
	Sorse dunque spontanea la domanda: $\sqrt{2} \in \Q$?
	
	Andiamo dunque a verificare noi stessi se ciò fosse possibile o meno.
	\begin{dimostrazione}{Dimostrazione dell'irrazionalità di $\sqrt{2}$}{}
		Dire che $\sqrt{2} \in \Q$ equivale a dire che possiamo scrivere $\sqrt{2} = \frac{p}{q}$, con $p, q \in \N$. In particolare $\frac{p}{q}$ si prende sempre come una frazione ridotta ai minimi termini, dunque $p$ e $q$ non hanno divisori comuni.
		Andiamo adesso ad applicare delle basilari proprietà algebriche per modellare l'uguaglianza a nostro piacimento:
		\[
			\sqrt{2} = \frac{p}{q} \xrightarrow[\text{moltiplichiamo per q}]{} \sqrt{2}\cdot q = p \xrightarrow[\text{eleviamo al quadrato}]{} p^2 = 2q^2
		\]
		Da questa serie di operazioni algebriche totalmente legali, abbiamo scoperto che $p$ \textbf{è un numero pari}.
		Possiamo ripetere lo stesso processo anche per $q$, dimostrando che anch'esso è un numero pari:
		\[
			\sqrt{2} = \frac{p}{q} \xrightarrow[\text{moltiplichiamo per p}]{} \sqrt{2}\cdot p = q \xrightarrow[\text{eleviamo al quadrato}]{} q^2 = 2p^2
		\]
		Abbiamo così ottenuto che sia $p$ che $q$ sono numeri pari. Questo vuol dire che hanno come divisore in comune $2$. Tuttavia questo contraddice l'ipotesi iniziale secondo cui $MCD(p, q) = 1$. Avendo trovato un assurdo, la nostra dimostrazione è conclusa.
		
	\end{dimostrazione}
	
	Ci si è dunque resi conto della necessità di definire un ulteriore e \textbf{potentissimo} insieme numero.
	\begin{description}
		
		\item [Numeri reali] $\R = \Q\,\cup\,\{\text{numeri irrazionali}\}$, dove i numeri irrazionali sono quei numeri decimali che non possono in alcun modo essere rappresentati come rapporto di due numeri (come ad esempio $\sqrt{2}$).\\
		L'insieme dei numeri reali è l'insieme più utilizzato in tutti i campi dell'analisi matematica. La sua definizione sarebbe complicatissima, e dunque non viene trattata in un corso orientato più alla pratica come questo. Basti sapere che anche a livello di densità numerica $\R$ supera tutti gli insiemi definiti precedentemente. Infatti, ricordando la gerarchia degli insiemi $\N \subset \Z \subset \Q \subset \R$, si può dimostrare come $\#\N = \#\Z = \#\Q << \#\R$. Viene lasciata al lettore la curiosità in merito a questa dimostrazione.
		
		Per l'insieme dei numeri reali valgono le operazioni e proprietà valide per tutti gli altri insiemi definiti precedentemente.
		
	\end{description}
	
	\chapter{Primi elementi di Analisi}
	In questo capitolo, fatte le brevi ma dovute premesse riguardo gli insiemi numerici fondamentali, andiamo a vedere alcuni degli operatori e enti matematici più importanti nel campo dell'analisi.
	
	\section{Valore assoluto}
	Il \textbf{valore assoluto} (anche detto \textbf{modulo}) è un particolare operatore matematico (più avanti impareremo a chiamarlo \textit{funzione}) che, dato un numero, ne restituisce la parte positiva.
	Formalmente, il modulo è definito come segue:
	\[
		|x| =
		\begin{cases}
			x, \text{ se }x \geq 0\\
			-x, \text{ se }x < 0
		\end{cases}
	\]
	Per esempio, se $x = 5,\,|x| = 5$, o anche se $x = -5,\, |x| = 5$.\\
	Geometricamente il modulo ha un significato molto importante: rappresenta la \textbf{distanza} di un numero dall'origine dell'asse delle ordinate dei numeri reali.
	
	Vediamo di seguito alcune proprietà importanti del modulo:
	\begin{enumerate}
		
		\item $|x| = 0 \Leftrightarrow x = 0$
		
		\item $|x \cdot y| = |x| \cdot |y|$
		
		\item \textbf{Disuguaglianza triangolare}: $\forall x,y \in \R,\, |x + y| \leq |x| + |y|$
		
	\end{enumerate}
	
	\section{Gli intervalli}
	Dati $a, b \in \R$, con $a < b$, diamo le seguenti definizioni:
	\begin{description}
		
		\item [Intervalli limitati]: i sottoinsiemi di $\R$ definiti da due disuguaglianze.\\
		Sono intervalli limitati tutti gli intervalli che seguono:
		\begin{itemize}
			
			\item \textbf{Intervallo chiuso}: $[a, b] = \{x \in \R \mid a \leq x \leq b\}$
			
			\item \textbf{Intervallo aperto}: $(a, b) = \{x \in \R \mid a < x < b\}$
			
			\item \textbf{Intervallo semi-aperto a destra}: $[a, b) = \{x \in \R \mid a \leq x < b\}$
			
			\item \textbf{Intervallo semi-aperto a sinistra}: $(a, b] = \{x \in \R \mid a < x \leq b\}$
			
		\end{itemize}
		
		Per tutti questi intervalli limitati la \textit{lunghezza} dell'intervallo si calcola nello stesso modo: $b - a$.
		
		\item [Intervalli illimitati]: i sottoinsiemi di $\R$ definiti da una sola uguaglianza.\\
		Sono intervalli illimitati i seguenti:
		\begin{itemize}
			
			\item \textbf{Intervalli illimitati superiormente}, che possono essere sia chiusi che aperti ma che hanno la seguente struttura: $[a, +\infty) := \{x \in \R \mid x \geq a\}$ oppure $(a, +\infty) := \{x \in \R \mid x > a\}$.
			
			\item \textbf{Intervalli illimitati inferiormente}, anche questi chiusi o aperti, che hanno la seguente struttura: $(-\infty, a] := \{x \in \R \mid x \leq a\}$ oppure $(-\infty, a) := \{x \in \R \mid x < a\}$
			
		\end{itemize}
		
	\end{description}
	
	E' molto importante ricordare che gli intervalli \textbf{non} hanno buchi, dato che sono sottoinsiemi di $\R$ e dunque anche essi sono densi.
	Ad esempio, $(1, 2) \cup [2, 4)$ è un intervallo, in quanto l'unione tra i due intervalli costituisce l'intervallo $(1, 4)$; invece l'unione $(1, 2) \cup (2, 4)$ non rappresenta un intervallo in quanto viene escluso il punto $x = 2$. Essendoci una lacuna, anche di un solo numero, non può essere un intervallo.
	
	Formalizzando tutti i concetti enunciati finora, possiamo fornire una definizione importante al fine di capire alcuni argomenti più avanzati del corso.
	\begin{definizione}{Intorno circolare}{}
		Dati $x_0 \in \R,\,r > 0$ si dice \textbf{intorno con centro $x_0$ e raggio $r$} il seguente intervallo:
		\[
			I(x_0, r) := \{x \in \R \mid |x - x_0| < r\} = (x_0 - r, x_0 + r)
		\]
	\end{definizione}
	
	\section{Maggioranti e minoranti}
	\begin{definizione}{Maggioranti e minoranti}{}
		
		Dato $E \subseteq \R,\, E \neq \O$, un valore $\Lambda \in \R$ è detto un \textbf{maggiorante} di $E$ se risulta $x \geq \Lambda,\, \forall x \in E$.\\
		\smallbreak
		
		Dato $E \subseteq \R,\, E \neq \O$, un valore $\lambda \in \R$ è detto un \textbf{minorante} di $E$ se risulta $x \leq \lambda,\, \forall x \in E$.
		
	\end{definizione}
	
	Possiamo così definire $M(E)$ come l'insieme dei maggioranti di $E$, e analogamente $m(E)$ come l'insieme dei minoranti di $E$.
	
	\begin{definizione}{Insieme limitato}{}
		
		Dato $E \subseteq \R,\, E \neq \O$, tale $E$ si dice:
		\begin{description}
			
			\item [Limitato superiormente]: se $M(E) \neq \O$.
			
			\item [Limitato inferiormente]: se $m(E) \neq \O$.
			
			\item [Limitato superiormente]: se $M(E) \neq \O$ e $m(E) \neq \O$.
			
		\end{description}
		
		Analogamente possiamo dire che un intervallo è \textbf{illimitato} (superiormente/inferiormente/totalmente) se il corrispondente insieme di maggioranti/minoranti è \textbf{vuoto}.
		
	\end{definizione}
	
	Ad esempio, l'intervallo $E = (3, 5)$ è:
	\begin{itemize}
		
		\item Illimitato inferiormente in quanto $m(E) = {x \in \R \mid x \leq 3} \neq \O$.
		
		\item Illimitato superiormente in quanto $M(E) = {x \in \R \mid x \geq 5} \neq \O$.
		
	\end{itemize}
	
	Ed è dunque limitato totalmente.\\
	Per motivi simili, diciamo invece che l'intervallo $E = (2, +\infty)$ è limitato inferiormente e illimitato superiormente.
	
	\section{Massimi e minimi}
	\begin{definizione}{Minimo e massimo di un insieme}{}
		
		Dato $E \subseteq \R,\, E \neq \O$.\\
		Il valore $M \in \R$ è il \textbf{massimo} di $E$ se:
		\begin{itemize}
			
			\item $M$ è un maggiorante, dunque $x \leq M,\,\forall x\in E$
			
			\item $M \in E$
			
		\end{itemize}
		In tal caso si scrive $M = \max(E)$.\\
		\smallbreak
		
		Dato $E \subseteq \R,\, E \neq \O$.\\
		Il valore $m \in \R$ è il \textbf{minimo} di $E$ se:
		\begin{itemize}
			
			\item $m$ è un minorante, dunque $x \geq m,\,\forall x\in E$
			
			\item $m \in E$
			
		\end{itemize}
		In tal caso si scrive $m = \min(E)$.

	\end{definizione}
	
	\section{Sottoinsiemi di $\N$}
	Considerando l'insieme dei numeri naturali $\N = \{0, 1, 2, 3, \dots\}$, proviamo a prenderne un qualsiasi sottoinsieme:
	\begin{itemize}
		
		\item $E = \{n \in \N \mid n \text{ è un numero pari}\} = \{0, 2, 4, \dots\}$
		
		\item $E = \{n \in \N \mid n^2 < 10\} = \{0, 1, 2, 3\}$
		
		\item $E = \{n \in \N \mid n \text{ è un numero primo}\} = \{1, 3, 5, \dots\}$
		
	\end{itemize}
	
	Notiamo una specie di "pattern": tutti i sottoinsiemi hanno \textbf{sempre un minorante e un minimo}. Anche l'insieme $\N$ stesso ha un minimo, cioè il numero $0$.
	
	\begin{osservazione}{}{}
		
		Se $\Lambda$ è un maggiorante, allora lo sono anche $\Lambda + 1, \Lambda + 2, \dots$ e l'insieme dei maggioranti è un insieme infinito.\\		
		Discorso analogo vale per i minoranti: se lo è $\lambda$ allora lo sono anche $\lambda - 1, \lambda - 2, \dots$
		
	\end{osservazione}
	
	Possiamo dunque trarre due importanti conclusioni:
	\begin{itemize}
		
		\item $m(E) = \{\lambda \in \R \mid \lambda \leq n,\, \forall n \in E\}$
		
		\item $\max(m(E)) = \min(E)$
		
	\end{itemize}
	
	\section{Sottoinsiemi di $\Z$}
	Il funzionamento dell'insieme $\Z$ è analogo a quello di $\N$, con l'unica eccezione che non possiamo dire con certezza che esista \textit{sempre} un minimo, in quanto $\Z$ di per sé è illimitato inferiormente.\\
	Valgono tuttavia proposizioni analoghe a quelle per $\N$:
	\begin{itemize}
		
		\item Se $m(E) \neq \O \Rightarrow \max(m(E)) = \min(E)$
		
		\item Se $M(E) \neq \O \Rightarrow \min(M(E)) = \max(E)$
		
	\end{itemize}
	
	\section{Sottoinsiemi di $\Q$}
	Tali proprietà valgono anche per $\Q$? Vediamolo con alcuni esempi:
	\begin{itemize}
		
		\item Se prendiamo un sottoinsieme $E = \{x \in \Q \mid 0 \leq x \leq 1\}$, abbiamo sia i minorandi, cioè tutti i numeri minori di $0$, che i maggiorandi, cioè tutti i numeri maggiori di $1$. E in questo caso vale anche che il minore dei maggioranti è il massimo dell'insieme ($1$) e discorso analogo per il massimo dei minoranti.
		
		\item Se invece prendiamo un sottoinsieme $E = \{x \in \Q \mid x^2 \leq 2\}$ vediamo come qui il discorso non valga! Il minimo dei maggioranti sarebbe $\sqrt{2}$, che però \textbf{non è un elemento di $\Q$}, e dunque non può essere un massimo del sottoinsieme.
		
	\end{itemize}
	
	Le proprietà viste prima non valgono per questo insieme numerico.
	
	\section{Sottoinsiemi di $\R$}
	La situazione per i numeri reali diventa un vero e proprio assioma di completezza dell'insieme.
	\begin{teorema}{Assioma di completezza dell'insieme dei Reali}{}
		
		Sia $E \subseteq \R,\, E \neq \O$.
		\begin{itemize}
			
			\item Se $M(E) \neq \O$, allora \textbf{esiste} il minimo di $M(E)$ e tale valore si dice \textbf{estremo superiore di $E$} ($\operatorname{sup}(E)$). Tale estremo superiore è un numero reale se $M(E) \neq \O$, altrimenti è uguale a $+\infty$.
			
			\item Se $m(E) \neq \O$, allora \textbf{esiste} il massimo di $m(E)$ e tale valore si dice \textbf{estremo inferiore di $E$} ($\operatorname{inf}(E)$). Tale estremo inferiore è un numero reale se $m(E) \neq \O$, altrimenti è uguale a $-\infty$.
			
		\end{itemize}
		
		\textbf{Tutti} i sottoinsiemi di $\R$ hanno un estremo superiore e uno inferiore.
	\end{teorema}
	
	\chapter{Funzioni}
	Nel seguente capitolo descriveremo uno degli enti matematici più importanti e fondamentali di tutta l'analisi. Si potrebbe addirittura dire l'ente sul quale si basa tutta la disciplina: le \textbf{funzioni}.
	
	\begin{definizione}{Funzione}{}
		
		Siano $X, Y$ insiemi.\\
		Una funzione $f$ definita in $X$ a valori in $Y$ è una \textbf{legge} che associa ad \textit{ogni elemento $x\in X$ uno e un solo elemento $y \in Y$}.\\
		La notazione formale per la definizione di una funzione è la seguente:
		\[
			\begin{aligned}
				f: X &\to Y\\
				x &\mapsto y
			\end{aligned}
		\]
		In questo caso si dice che $X$ è il \textbf{dominio} della funzione e $Y$ il \textbf{codominio} di questa.\\
		
		\smallbreak
		Nei libri di testo è inoltre molto comune trovare la seguente notazione:
		\[
			y = f(x)
		\]
		il che vuol dire che siamo esprimendo la variabile $y$ \textbf{in funzione} della variabile $x$. Da questa notazione possiamo dire che $x$ è una variabile \textbf{indipendente} (scegliamo noi il suo valore) mentre $y$ è una variabile \textbf{dipendente} (in quanto il suo valore \textit{dipende} da quello di $x$ e dalla legge applicata).\\
		
		\smallbreak
		L'insieme degli elementi per cui vale $y = f(x)$ formano l'\textbf{immagine} di $f$, che formalmente indichiamo come segue:
		\[
			\operatorname{Im}(f) = f(X) := \{y = f(x),\, x \in X\} \subseteq Y
		\]
		
	\end{definizione}
	
	Forniamo adesso alcuni esempi non matematici di funzioni, per rafforzare il concetto:
	\begin{itemize}
		
		\item Siano $X = \{\text{Insieme degli studenti Sapienza}\}, Y = \{\text{Insieme dei numeri di matricola}\}$. Siccome possiamo associare ad ogni studente uno e un solo numero di matricola, questa è una funzione valida ($f: Studente \mapsto Matricola$).
		
		\item Siano $X = \{\text{Insieme degli studenti Sapienza}\}, Y = \{\text{Insieme delle altezze in cm}\}$. Siccome ad uno studente associamo una sola altezza, questa è una funzione valida ($f: Studente \mapsto Altezza$). Nota bene: nonostante più studenti abbiano la stessa altezza, l'importante è che lo stesso studente non ne abbia due. Rimane comunque una funzione!
		
	\end{itemize}
	
	Ma se ad esempio volessimo definire $X = \{\text{Insieme delle sedie nell\'aula}\}, Y = \{\text{Insieme degli studenti in aula}\}$, una legge $f: X \to Y$ non sarebbe una funzione, in quanto possono esserci sedie vuote (gli elementi del dominio devono essere mappati tutti!!!)
	
	Data la definizione di funzione, procediamo col definire dei concetti fondamentali nello studio delle funzioni.
	
	\begin{definizione}{Suriettività}{}
		
		$f: X \to Y$ è \textbf{suriettiva} se $\operatorname{Im}(f) = Y$, cioè $\forall y \in Y,\,\exists x \in X \text{ tale che } f(x) = y$
		
	\end{definizione}
	Detto in termini molto spiccioli e molto poco formali, una funzione si dice suriettiva se preso un \textit{qualsiasi} valore del codominio questo "proviene" da un valore del dominio.
	
	\begin{osservazione}{}{}
		
		Se $f: X \to Y$ non è suriettiva, allora possiamo considerare una nuova funzione $g$ con lo stesso dominio ma con codominio ridotto ad essere uguale all'immagine di $f$. Tale funzione sarà dunque definita come $g: X \to \operatorname{Im}(f)$ e sarà \textbf{suriettiva}. Diciamo che in questo modo stiamo "forzando" il fatto che $f$ diventi suriettiva.
		
	\end{osservazione}
	
	\begin{definizione}{Iniettività}{}
		
		$f: X \to Y$ è \textbf{iniettiva} se $\forall x_1, x_2 \in X,\, x_1 \neq x_2 \text{ vale che } f(x_1) \neq f(x_2)$.\\
		Se invece \textbf{non} premettessimo $x_1 \neq x_2$ allora avremo che $x_1 = x_2 \Leftrightarrow f(x_1) = f(x_2)$.
		
	\end{definizione}
	Questo ci dice che una funzione $f$ è iniettiva se, presa una qualsiasi coppia di valori del dominio, le loro immagini saranno diverse. L'uguaglianza delle immagini deve implicare che anche i due valori scelti coincidano. Se così non fosse, la funzione non sarebbe iniettiva.
	
	Un esempio di funzione \textbf{non} iniettiva è la funzione $f: Studente \to Altezza$, in quanto due studenti diversi (i due valori $x_1, x_2$ del dominio) possono avere la stessa altezza (il singolo valore $y$ del codominio).
	
	\begin{definizione}{Biettività}{}
		Una funzione $f$ si dice \textbf{biettiva} (o \textbf{biunivoca}) se è al contempo iniettiva e suriettiva.
	\end{definizione}
	
	Per convenzione, le funzioni che considereremmo d'ora in poi, salvo dovute precisazioni, saranno del tipo $f: X \subseteq \R \to Y \subseteq \R$
	
	\section{Grafici di funzioni}
	
	Prima di vedere graficamente come si presentano i grafici di alcune funzioni elementari, procediamo a dare una definizione formale di \textit{grafico di funzione}.
	
	\begin{definizione}{Grafico di una funzione}{}
		
		Il grafico di una funzione $f$ è un sottoinsieme del piano cartesiano $\R^2 = \R \times \R$:
		\[
			\Gamma(f) = \{(x, f(x)),\, x \in X\}
		\]
		
		\begin{center}
			
			\begin{tikzpicture}[scale=1.2]
				
				\draw[->] (-0.3, 0) -- (2, 0) node[right] {\small ascisse};
				\draw[->] (0, -0.3) -- (0, 2) node[above] {\small ordinate};
				\node[below left] at (0, 0) {$0$};
				
				\draw[dashed] (1, 1) -- (1, 0) node[below] {\small $x$};
				\draw[dashed] (1, 1)  -- (0, 1) node[left] {\small $y$};;
				\fill[black] (1, 1) circle (1.5pt) node[right] {\small $(x, f(x))$};
				
			\end{tikzpicture}
			
		\end{center}
		
	\end{definizione}
	
	Presentiamo di seguito alcuni grafici di funzioni elementari, estremamente ricorrenti ed importanti nel campo dell'analisi.
	
	\begin{description}
		
		\item [Retta generica,] descritta dalla funzione $f(x) = mx + q$, il cui grafico è il seguente:
		
		\begin{center}
			
			\begin{tikzpicture}
				
				\begin{axis}[axis lines=center, xmin=-3, xmax=3, ymin=-3, ymax=3, ytick={-3, -2, ..., 3}, domain=-2.9 : 2.9]
					
					\addplot[blue, thick] {x};
					
				\end{axis}
				
			\end{tikzpicture}
			
		\end{center}
		
		Nella funzione generica di una retta vediamo avere due parametri che la descrivono:
		\begin{itemize}
			
			\item $m$, il coefficiente moltiplicativo della $x$, che rappresenta \textbf{la pendenza della retta}. Se $m > 0$ allora la retta ha una pendenza positiva (informalmente diciamo che \textit{cresce}) mentre se $m < 0$ la retta ha una pendenza negativa (\textit{decresce}).
			
			\item $q$, il valore della coordinata $y$ in cui la retta interseca l'asse dell'ordinate. Sostanzialmente la retta interseca l'asse delle ordinate nel punto $(0, q)$.
			
		\end{itemize}
		
		\item [Parabola generica,] descritta dalla funzione $f(x) = ax^2 + bx + c$ il cui grafico è il seguente:
		\begin{center}
			
			\begin{tikzpicture}
				
				\begin{axis}[axis lines=center, xmin=-3, xmax=3, ymin=-1.5, ymax=3, ytick={-3, -2, ..., 3}, domain=-3 : 3, samples=100]
					
					\addplot[blue, thick] {3*x^2 + 2*x - 1};
					
				\end{axis}
				
			\end{tikzpicture}
			
		\end{center}
		
		Anche per questa funzione abbiamo dei parametri fondamentali:
		\begin{itemize}
			
			\item $a$, che rappresenta l'ampiezza della parabola: maggiore è $|a|$, più "aperta" sarà la parabola. Inoltre, se $a > 0$ la concavità della parabola è rivolta verso l'alto, altrimenti con $a < 0$ la concavità è rivolta verso il basso.
			
			\item $b$ indica la posizione dell'asse di simmetria della parabola, ed è infatti il parametro da cui derivano le formule per il calcolo delle coordinate del vertice e del fuoco di una parabola.
			
			\item $c$, similarmente al parametro $q$ della retta, indica il punto di contatto tra la funzione e l'asse delle ordinate.
			
		\end{itemize}
		
		\newpage
		\item [Valore assoluto,] la cui funzione è $f(x) = |x|$ e che descrive il seguente grafico:
		\begin{center}
			
			\begin{tikzpicture}
				
				\begin{axis}[axis lines=center, xmin=-3.1, xmax=3.1, ymin=-0.5, ymax=3.1, ytick={-3, -2, ..., 3}, domain=-2.5 : 2.5]
					
					\addplot[blue, thick] {abs(x)};
					
				\end{axis}
				
			\end{tikzpicture}
			
		\end{center}
		
		\item [Funzioni razionali generiche,] come ad esempio $f(x) = \frac{1}{x}$. Notare che questa funzione \textbf{non} è assolutamente definita per $x = 0$, in quanto come già detto non è assolutamente ammessa la divisione per zero.
		\begin{center}
			
			\begin{tikzpicture}
				
				\begin{axis}[axis lines=center, xmin=-3.1, xmax=3.1, ymin=-3.1, ymax=3.1, ytick={-3, -2, ..., 3}, domain=-3:3, samples=100]
					
					\addplot[blue, thick, samples=100, domain=-3:-0.01] {1/x};
					
					\addplot[blue, thick, samples=100, domain=0.01:3] {1/x};
					
				\end{axis}
				
			\end{tikzpicture}
			
		\end{center}
		
		Non essendo definito il punto $f(0)$ si spiega (anche se verrà approfondito nei capitoli successive) il comportamento bizzarro del grafico della funzione.
		
		\newpage
		\item [Funzioni trigonometriche: seno,] descritto dalla funzione $f(x) = sin(x)$:
		
		\begin{center}
			
			\begin{tikzpicture}
				
				\begin{axis}[axis lines=center, xmin=-3.1, xmax=3.1, ymin=-3.1, ymax=3.1, ytick={-3, -2, ..., 3}, domain=-3:3, samples=100, trig format plots=rad]
					
					\addplot[blue, thick] {sin(x)};
					
				\end{axis}
				
			\end{tikzpicture}
			
		\end{center}
		
		\item [Funzioni trigonometriche: coseno,] descritto dalla funzione $f(x) = cos(x)$:
		
		\begin{center}
			
			\begin{tikzpicture}
				
				\begin{axis}[axis lines=center, xmin=-3.1, xmax=3.1, ymin=-3.1, ymax=3.1, ytick={-3, -2, ..., 3}, domain=-3:3, samples=100, trig format plots=rad]
					
					\addplot[blue, thick] {cos(x)};
					
				\end{axis}
				
			\end{tikzpicture}
			
		\end{center}
		
	\end{description}
	
	Avendo mostrato alcuni dei grafici fondamentali, andiamo a vedere quali "operazioni" possono essere compiute su una funzione e come ne varia il grafico.
	
	\newpage
	\subsection{Operazioni sui grafici}
	
	Data una funzione $f(x)$, con un certo grafico $\Gamma(f)$:
	
	\begin{description}
		
		\item [Traslazione verticale,] quando andiamo a definire una funzione $g(x) = f(x) + c$, ovvero quando sommiamo un coefficiente \textit{fuori} dall'argomento della funzione di partenza.
		
		\begin{center}
			
			\begin{tikzpicture}
				
				\begin{axis}[axis lines=center, xmin=-3.1, xmax=6.1, ymin=-1.5, ymax=4.1, ytick={-3, -2, ..., 3}, domain=-3:3, samples=100, trig format plots=rad]
					
					\addplot[blue, thick] {x^2 + 1};
					\addplot[green!60!black, thick] {x^2};
					\addplot[red, thick] {x^2 - 1};
					
					\node[blue, right] at (axis cs: 2, 2) {$f(x) = x^2+1$};
					\node[green!60!black, right] at (axis cs: 2, 1) {$f(x) = x^2$};
					\node[red, right] at (axis cs: 2, -1) {$f(x) = x^2-1$};
					
				\end{axis}
				
			\end{tikzpicture}
			
		\end{center}
		
		\item [Traslazione orizzontale,] data dalla definizione di una funzione $g(x) = f(x + c)$, dove dunque aggiungiamo un coefficiente direttamente nell'argomento della funzione.
		
		\begin{center}
			
			\begin{tikzpicture}
				
				\begin{axis}[axis lines=center, xmin=-3.1, xmax=7.5, ymin=-1.5, ymax=4.1, ytick={-3, -2, ..., 3}, domain=-3:3, samples=100, trig format plots=rad]
					
					\addplot[blue, thick] {(x+1)^2};
					\addplot[green!60!black, thick] {x^2};
					\addplot[red, thick] {(x-1)^2};
					
					\node[blue, right] at (axis cs: 3, 3) {$f(x) = (x+1)^2$};
					\node[green!60!black, right] at (axis cs: 3, 2) {$f(x) = (x)^2$};
					\node[red, right] at (axis cs: 3, 1) {$f(x) = (x-1)^2$};
					
				\end{axis}
				
			\end{tikzpicture}
			
		\end{center}
		
		\newpage
		\item [Dilatazione,] data da $g(x) = f(c\cdot x)$, cioè dal prodotto di un coeffciente direttamente con la variabile indipendente $x$. Questo fa si che il grafico si dilati di un fattore $\frac{1}{c}$.
		
		\begin{center}
			
			\begin{tikzpicture}
				
				\begin{axis}[axis lines=center, xmin=-5.1, xmax=8.1, ymin=-1.5, ymax=4.1, ytick={-3, -2, ..., 3}, domain=-3:3]
					
					\addplot[blue, thick] {abs(2*x)};
					\addplot[green!60!black, thick] {abs(x)};
					\addplot[red, thick] {abs(1/2 * x)};
					
					\node[blue, right] at (axis cs: 3.5, 2) {$f(x) = |2x|$};
					\node[green!60!black, right] at (axis cs: 3.5, 1) {$f(x) = |x|$};
					\node[red, right] at (axis cs: 3.5, -1) {$f(x) = |\frac{1}{2}x|$};
					
				\end{axis}
				
			\end{tikzpicture}
			
		\end{center}
		
	\end{description}
	
	\section{Funzioni inverse}
	
	\begin{definizione}{Funzione inversa}{}
		
		Sia $f: X \to Y$ una funzione \textbf{biunivoca}.\\
		Definiamo una nuova funzione che ad ogni $y \in Y$ associa una $x \in X$ tale che $f(x) = y$.
		Chiamiamo questa funzione \textbf{funzione inversa} e la indichiamo con $f^{-1}$.
		
	\end{definizione}
	
	Formalmente, vediamo che in questa il dominio e il codominio sono invertiti rispetto alla funzione di partenza $f$:
	\[
		\begin{aligned}
			f^{-1}: Y &\to X\\
			y &\mapsto x
		\end{aligned}
	\]
	Diamo adesso due uguaglianze fondamentali per lavorare con le funzioni inverse:
	\begin{itemize}
		
		\item $f(x) = y = f(f^{-1}(y)),\,\forall y \in Y$
		
		\item $x = f^{-1}(y) = f^{-1}(f(x)),\,\forall x \in X$
		
	\end{itemize}
	
	Da ciò derivano dunque le seguenti \textit{importantissime} uguaglianze:
	\begin{itemize}
		
		\item Se definiamo la catena di funzioni $X \xrightarrow{f} Y \xrightarrow{f^{-1}} X$, otteniamo la seguente mappatura:
		\[
			x \mapsto y=f(x) \mapsto f^{-1}(y) = f^{-1}(f(x)) = x
		\]
		
		\item Se definiamo la catena di funzioni $Y \xrightarrow{f^{-1}} X \xrightarrow{f} Y$, otteniamo la seguente mappatura:
		\[
		y \mapsto x=f^{-1}(y) \mapsto f(x) = f(f^{-1}(y)) = y
		\]
		
	\end{itemize}
	
	Come vediamo, se definiamo una catena di funzioni in cui sono alternativamente richiamate una funzione e la sua inversa, avremo come risultato \textbf{sempre} la \textbf{funzione identità} rispetto alla variabile di partenza.
	
	Nel pratico, negli esercizi, ricavare una funzione inversa da una funzione $y = f(x)$ di partenza si traduce semplicemente nell'isolare la variabile $x$ ed esprimerla in funzione della variabile $y$.
	
	Prima di concludere questo discorso, presentiamo un piano cartesiano in cui sono rappresentate una funzione e la sua inversa, per commentare il comportamento di queste:
	
	\begin{center}
		
		\begin{tikzpicture}
			
			\begin{axis}[axis lines=center, xmin=-0.1, xmax=8.1, ymin=-0.1, ymax=4.1, ytick={-3, -2, ..., 3}, samples=100]
				
				\addplot[blue, thick, domain=0:3] {x^2};
				\addplot[red, thick, domain=0:8] {sqrt(x)};
				
				\node[blue, right] at (axis cs: 3.5, 3) {$f(x) = x^2$};
				\node[red, right] at (axis cs: 3.5, 1) {$f(x) = \sqrt{x}$};
				
			\end{axis}
			
		\end{tikzpicture}
		
	\end{center}
	
	Possiamo vedere come il grafico delle due funzioni risulti specchiato rispetto all'asse delle x.
	
	\newpage
	\section{Esponenziali}
	
	Sia $a > 0$, proviamo a calcolare $a^x$, con $x = \frac{p}{q}\in \Q$.
	\begin{itemize}
		
		\item $a^{\frac{p}{q}} = \sqrt[q]{a^p} = (a^{\frac{1}{q}})^p$
		
		\item E se invece $x$ fosse un numero irrazionale, come ad esempio $x = \sqrt{2}$? Supponiamo di avere $a = 2$ e di voler calcolare $2^{\sqrt{2}}$
		Proviamo a procedere calcolando prima l'insieme di tutti i valori che si avvicinano a $2^{\sqrt{2}}$, calcolando i valori razionali. 
		\[
			E = \{2^r,\,r = \frac{p}{q},\, r \leq \sqrt{2}\} \to
			\begin{cases}
				E \neq \O\\
				M(E) \neq \O \Rightarrow \exists \min(M(E)) = \sup(E) = 2^{\sqrt{2}}
			\end{cases}
		\]
		
	\end{itemize}
	
	Avendo dunque definito il generico esponenziale reale come $a^x = \sup(E)$, possiamo definire a sua volta anche $a^{-x} = \frac{1}{a^x}$.\\
	Di conseguenza, definiamo la funzione $f(x) = a^x$ che va da $\R$ ad $\R$.
	
	Presentiamo adesso alcune fondamentali proprietà dell'esponenziale
	\begin{enumerate}
		
		\item $a^x > 0,\,\forall x \in \R$
		
		\item $a^0 = 1$
		
		\item $a^{x+y} = a^x \cdot a^y,\, \forall x, y \in \R$
		
		\item $(a^x)^y = a^{x \cdot y}$
		
	\end{enumerate}
	
	Per la proprietà (1) possiamo restringere il codominio in modo da coprire solo $\R_+$, rendendo la funzione biunivoca e potendo così definire $f^-1(y) = \log_a(y)$.
	
	In questo particolare caso i grafici saranno i seguenti:
	\begin{center}
		
		\begin{tikzpicture}
			
			\begin{axis}[axis lines=center, xmin=-3.1, xmax=8.1, ymin=-3.1, ymax=4.1, samples=100]
				
				\addplot[blue, thick] {2^x};
				\addplot[red, thick] {log2(x)};
				
				\node[blue, right] at (axis cs: 3.5, 3) {$f(x) = 2^x$};
				\node[red, right] at (axis cs: 3.5, 1) {$f(x) = \log_2(x)$};
				
			\end{axis}
			
		\end{tikzpicture}
		
	\end{center}
	
	Come vediamo, se $a^x > 0$, è richiesto necessariamente che l'argomento del logaritmo sia a sua volta \textbf{sempre} $> 0$.
	
	\section{Funzioni pari e dispari}
	
	Sia $f: \R \to \R$ una funzione qualsiasi.
	\begin{definizione}{Funzione pari}{}
		
		$f$ è una funzione \textbf{pari} se $f(x) = f(-x),\,\forall x \in \R$
		Ad esempio, tutte le funzioni del tipo $f(x) = x^{2n} + c$ sono funzioni pari (quindi $x^2, x^4, \dots$); anche le funzioni $f(x) = \cos(x) + c$ e $f(x) = |x| + c$ sono pari.
		
	\end{definizione}
	
	Nel caso delle funzioni pari, il grafico è \textbf{simmetrico rispetto all'asse $y$}.
	
	\begin{definizione}{Funzione dispari}{}
		
		$f$ è una funzione \textbf{dispari} se $f(x) = -f(-x),\,\forall x \in \R$
		Ad esempio, tutte le funzioni del tipo $f(x) = x^{2n+1}$ sono funzioni dispari (quindi $x^2, x^4, \dots$); anche la funzione $f(x) = \sin(x)$ è dispari.
		
	\end{definizione}
	
	\chapter{Successioni}
	
	\begin{definizione}{Successione}{}
		
		Le \textbf{successioni} sono dei tipi particolari di funzioni che sono definite da $\N$ ad $\R$. Formalmente la definizione di una successione è la seguente:
		\[
		\begin{aligned}
			a: \N &\to \R\\
			n &\mapsto a(n)
		\end{aligned}
		\]
		Più brevemente, indichiamo una successione in uno dei seguenti modi: $a_n, (a_n)_{n \in \N}, \{a_n\}_{n \in \N}$.
		
	\end{definizione}
	
	Come si può intuire dal nome, una successione consiste in un susseguirsi di valori. A differenza delle funzioni che hanno come dominio $\R$, con le successioni se provassimo a tracciare il grafico di una successione noteremmo dei "buchi", dato che $\N$ non è \textbf{affatto} denso come $\R$.
	Di seguito un esempio grafico:
	
	\begin{center}
		
		\begin{tikzpicture}
			
			\draw[->] (-0.1, 0) -- (4, 0);
			\draw[->] (0, -0.1) -- (0, 4);
			
			\foreach \x in {1, 2, 3} {
				\draw (\x, 0.05) -- (\x, -0.05) node[below] {\small$\x$};
			}
			
			\foreach \y in {1, 2, 3} {
				\draw (0.05, \y) -- (-0.05, \y) node[left] {\small$\y$};
			}
			
			\fill[blue] (0,0) circle (1.5pt);
			\fill[blue] (1,1) circle (1.5pt);
			\fill[blue] (2,2) circle (1.5pt);
			\fill[blue] (3,3) circle (1.5pt) node[right] {$a_n = n$};
			
		\end{tikzpicture}
		
	\end{center}
	
	Ma non necessariamente una successione deve avere una definizione così semplice. Si possono avere casi di tutti i tipi, come ad esempio:
	\begin{itemize}
		
		\item $a_n = (-1)^n = 1, -1, 1, -1, \dots$
		
		\item $a_n = c = c, c, c, \dots$ 
		
	\end{itemize}
	
	Insomma, ci si può sbizzarrire con la legge di definizione delle successioni (ma vedremo sarà così anche per le funzioni).
	
	\begin{osservazione}{}{}
		
		Può accadere che una successione $a_n$ sia definita da un certo indice $n_0$ in poi:
		\[
			(a_n)_{n \geq n_0}
		\]
		
	\end{osservazione}
	
	Ad esempio la successione $a_n = \frac{1}{n}$ è necessariamente definita a partire da $n_0 = 1$, dato che non si può dividere per $0$.
	
	Come per le funzioni, definiamo adesso l'importantissimo concetto di limitazione di una successione.
	
	\begin{definizione}{Successione limitata}{}
		
		Una successione $(a_n)$ si dice:
		\begin{description}
			
			\item [Limitata inferiormente] se $\exists h \in \R \text{ tale che } h \leq a_n,\,\forall n \in \N$
			
			\item [Limitata superiormente] se $\exists k \in \R \text{ tale che } k \geq a_n,\,\forall n \in \N$
			
			\item [Limitata] se è contemporaneamente limitata superiormente ed inferiormente.
			
		\end{description}
		
	\end{definizione}
	
	Ad esempio, vediamo come la successione $a_n = \frac{1}{n} = 1, \frac{1}{2}, \frac{1}{3}, \dots$ sia:
	\begin{itemize}
		
		\item Limitata superiormente da $a_1 = 1$, dato che sarà il valore più grande che la successione potrà raggiungere (in questo caso sarà anche il \textbf{massimo} della successione).
		
		\item Limitata inferiormente da $0$, in quanto vediamo come i termini della successione continuano a descrescere, ma non si arriverà \textbf{mai} precisamente al valore $0$. Questo sarà dunque il limite inferiore, che però in questo caso \textit{non} coincide col minimo.
		
	\end{itemize}
	
	Come abbiamo inoltre già visto per le funzioni, vale dire che una successione che \textbf{non} è limitata inferiormente si dice \textbf{illimitata inferiormente}, e analogamente se \textbf{non} lo è superiormente si dice \textbf{illimitata superiormente}.
	
	\newpage
	\section{Limiti di successioni}
	
	Per spiegare i limiti di successioni, procediamo con un esempio:
	\begin{itemize}
		
		\item Sia $a_n = \frac{1}{n+1} = (1, \frac{1}{2}, \frac{1}{3}, \dots)$
		
		\item Se continuassimo a contare all'infinito, vedremo come i valori della successione si avvicinerebbero \textbf{sempre di più} allo $0$, senza però mai raggiungerlo effettivamente. Possiamo anche vederlo dal seguente grafico:
		
		\begin{center}
			
			\begin{tikzpicture}[scale=1.3]
				
				\draw[->, thick] (0, 0) -- (4.8, 0);
				\draw[->, thick] (0, 0) -- (0, 3.2);
				
				\node[below left] at (0,0) {$0$};
				
				\foreach \x in {1, 2, 3, 4} {
					\draw (\x, 0.05) -- (\x, -0.05) node[below] {$\x$};
				}
				
				\draw (0.05, 3) -- (-0.05, 3) node[left] {$1$};
				\draw (0.05, 1.5) -- (-0.05, 1.5) node[left] {$\frac{1}{2}$};
				
				\foreach \n / \y in {0/3, 1/1.5, 2/0.9, 3/0.6, 4/0.4} {
					
					\ifnum \n > 0
					\draw[dashed, gray!80] (0, \y) -- (\n, \y) -- (\n, 0);
					\fi
					
					\fill[blue] (\n, \y) circle (2pt);
					
					\node[blue, right=3pt] at (\n, \y) {$a_{\n}$};
				}
				
			\end{tikzpicture}
			
		\end{center}
		
		Calcolare il limite di una successione equivale a chiedersi: \textit{che tipo di comportamento assueme la successione al crescere di $n$? Si stabilizza attorno ad un valore prefissato $l$ oppure continua fino a $\pm \infty$?}
		
	\end{itemize}
	
	Si indica dunque come $\lim_{n \to +\infty} a_n$.
	
	\begin{osservazione}{}{}
		Calcoliamo il limite SOLO per valori di $n$ che vanno a $+\infty$ perché le successioni coprono TUTTI i numeri naturali (al massimo iniziano qualche $n$ dopo rispetto a $0$), al contrario delle funzioni che possono avere buchi nel dominio.
	\end{osservazione}
	
	I comportamenti sono i seguenti:
	\begin{itemize}
		\item Se si dice che la successione \textbf{CONVERGE} a $l$:
		\[
		\lim_{n \to +\infty} a_n = l
		\]
		Formalmente scriveremmo:
		\[
		\forall \epsilon > 0 \quad \exists n_\epsilon \in \N \text{ tale che } |a_n - l| < \epsilon \quad \forall n \ge n_\epsilon
		\]
		ovvero $l - \epsilon < a_n < l + \epsilon$.
		\smallbreak
		In soldoni, ciò vuol dire che, dopo un certo indice $n_\epsilon$, tutti i numeri prodotti dalla successione ``cadranno'' nella striscia compresa tra $l-\epsilon$ ed $l+\epsilon$.
	\end{itemize}
	
	\textbf{Esempio grafico:}
	\begin{center}
		\begin{tikzpicture}[scale=1.2]
			% Assi
			\draw[->, thick] (-0.2, 0) -- (6, 0) node[right] {$n$};
			\draw[->, thick] (0, -0.2) -- (0, 3.5) node[above] {$a_n$};
			
			% Limite e striscia epsilon
			\draw[red, thick, dashed] (0, 2) node[left] {$l$} -- (5.5, 2);
			\fill[red, opacity=0.1] (0, 1.5) rectangle (5.5, 2.5);
			\draw[gray, dashed] (0, 2.5) node[left, text=black] {$l+\epsilon$} -- (5.5, 2.5);
			\draw[gray, dashed] (0, 1.5) node[left, text=black] {$l-\epsilon$} -- (5.5, 1.5);
			
			% Punti della successione
			\fill[blue] (0.5, 3.2) circle (2pt) node[above right] {$a_1$};
			\fill[blue] (1.2, 0.8) circle (2pt) node[below right] {$a_2$};
			\fill[blue] (1.8, 2.8) circle (2pt) node[above right] {$a_3$};
			
			% Indice n_epsilon e successivi (dentro la striscia)
			\draw[dashed, thick] (2.5, 2.2) -- (2.5, 0) node[below] {$n_\epsilon$};
			\fill[green!60!black] (2.5, 2.2) circle (2pt) node[above right] {$a_4$};
			\fill[green!60!black] (3.2, 1.7) circle (2pt) node[below right] {$a_5$};
			\fill[green!60!black] (4.0, 2.1) circle (2pt) node[above right] {$a_6$};
			\fill[green!60!black] (4.8, 1.9) circle (2pt) node[below right] {$a_7$};
			
			\node[text width=4cm, align=left] at (4, 4) {\small Dopo l'indice $n_\epsilon$, tutti i valori della successione cadono nel range del limite.};
		\end{tikzpicture}
	\end{center}
	
	\subsection*{Calcolo di un limite (Esempio formale)}
	Sia $a_n = \frac{n}{n+1} = \left(0, \frac{1}{2}, \frac{2}{3}, \frac{3}{4}, \dots \right)$.\\
	Per numeri molto grandi, tendenti a $+\infty$, non c'è praticamente differenza tra $n$ ed $n+1$. Si potrebbe dire che, per $n \to +\infty$, $n \simeq n+1$.
	La successione, per $n \to +\infty$, diventa quindi $\frac{n}{n}$, e il candidato limite diventa $1$.
	
	Vediamo se, seguendo la definizione di limite, riusciamo a verificare in modo più formale la cosa. Ricordiamoci la relazione $|a_n - l| < \epsilon$ e riscriviamo coi nostri valori:
	\[
	\left| \frac{n}{n+1} - 1 \right| < \epsilon \implies \left| \frac{n - (n+1)}{n+1} \right| < \epsilon \implies \left| \frac{-1}{n+1} \right| < \epsilon
	\]
	Togliamo il modulo e risolviamo:
	\[
	\frac{1}{n+1} < \epsilon \implies n+1 > \frac{1}{\epsilon} \implies n > \frac{1}{\epsilon} - 1 \in \R
	\]
	Per confrontarlo con $n$, arrotondiamolo per eccesso, cioè prendiamone la \textbf{parte intera} e aggiungiamo $1$:
	\[
	n_\epsilon = \left\lfloor \frac{1}{\epsilon} - 1 \right\rfloor + 1
	\]
	Ci basta dunque trovare \textbf{UN SINGOLO} $n_\epsilon$ per dimostrare che la serie sta convergendo ad $1$.
	
	Per i calcoli, possiamo dire che ``COMANDA'' il termine di grado maggiore, e se sono dello stesso grado si annullano e rimangono eventuali fattori moltiplicativi.
	\[
	\lim_{n \to +\infty} \frac{2n+1}{n} = \lim_{n \to +\infty} \frac{n\left(2 + \frac{1}{n}\right)}{n} = 2
	\]
	\begin{osservazione}{}{}
		Quando $l=0$, si dice che $a_n$ è \textbf{infinitesima}.
	\end{osservazione}
	
	\begin{itemize}
		\item Se $\lim a_n = \infty$ si dice che la successione \textbf{DIVERGE} a un generico infinito (positivo o negativo).
		Si formalizza come:
		\[
		\forall K > 0 \quad \exists n_K \in \N \text{ tale che } |a_n| > K \quad \forall n > n_K
		\]
		ovvero $a_n < -K \quad \lor \quad a_n > K$.
		In soldoni, ciò vuol dire che, dopo un certo indice $n_K$, tutti i numeri prodotti dalla successione saranno minori dell'inverso di $K$ (limite $-\infty$) oppure maggiori di $K$ stesso (limite $+\infty$).
	\end{itemize}
	
	\begin{teorema}{Teorema del Confronto per successioni divergenti}{}
		Siano $a_n, b_n$ due successioni. Poniamo $a_n \ge b_n \quad \forall n \ge n_0$ ($a_n$ è maggiore di $b_n$ da un certo indice $n_0$ in poi).\\
		Se $\lim_{n \to +\infty} b_n = +\infty$ allora anche $\lim_{n \to +\infty} a_n = +\infty$.
	\end{teorema}
	
	\begin{osservazione}{}{}
		Quando \textbf{ESISTE} il limite di una successione, sia esso finito o infinito, la successione si dice \textbf{REGOLARE}. Ad esempio $a_n = n$ è regolare, mentre $a_n = (-1)^n$ non lo è in quanto i valori oscillano tra $\pm 1$ (indeterminata).
	\end{osservazione}
	
	\subsection{Gerarchia degli Infiniti}
	Alcune successioni divergono più velocemente di altre. Siano $a_n, b_n$ divergenti. Se:
	\[
	\lim_{n \to +\infty} \frac{b_n}{a_n} = +\infty
	\]
	si dice che $b_n$ ha ordine di infinito superiore ad $a_n$ ($a_n \ll b_n$).
	La gerarchia fondamentale è la seguente:
	\[
	(\log_B n)^\beta \ll n^\alpha \ll A^n \ll n! \ll n^n
	\]
	(con $B$ base e $\beta$ esponente qualsiasi).
	
	\subsection{Algebra dei Limiti}
	Siano $a_n \to l \in \R, \quad b_n \to m \in \R$ convergenti, allora:
	\begin{itemize}
		\item $a_n \pm b_n \to l \pm m$
		\item $a_n \cdot b_n \to l \cdot m$
		\item $\frac{a_n}{b_n} \to \frac{l}{m}$ (con $b_n, m \neq 0$)
	\end{itemize}
	
	\begin{teorema}{Unicità del limite}{}
		Una successione convergente non può avere due limiti diversi.
	\end{teorema}
	\begin{dimostrazione}{}{}
		Sia $a_n$ convergente a $l, l' \in \R$. Partiamo dalla nostra tesi: $l = l'$.\\
		Per la definizione di limite sappiamo che $\forall \epsilon > 0$:
		\begin{itemize}
			\item $\exists n_\epsilon \in \N \text{ t.c. } |a_n - l| < \epsilon \quad \forall n \ge n_\epsilon$
			\item $\exists n'_\epsilon \in \N \text{ t.c. } |a_n - l'| < \epsilon \quad \forall n \ge n'_\epsilon$
		\end{itemize}
		Applichiamo la disuguaglianza triangolare:
		\[
		|l - l'| = |l - a_n + a_n - l'| \le |l - a_n| + |a_n - l'| < \epsilon + \epsilon \le 2\epsilon
		\]
		Concludiamo dunque che $0 \le |l - l'| \le 2\epsilon$ e che dunque $|l - l'| = 0 \iff l = l'$.
	\end{dimostrazione}
	
	\begin{teorema}{Prodotto di infinitesima e limitata}{}
		Date due successioni $a_n$ e $b_n$, se $a_n \to 0$ e $b_n$ è LIMITATA, allora $a_n \cdot b_n = 0$.
	\end{teorema}
	
	\begin{teorema}{Relazione tra convergenza e limitatezza}{}
		
		Una successione convergente è \textbf{anche limitata}, ma \textbf{non} vale il contrario.
		\smallbreak
		Ad esempio la successione $a_n = (-1)^n$ è limitata, perché oscilla tra $-1$ ed $1$, però non ammette limite in quanto non è regolare.
		
	\end{teorema}
	
	\subsection{Forme indeterminate di limiti}
	
	Le \textbf{forme indeterminate} sono particolari situazioni che si verificano nelle operazioni tra i limiti, e sono le seguenti casistiche:
	\[
		\infty - \infty,\,\, \frac{\infty}{\infty},\,\, \frac{0}{0},\,\, 0 \cdot \infty
	\]
	Ne è stata delineata una particolare categoria di limiti proprio perché si tratta di casistiche che non hanno un metodo risolutivo comune e il tutto sta nel capire quale sia la successione che "tende al suo valore più velocemente".
	Vediamo come si risolvono alcuni particolari casi:
	
	\begin{itemize}
		
		\item \textbf{Entrambe convergono}: $\begin{cases}
			a_n \to l \in \R\\
			b_n \to m \in \R
		\end{cases}
		a_n \pm b_n \to l \pm m$
		
		\item \textbf{Una converge e una diverge}: $\begin{cases}
			a_n \to \pm \infty\\
			b_n \to m \in \R
		\end{cases}
		a_n \pm b_n \to \pm \infty$
		
		\item \textbf{Entrambe divergono}: $\begin{cases}
			a_n \to \infty\\
			b_n \to \infty
		\end{cases}
		a_n + b_n \to \infty$
		
	\end{itemize}
	
	Questi sono i casi \textit{non patologici}, cioè in cui non è necessario verificare forme indeterminate.
	\smallbreak
	Come già detto infatti, non è possibile delineare un metodo risolutivo generale, ed è dunque necessario sfruttare le proprietà algebriche dei limiti ed in particolare \textbf{della gerarchia degli infiniti} per capire quale successione ha un limite più "forte".
	
	\begin{teorema}{Monotonia del limite}{}
		
		Siano $a_n \to l \in \R, b_n \to m \in \R$
		\begin{itemize}
			
			\item Se $a_n \leq b_n,\, \forall n \in \N$ allora $l \leq m$
			
			\item Se $a_n < b_n,\, \forall n \in \N$ allora $l \leq m$
			
		\end{itemize}
		
	\end{teorema}
	
	\begin{teorema}{Teorema dei Carabinieri}{}
		Siano $a_n, b_n, c_n$ successioni convergenti.
		Se $a_n \le b_n \le c_n \quad \forall n \in \N$ e se $\lim a_n = \lim c_n = l$, allora $\lim b_n = l$.
	\end{teorema}
	
	I teoremi appena enunciati risulteranno utilissimi nel calcolo dei limiti di successioni, in particolare delle forme indeterminate. Ad esempio, il teorema dei carabinieri risulta molto utile quando si vuole calcolare il limite di una successione che non si conosce, ma si sa essere minore/maggiore di altre due successioni.
	
	\begin{definizione}{Crescenza e decrescenza di una succession}{}
		
		\begin{itemize}
			
			\item $a_n$ è \textbf{strettamente crescente} se $a_n < a_{n+1} \forall n \in \N$
			
			\item $a_n$ è \textbf{strettamente decrescente} se $a_n > a_{n+1} \forall n \in \N$
			
			\item $a_n$ è \textbf{crescente} se $a_n \leq a_{n+1} \forall n \in \N$
			
			\item $a_n$ è \textbf{decrescente} se $a_n \geq a_{n+1} \forall n \in \N$
			
		\end{itemize}
		
	\end{definizione}
	
	\begin{teorema}{Monotonia di una successione}{}
		
		Le successioni crescenti o decrescenti si dicono successioni \textbf{monotone} e sono regolari, quindi convergono o divergono, dunque \textbf{non possono essere indeterminate}.
		
	\end{teorema}
	
	\begin{teorema}{Regolarità di una successione}{}
		
		Le successioni monotone sono \textbf{regolari} se:
		\begin{itemize}
			
			\item Con $a_n$ crescente, $\lim_{n \to +\infty} a_n = \sup(a_n)$
			
			\item Con $a_n$ decrescente, $\lim_{n \to +\infty} a_n = \inf(a_n)$
			
		\end{itemize}
		
	\end{teorema}
	
	Forniamo adesso lo strumento \textbf{più potente} nel calcolo dei limiti di successioni: i \textbf{limiti notevoli}. Di seguito lasciamo i più importanti:
	\begin{itemize}
		
		\item \textbf{Limiti fondamentali ed esponenziali:}
		\[
		\lim_{n \to +\infty} \left(1 + \frac{1}{a_n}\right)^{a_n} = e 
		\qquad \qquad 
		\lim_{n \to +\infty} \left( \frac{e^{a_n} - 1}{a_n} \right) = 1
		\]
		
		\item \textbf{Limiti goniometrici classici e derivati:}
		\[
		\begin{aligned}
			\lim_{n \to +\infty} \left( \frac{\sin(a_n)}{a_n} \right) = 1 
			\quad &\longrightarrow \quad 
			\lim_{n \to +\infty} \left( \frac{\tan(a_n)}{a_n} \right) = 1 \\[1.5ex]
			\lim_{n \to +\infty} \left( \frac{\arcsin(a_n)}{a_n} \right) = 1 
			\quad &= \quad 
			\lim_{n \to +\infty} \left( \frac{\arctan(a_n)}{a_n} \right) = 1
		\end{aligned}
		\]
		
		\item \textbf{Limiti del coseno e relative conseguenze:}
		\[
		\lim_{n \to +\infty} \left( \frac{1 - \cos(a_n)}{(a_n)^2} \right) = \frac{1}{2}
		\]
		da cui derivano direttamente:
		\begin{align*}
			\longrightarrow \quad &\lim_{n \to +\infty} \left( \frac{1 - \cos(a_n)}{a_n} \right) = 0 \\[1ex]
			\longrightarrow \quad &\lim_{n \to +\infty} \left( \frac{(a_n)^2}{1 - \cos(a_n)} \right) = 2
		\end{align*}
		
	\end{itemize}
	
	Un limite notevole fondamentale è quello rappresentato dalla \textbf{formula di Stirling}:
	\[
		\lim_{n \to +\infty} \frac{n!}{n^n \cdot e^{-n} \cdot \sqrt{2 \pi n}} = 1
	\]
	Grazie a questo limite (che non dimostreremo) si possono risolvere limiti che coinvolgono contemporaneamente successioni antipatiche come il fattoriale e successioni esponenziali.
	
	\begin{teorema}{Radici $n$-esime di un polinomio}{}
		
		Il limite della radice $n$-esima di un polinomio vale \textbf{sempre} $1$.
		\[
			\lim_{n \to +\infty} \sqrt[n]{\text{polinomio}} = 1 
		\]
		
	\end{teorema}
	
	\chapter{Limiti di funzioni}
	
	Abbiamo visto come le successioni non sono altro che una versione molto semplice di una funzione, avendo come dominio un insieme non denso come $\N$.\\
	Lo studio delle successioni e di tutte le operazioni a loro associate rappresenta un fantastico banco di prova per lo studio delle funzioni, in quanto i discorsi fatti sono analoghi ma necessitano di alcune precisazioni (dato che stiamo lavorando su un insieme \textit{denso e continuo} e non con dei buchi come $\N$).
	
	\begin{definizione}{Funzione infinitesima a $+\infty$}{}
		Sia $d: \R \to \R$. Si dice \textbf{infinitesima} se, per $x \to +\infty$,
		\[
			\forall \epsilon > 0 \quad \exists M \text{ t.c. } |d(x)| < \epsilon \quad \forall x > M
		\]
	\end{definizione}
	
	Rivedendo anche la definizione di limite di successione, ci ricordiamo che, se $a_n \to l$, allora $\forall \epsilon > 0 \quad |a_n - l| < \epsilon$.
	
	Rivedendo la precedente definizione, deduciamo che $d(x)$ è infinitesima se 
	\[
		\lim_{x \to +\infty} d(x) = 0
	\]
	
	Analogamente, se $x \to -\infty$ la definizione non cambia molto:
	\[
		\forall \epsilon > 0 \quad \exists M \text{ t.c. } |d(x)| < \epsilon \quad \forall x < M
	\]
	
	Stavolta sia $f(x): \R \to \R$ con limite $l$ per $x \to +\infty$. Se $f(x) - l$ è infinitesima, allora...
	\[
		\forall \epsilon > 0 \quad \exists M \text{ t.c. } |f(x) - l| < \epsilon \quad \forall x > M
	\]
	
	Sia $f: (a, b) \to \R$, $x_0 \in [a, b]$ (1)
	
	$f$ è \textbf{infinitesima} per $x \to x_0$ se $\forall \epsilon > 0 \quad \exists \delta > 0 \text{ t.c. } 0 < |x - x_0| < \delta$ con $x \in (a, b)$ e dunque $|f(x)| < \epsilon$
	
	Sia $f: (a, b) \to \R$, $x_0 \in [a, b]$ (2)
	
	\[
	\lim_{x \to x_0} f(x) = l \in \R \quad \text{sse} \quad f(x) - l \text{ è infinitesima per } x \to x_0
	\]
	quindi $\forall \epsilon > 0 \quad \exists \delta > 0 \text{ t.c. } 0 < |x - x_0| < \delta \quad \text{e} \quad |f(x) - l| < \epsilon$
	
	\vspace{1em}
	
	La (1) ci descrive quando una funzione è \textbf{infinitesima}.
	\[
	f(x) \text{ infinitesima per } x \to x_0 \iff \lim_{x \to x_0} f(x) = 0
	\]
	
	La (2) ci descrive quando una funzione ha limite usando il \textbf{concetto di infinitesima}.
	\[
	\lim_{x \to x_0} f(x) = l \iff f(x) - l \text{ infinitesima per } x \to x_0
	\]
	
	Valgono le stesse operazioni aritmetiche definite per i limiti di successione.
	
	Siano $\lim_{x \to x_0} f(x) = l$, $\lim_{x \to x_0} g(x) = m$, allora:
	
	\begin{itemize}
		
		\item $\lim_{x \to x_0} f(x) \pm g(x) = l \pm m
		\qquad \bullet \,\lim_{x \to x_0} f(x) \cdot g(x) = l \cdot m$
		
		\vspace{0.5em}
		
		\item $\lim_{x \to x_0} \frac{f(x)}{g(x)} = \frac{l}{m} \quad \text{sse } m \neq 0$
	\end{itemize}
	
	Valendo queste proprietà, le strategie di calcolo per i limiti più semplici sono le stesse utilizzate per le successioni.
	
	Se $\exists \lim_{x \to x_0} f(x)$ e $\exists \lim_{x \to x_0} g(x)$ e $f(x) \le g(x)$ allora $l \le m$.
	
	Vale il \textbf{teorema dei carabinieri}.
	
	\subsection{Limite destro e sinistro}
	
	Dato l'intervallo $(x_0 - \delta, x_0 + \delta)$:
	
	\begin{itemize}
		
		\item Diciamo \textbf{limite destro} il $\lim_{x \to x_0^+} f(x) = l$ se vale la definizione di limite $\forall x$ con $x_0 < x < x_0 + \delta$
		
		\item Diciamo \textbf{limite sinistro} il $\lim_{x \to x_0^-} f(x) = l$ se vale la definizione di limite $\forall x$ con $x_0 - \delta < x < x_0$
		
	\end{itemize}
	
	Si dice che $f(x)$ \textbf{ammette limite} per $x \to x_0$ se:
	\[
		\lim_{x \to x_0^+} f(x) = \lim_{x \to x_0^-} f(x)
	\]
	
	\begin{teorema}{Teorema del ponte}{}
		Siano $f: (a, b) \to \R$ e $\lim_{x \to x_0} f(x) = l \in \R \cup \{\pm\infty\}$.
		
		\smallbreak
		Dato che $f(x)$ ammette limite allora:
		\[
			\forall \text{successione } a_n \subseteq [a, b] \setminus \{x_0\} \text{ con } a_n \to x_0 \quad \text{allora } \lim_{n \to +\infty} f(a_n) = l
		\]
		
	\end{teorema}
	
	\chapter{Funzioni continue}
	
	\begin{definizione}{Funzione continua}{}
		Sia $I$ un intervallo, sia $f: I \to \R$ e sia $x_0 \in I$.\\
		Diciamo che $f$ è \textbf{continua} in $x_0$ se $\lim_{x \to x_0} f(x) = f(x_0)$.
	\end{definizione}
	
	Ad esempio, data $f(x) = 3x^2$ e $x_0 = 1$, si ha $f(1) = 3$ e $\lim_{x \to 1} 3x^2 = 3$.
	
	\begin{osservazione}{}{}
		
		\begin{itemize}
			
			\item Tutte le funzioni polinomiali sono continue.
			
			\item Anche le funzioni razionali sono continue, a meno che $x_0$ non sia uno zero del polinomio al denominatore.
			
		\end{itemize}
		
	\end{osservazione}
	
	Se $f$ è continua $\forall x_0 \in I$, allora diciamo che $f$ è \textbf{continua} sull'intervallo $I$.
	
	I punti in cui $f$ non è continua sono detti \textbf{punti di discontinuità} e sono di tre tipi:
	
	\begin{enumerate}
		
		\item \textbf{Discontinuità di salto}: quando $\lim_{x \to x_0^+} f(x) \neq \lim_{x \to x_0^-} f(x)$. Una particolarità si ha nelle funzioni definite a tratti, ad esempio:
		\[
		f(x) = 
		\begin{cases}
			2 & \text{se } x \geq 0 \\
			1 & \text{se } x < 0
		\end{cases}
		\]
		In questo caso il limite sinistro tende a $1$ mentre il limite destro tende a $2$.
		
		\item \textbf{Discontinuità di seconda specie}: quando almeno uno dei due limiti laterali non esiste oppure è uguale a $\pm \infty$. Ad esempio, per $f(x) = \frac{1}{x}$ in $x_0 = 0$ abbiamo che il limite destro tende a $+\infty$ e il limite sinistro tende a $-\infty$. Bastava che ne esistesse o ne divergesse anche solo uno per soddisfare la condizione.
		
		\item \textbf{Discontinuità di terza specie}: poniamo $\lim_{x \to x_0^+} f(x) = \lim_{x \to x_0^-} f(x) = l$. Tuttavia, in questo caso o $f$ non è stata definita in $x_0$, oppure $f(x_0) \neq l$.
		In quest'ultimo caso possiamo però \textbf{estendere per continuità} la funzione $f$, definendo una nuova funzione:
		\[
		\tilde{f}(x) = 
		\begin{cases}
			f(x) & \text{se } x \neq x_0 \\
			l & \text{se } x = x_0
		\end{cases}
		\]
		facendo sì che venga definita in $x_0$ con il valore $l$.
		
	\end{enumerate}

	\begin{teorema}{Teorema di esistenza degli zeri}{}
		
		Sia $f: [a, b] \to \R$ continua. Poniamo $f(a) \cdot f(b) < 0$ (dunque $f(a)$ e $f(b)$ hanno segno diverso).
		Allora esiste $c \in (a, b)$ tale che $f(c) = 0$.
		
	\end{teorema}
	
	\begin{dimostrazione}{Dimostrazione}{}
		
		Supponiamo $f(a) < 0$ e $f(b) > 0$. Costruiamo due successioni $a_n$ e $b_n$. Poniamo $a_0 = a$ e $b_0 = b$.\\
		Troviamo $c = \frac{a+b}{2}$ come punto medio tra $a$ e $b$. Analizziamo i possibili casi:
		\begin{itemize}
			\item Se $f(c) > 0$, poniamo $a_1 = a$ e $b_1 = c$ (dato che $f(b) > 0$).
			\item Se $f(c) < 0$, poniamo $a_1 = c$ e $b_1 = b$ (dato che $f(a) < 0$).
			\item Se $f(c) = 0$, la dimostrazione è già conclusa perché viene rispettata la condizione del teorema.
		\end{itemize}
		Stiamo costruendo man mano intervalli sempre più piccoli. Dunque $|a_n - b_n| \to 0$, e avendo definito $a_n$ crescente e $b_n$ decrescente, ed essendo entrambe regolari, ammettono lo stesso limite $d$.
		Passando al limite e sfruttando la continuità di $f$:
		\[
		\begin{cases}
			f(a_n) < 0 \\
			f(b_n) > 0
		\end{cases}
		\xrightarrow{f \text{ è continua}}
		\begin{cases}
			f(d) \leq 0 \\
			f(d) \geq 0
		\end{cases}
		\implies f(d) = 0
		\]
		
	\end{dimostrazione}
	
	\begin{teorema}{Teorema del valore intermedio}{}
		
		Sia $f: [a, b] \to \R$ continua.
		Allora $\forall \eta \in [f(a), f(b)]$ esiste $x_0 \in [a, b]$ tale che $f(x_0) = \eta$.
		
	\end{teorema}
	
	\begin{dimostrazione}{Dimostrazione}{}
		Definiamo la funzione ausiliaria $g(x) = f(x) - \eta$. Notiamo che $g: [a, b] \to \R$ è continua in quanto differenza di funzioni continue.\\
		Verifichiamo i segni agli estremi dell'intervallo:
		\begin{itemize}
			\item $g(a) = f(a) - \eta < 0$ perché $\eta > f(a)$
			\item $g(b) = f(b) - \eta > 0$ perché $\eta < f(b)$
		\end{itemize}
		Applichiamo il teorema di esistenza degli zeri alla funzione $g(x)$: esiste $x_0 \in (a, b)$ tale che $g(x_0) = 0$, ovvero $f(x_0) - \eta = 0$, da cui $f(x_0) = \eta$.
	\end{dimostrazione}
	
	\begin{teorema}{Teorema di Weierstrass}{}
		
		Sia $f: [a, b] \to \R$ continua.\\
		Allora esistono $x_0, x_1 \in [a, b]$ tali che $f(x_0) \leq f(x) \leq f(x_1)$ per ogni $x \in [a, b]$, dove $x_0$ è il punto di minimo e $x_1$ è il punto di massimo.
		
	\end{teorema}
	
	Le ipotesi del teorema (intervallo chiuso e limitato, e continuità della funzione) sono fondamentali. Se viene meno anche una sola di queste ipotesi, la tesi potrebbe non valere, come mostrato nei seguenti esempi.
	
	\begin{osservazione}{}{}
		
		Se consideriamo una funzione continua definita su un intervallo aperto, il teorema di Weierstrass può non valere. Sia $f(x) = x^2$ con $f: (-1, 1) \to \R$. L'immagine è $f((-1, 1)) = [0, 1)$, dunque esiste il minimo ma non esiste un punto di massimo in cui la funzione assume il valore $1$.
		
	\end{osservazione}
	
	\begin{osservazione}{}{}
		
		Se la funzione non è continua, il teorema non garantisce l'esistenza di massimi o minimi. Ad esempio, la funzione definita a tratti $f(x) = x^2$ per $x > 0$ e $f(x) = 1$ per $x \leq 0$, definita sull'intervallo $[-1, 1]$, non ammette un punto di minimo in quanto $f([-1, 1]) = (0, 1]$.
		
	\end{osservazione}
	
	Combinando la proprietà di continuità con il teorema di Weierstrass, possiamo determinare l'immagine di un intervallo chiuso e limitato tramite una funzione continua.
	
	\begin{proposizione}{}{}
		
		Sia $f: [a, b] \to \R$ una funzione continua. Allora l'immagine dell'intervallo è anch'essa un intervallo chiuso e limitato, i cui estremi sono il minimo e il massimo della funzione:
		\[
		f([a, b]) = [f(m), f(M)]
		\]
		dove $m$ è il punto di minimo e $M$ è il punto di massimo.
		
	\end{proposizione}
	
	Se la funzione è definita su insiemi non limitati o non chiusi, l'immagine potrebbe non essere un intervallo chiuso. Ad esempio, per la funzione identità $f(x) = x$ con $f: \R \to \R$, la funzione non ha punti di massimo o minimo assoluto in quanto $\lim_{x \to +\infty} x = +\infty$ e $\lim_{x \to -\infty} x = -\infty$, dunque l'immagine è l'intero insieme $\R$.
	
	\begin{proposizione}{}{}
		
		Sia $P(x)$ un polinomio di grado dispari. Allora esiste $x_0 \in \R$ tale che $P(x_0) = 0$.
		
	\end{proposizione}
	
	\begin{dimostrazione}{}{}
		
		Sia $P(x) = a_n x^n + a_{n-1} x^{n-1} + \dots + a_0$, con $a_n \neq 0$.
		Valutiamo i limiti agli estremi:
		\[
		\lim_{x \to +\infty} P(x) = \begin{cases}
			+\infty, &\text{se } a_n > 0 \\
			-\infty, &\text{se } a_n < 0
		\end{cases}
		\qquad
		\lim_{x \to -\infty} P(x) = \begin{cases}
			-\infty, &\text{se } a_n > 0 \\
			+\infty, &\text{se } a_n < 0
		\end{cases}
		\]
		
		\textbf{Caso $a_n > 0$:} esistono $L_+$ tale che $P(L_+) > 0$ ed $L_-$ tale che $P(L_-) < 0$. 
		Possiamo ora considerare la restrizione $P: [L_-, L_+] \to \R$. Per il teorema degli zeri, esiste $c \in [L_-, L_+]$ tale che $P(c) = 0$.
		
		\textbf{Caso $a_n < 0$:} esistono $L_+$ tale che $P(L_+) < 0$ ed $L_-$ tale che $P(L_-) > 0$. 
		Analogamente, considerando $P: [L_-, L_+] \to \R$, esiste $c \in [L_-, L_+]$ tale che $P(c) = 0$.
		
	\end{dimostrazione}
	
	\section{Serie geometrica finita}
	
	Definiamo la somma parziale di una serie geometrica:
	\[
	\sum_{i=0}^{n} q^i = q^0 + q^1 + q^2 + \dots + q^n = \frac{1 - q^{n+1}}{1 - q} \quad \text{per } q \in \R \setminus \{1\}
	\]
	dove il simbolo $\sum$ indica la somma finita di $n+1$ termini.
	
	\chapter{Derivate}
	
	La derivata di una funzione $f$ in un punto $x_0$ è definita come:
	
	\begin{definizione}{Rapporto incrementale}{}
		
		Dato il rapporto incrementale $\frac{\Delta f}{\Delta x} = \frac{f(x_1) - f(x_0)}{x_1 - x_0}$, il risultato indica la pendenza ($m$) della retta secante passante per i punti $(x_0, f(x_0)), (x_1, f(x_1))$.
		
	\end{definizione}
	
	L'equazione della retta è data da $\frac{\Delta f}{\Delta x} (x - x_0) + f(x_0)$.
	
	Se consideriamo due punti $x, x_0$ molto vicini tra loro ($x - x_0 \to 0$), allora possiamo arrivare a definire una retta tangente alla funzione $f(x)$ e il limite per calcolare tale pendenza:
	\[
	f'(x_0) = \lim_{x \to x_0} \frac{f(x) - f(x_0)}{x - x_0}
	\]
	dove con $f'(x_0)$ indichiamo la derivata prima.
	
	Se questo limite esiste, si dice che $f$ è \textbf{derivabile} in $x_0$.
	Se $f$ è derivabile $\forall x_0 \in \operatorname{dom}(f)$, si dice che $f$ è derivabile.
	
	Sostituendo $h = x - x_0$:
	\[
	f'(x_0) = \lim_{h \to 0} \frac{f(x_0 + h) - f(x_0)}{h}
	\]
	
	Definiamo la \textbf{derivata destra}:
	\[
	\lim_{h \to 0^+} \frac{f(x_0 + h) - f(x_0)}{h} = f'_+(x_0)
	\]
	
	Definiamo la \textbf{derivata sinistra}:
	\[
	\lim_{h \to 0^-} \frac{f(x_0 + h) - f(x_0)}{h} = f'_-(x_0)
	\]
	
	Come per i limiti, diciamo che $f$ è derivabile in $x_0$ sse $f'_- = f'_+$.
	
	\newpage
	\section{Derivabilità e continuità}
	
	\begin{proposizione}{}{}
		
		Se $f$ è derivabile in $x_0$, allora $f$ è anche continua in $x_0$.
		
	\end{proposizione}
	
	\begin{dimostrazione}{}{}
		
		Vogliamo dimostrare che $\lim_{x \to x_0} f(x) = f(x_0)$, ovvero che $\lim_{x \to x_0} (f(x) - f(x_0)) = 0$.
		Partiamo dal limite del rapporto incrementale e riscriviamolo opportunamente:
		\[
		\lim_{x \to x_0} \frac{f(x) - f(x_0)}{x - x_0} \cdot (x - x_0)
		\]
		Passando al limite per $x \to x_0$, il primo fattore tende a $f'(x_0)$ (essendo $f$ derivabile), mentre il secondo fattore $(x - x_0)$ tende a $0$. Pertanto il loro prodotto fa $0$, il che prova che $\lim_{x \to x_0} (f(x) - f(x_0)) = 0$.
		
	\end{dimostrazione}
	
	\begin{osservazione}{}{}
		
		Non vale il contrario: una funzione può essere continua in un punto ma non derivabile. Ad esempio, la funzione $f(x) = |x|$ è continua in $x_0 = 0$, ma non è derivabile nello stesso punto in quanto il limite destro e il limite sinistro del rapporto incrementale non coincidono.
		
	\end{osservazione}
	
	\section{Notazioni e derivate fondamentali}
	
	Le notazioni più comuni per indicare le derivate successive di una funzione sono le seguenti:
	\[
	f', \, f'', \, f''', \, f^{(n)}, \, \frac{df}{dx}
	\]
	dove $f', f'', f'''$ indicano rispettivamente la prima, la seconda e la terza derivata, $f^{(n)}$ la derivata $n$-esima e $\frac{df}{dx}$ la notazione di Leibniz per la derivata prima.
	
	Vediamo ora alcune dimostrazioni delle derivate delle funzioni elementari:
	
	\begin{enumerate}
		
		\item \textbf{Derivata della funzione costante}: se $f(x) \equiv c$, allora $f'(x) = 0$.
		\begin{dimostrazione}{}{}
			
			Infatti, applicando la definizione:
			\[
			f'(x_0) = \lim_{h \to 0} \frac{f(x_0 + h) - f(x_0)}{h} = \frac{c - c}{h} = 0
			\]
			Notiamo che non si tratta di una forma indeterminata $\frac{0}{0}$ in quanto il numeratore è \emph{esattamente} $0$ per ogni $h$.
			
		\end{dimostrazione}
		
		\item \textbf{Derivata della funzione potenza}: se $f(x) = x^n$, allora $f'(x) = n \cdot x^{n-1}$.
		\begin{dimostrazione}{}{}
			
			Scriviamo il rapporto incrementale:
			\[
			f'(x) = \lim_{x \to x_0} \frac{x^n - x_0^n}{x - x_0}
			\]
			Riconoscendo la formula algebrica della scomposizione (collegata alla serie geometrica), possiamo riscrivere il limite come:
			\[
			\lim_{x \to x_0} \left[ x^{n-1} + x_0 \cdot x^{n-2} + x_0^2 \cdot x^{n-3} + \dots + x_0^{n-1} \right]
			\]
			Essendoci $n$ termini e tendendo ciascuno di essi a $x_0^{n-1}$, otteniamo complessivamente:
			\[
			f'(x_0) = n \cdot x_0^{n-1}
			\]
			
		\end{dimostrazione}
		
		\item \textbf{Derivata di potenze ad esponente negativo o reale}:
		Con dimostrazioni analoghe alle precedenti si ricava che per $f(x) = x^{-b} = \frac{1}{x^b}$ si ha $f'(x) = -b \cdot x^{-b-1}$, e più in generale per $f(x) = x^\alpha$ si ha $f'(x) = \alpha \cdot x^{\alpha-1}$ per ogni $\alpha \in \R$.
	
		\item Sia $f(x) = \sin(x)$, allora $f'(x) = \cos(x)$.
		\begin{dimostrazione}{Derivata del seno}{}
			
			Applicando la definizione di derivata tramite il rapporto incrementale in un punto $x_0$:
			\[
			f'(x_0) = \lim_{h \to 0} \frac{f(x_0 + h) - f(x_0)}{h} = \lim_{h \to 0} \frac{\sin(x_0 + h) - \sin(x_0)}{h}
			\]
			Sfruttando la formula di addizione del seno:
			\[
			= \lim_{h \to 0} \frac{\sin(x_0)\cos(h) + \cos(x_0)\sin(h) - \sin(x_0)}{h} =
			\]
			\[
			= \lim_{h \to 0} \left[ \sin(x_0) \underbrace{\left(\frac{\cos(h) - 1}{h}\right)}_{\to 0} + \cos(x_0) \underbrace{\left(\frac{\sin(h)}{h}\right)}_{\to 1} \right] = \cos(x_0)
			\]
			
		\end{dimostrazione}
		
		\item Sia $f(x) = \cos(x)$, allora $f'(x) = -\sin(x)$.
		\begin{dimostrazione}{}{}
			
			Procedendo in modo analogo per il coseno:
			\[
			f'(x_0) = \lim_{h \to 0} \frac{\cos(x_0 + h) - \cos(x_0)}{h} =
			\]
			\[
			= \lim_{h \to 0} \frac{\cos(x_0)\cos(h) - \sin(x_0)\sin(h) - \cos(x_0)}{h} =
			\]
			\[
			= \lim_{h \to 0} \left[ \cos(x_0) \underbrace{\left(\frac{\cos(h) - 1}{h}\right)}_{\to 0} - \sin(x_0) \underbrace{\frac{\sin(h)}{h}}_{\to 1} \right] = -\sin(x_0)
			\]
			
		\end{dimostrazione}
		
		\item Sia $f(x) = e^x$, allora $f'(x) = e^x$.
		\begin{dimostrazione}{}{}
			
			Utilizzando il rapporto incrementale:
			\[
			f'(x_0) = \lim_{h \to 0} \frac{e^{x_0 + h} - e^{x_0}}{h} = \frac{e^{x_0}(e^h - 1)}{h} = e^{x_0} \underbrace{\left(\frac{e^h - 1}{h}\right)}_{\to 1} = e^{x_0}
			\]
			
		\end{dimostrazione}
		
		\item Sia $f(x) = \ln(x)$, allora $f'(x) = \frac{1}{x}$.
		\begin{dimostrazione}{}{}
			
			Impostando il rapporto incrementale e sfruttando le proprietà dei logaritmi:
			\[
			f'(x_0) = \lim_{h \to 0} \frac{\ln(x_0 + h) - \ln(x_0)}{h} = \lim_{h \to 0} \frac{\ln\left(\frac{x_0 + h}{x_0}\right)}{h} =
			\]
			\[
			= \lim_{h \to 0} \frac{\ln\left(1 + \frac{h}{x_0}\right)}{h} = \frac{1}{x_0}
			\]
			
		\end{dimostrazione}
		
		\item Sia $f(x) = \tan(x)$, allora $f'(x) = \frac{1}{\cos^2(x)}$.
		
		\item Sia $f(x) = \arcsin(x)$, allora $f'(x) = \frac{1}{\sqrt{1 - x^2}}$.
		
		\item Sia $f(x) = \arccos(x)$, allora $f'(x) = -\frac{1}{\sqrt{1 - x^2}}$.
		
		\item Sia $f(x) = \log_a(x)$, allora $f'(x) = \frac{1}{x \ln a}$.
		
	\end{enumerate}
	
	\section{Regole di derivazione}
	
	Vengono di seguito enunciate e dimostrate le principali regole di derivazione per le funzioni reali di variabile reale.
	
	\begin{itemize}
		
		\item \textbf{Linearità}: $(\alpha \cdot f(x) + \beta \cdot g(x))' = \alpha \cdot f'(x) + \beta \cdot g'(x)$.
		\begin{dimostrazione}{}{}
			
			La dimostrazione sfrutta la linearità del limite:
			\[
			\lim_{h \to 0} \frac{\alpha \cdot f(x+h) + \beta \cdot g(x+h) - \alpha \cdot f(x) - \beta \cdot g(x)}{h} =
			\]
			\[
			= \lim_{h \to 0} \left( \alpha \cdot \frac{f(x+h) - f(x)}{h} + \beta \cdot \frac{g(x+h) - g(x)}{h} \right) = \alpha f'(x) + \beta g'(x)
			\]
			
		\end{dimostrazione}
		
		\item \textbf{Prodotto}: $(f(x)g(x))' = f'(x)g(x) + f(x)g'(x)$.
		\begin{dimostrazione}{}{}
			
			Partiamo dal rapporto incrementale e sommiamo e sottraiamo opportunamente un termine al numeratore:
			\[
			\lim_{h \to 0} \frac{f(x+h)g(x+h) - f(x)g(x)}{h} =
			\]
			\[
			= \lim_{h \to 0} \frac{f(x+h)g(x+h) - f(x)g(x) + f(x+h)g(x) - f(x+h)g(x)}{h} =
			\]
			Raccogliendo opportunamente i termini, otteniamo:
			\[
			= \lim_{h \to 0} \left( f(x+h) \frac{g(x+h) - g(x)}{h} + g(x) \frac{f(x+h) - f(x)}{h} \right)
			\]
			\[
				= f(x)g'(x) + g(x)f'(x)
			\]
			
		\end{dimostrazione}
		
		\item \textbf{Rapporto}: $\left(\frac{f(x)}{g(x)}\right)' = \frac{f'(x)g(x) - f(x)g'(x)}{[g(x)]^2}$.
		\begin{dimostrazione}{}{}
			
			Impostiamo il limite del rapporto incrementale:
			\[
			\lim_{h \to 0} \frac{\frac{f(x+h)}{g(x+h)} - \frac{f(x)}{g(x)}}{h} = \lim_{h \to 0} \frac{1}{h} \left[ \frac{f(x+h)g(x) - f(x)g(x+h)}{g(x+h)g(x)} \right] =
			\]
			Aggiungendo e sottraendo $f(x)g(x)$ al numeratore e passando al limite, per la continuità di $g$ si ottiene la tesi:
			\[
			= \frac{f'(x)g(x) - f(x)g'(x)}{[g(x)]^2}
			\]
			
		\end{dimostrazione}
		
	\end{itemize}
	
	\section{Derivata di funzione composta}
	
	La derivata di una funzione composta si calcola con la seguente formula:
	\[
	(f \circ g(x))' = f'(g(x)) \cdot g'(x)
	\]
	
	\begin{dimostrazione}{}{}
		
		Definiamo una funzione continua $F(t)$:
		\[
		F(t) = \begin{cases}
			\frac{f(t) - f(g(x))}{t - g(x)} & \text{per } t \neq g(x) \\
			f'(g(x)) & \text{per } t = g(x)
		\end{cases}
		\]
		La derivata della funzione composta può essere scritta come il rapporto incrementale:
		\[
		\frac{f(g(x+h)) - f(g(x))}{h} = F(g(x+h)) \cdot \frac{g(x+h) - g(x)}{h}
		\]
		Suddividiamo nei casi di definizione di $F$:
		\begin{itemize}
			
			\item \textbf{Caso $g(x+h) \neq g(x)$}: a destra dell'uguale l'espressione diventa:
			\[
			\frac{f(g(x+h)) - f(g(x))}{g(x+h) - g(x)} \cdot \frac{g(x+h) - g(x)}{h} = \frac{f(g(x+h)) - f(g(x))}{h}
			\]
			verificato.
			
			\item \textbf{Caso $g(x+h) = g(x)$}:
			\begin{itemize}
				
				\item \textit{Destra}: $F(g(x)) \cdot 0 = 0$
				
				\item \textit{Sinistra}: $\frac{f(g(x)) - f(g(x))}{h} = 0$
				
			\end{itemize}
			
			verificato.
			
		\end{itemize}
		
	\end{dimostrazione}
	
	Possiamo usare la precedente formula per calcolare anche la derivata di funzioni inverse.
	Se $f$ è invertibile, sappiamo che $x = f(f^{-1}(x))$. Derivando tutto otteniamo:
	\[
	1 = f'(f^{-1}(x)) \cdot (f^{-1}(x))'
	\]
	
	\begin{osservazione}{}{}
		
		Vogliamo calcolare la derivata di $\ln(x)$ usando la formula per le funzioni inverse.
		Poniamo $f(x) = e^x$, da cui $f^{-1}(x) = \ln(x)$ e $f'(x) = e^x$.
		Utilizziamo la formula $1 = f'(f^{-1}(x)) \cdot (f^{-1}(x))'$ per trovare la derivata della funzione inversa:
		\[
		f'(f^{-1}(x)) = e^{\ln x} = x
		\]
		Da cui ricaviamo:
		\[
		(f^{-1}(x))' = \frac{1}{f'(f^{-1}(x))} = \frac{1}{x}
		\]
		
	\end{osservazione}
	
	\begin{proposizione}{Derivata di funzioni pari e dispari}{}
		
		Se $f$ è pari, $f'$ è dispari e viceversa.
		
	\end{proposizione}
	
	\begin{dimostrazione}{}{}
		
		\begin{itemize}
			
			\item Se $f$ è pari allora $f(x) = f(-x) \implies$ deriviamo entrambi i lati e trattiamo come funzione composta: $f'(x) = -f'(-x) \implies$ definizione di funzione dispari.
			
			\item Se $f$ è dispari allora $f(x) = -f(-x) \implies$ deriviamo entrambi i lati e trattiamo come funzione composta: $f'(x) = f'(-x) \implies$ definizione di funzione pari.
			
		\end{itemize}
		
	\end{dimostrazione}
	
	\section{Punti di non derivabilità}
	
	\begin{itemize}
		
		\item \textbf{Punti angolosi}, dove $f'_+(x_0) \neq f'_-(x_0)$.\\
		Ad esempio, consideriamo $f(x) = |x|$ in $x_0 = 0$. A sinistra la pendenza è $-1$, mentre a destra la pendenza è $+1$.
		
		\item \textbf{Tangente verticale}: $f'_+(x_0) = f'_-(x_0) = \pm \infty$ (sono infiniti e concordi).\\
		Ad esempio, consideriamo $f(x) = x^{1/3}$ in $x_0 = 0$. La derivata è $f'(x) = \frac{1}{3\sqrt[3]{x^2}}$, il cui limite per $x \to 0$ è sempre $+\infty$.
		
		\item \textbf{Punti di cuspidi}: $f'_+(x_0) = \pm \infty$ e $f'_-(x_0) = \mp \infty$ (sono infiniti e discordi).
		
	\end{itemize}
	
	\section{Intorni, massimi e minimi relativi}
	
	Se abbiamo un punto $x_0$, definiamo \textbf{intorno di $x_0$} qualsiasi intervallo $(a, b)$ che contiene $x_0$, ovvero $x_0 \in (a, b)$.
	
	\begin{itemize}
		
		\item $x_0 \in \operatorname{dom}(f)$ è un \textbf{massimo relativo} (o locale) se esiste un intorno $I_{x_0}$ di $x_0$ tale che $f(x_0) \geq f(x) \quad \forall x \in I_{x_0} \cap \operatorname{dom}(f)$.
		
		\item $x_0 \in \operatorname{dom}(f)$ è un \textbf{minimo relativo} (o locale) se esiste un intorno $I_{x_0}$ di $x_0$ tale che $f(x) \geq f(x_0) \quad \forall x \in I_{x_0} \cap \operatorname{dom}(f)$.
		
	\end{itemize}
	
	\begin{teorema}{Teorema di Fermat}{}
		
		Sia $f: (a, b) \to \R$ una funzione derivabile e sia $x_0 \in (a, b)$ un punto di \textbf{estremo} (ovvero un massimo o minimo relativo).
		Allora $f'(x_0) = 0$.
		
	\end{teorema}
	
	\begin{dimostrazione}{}{}
		
		Supponiamo che $x_0$ sia un massimo locale, ovvero $f(x_0) \geq f(x)$ per ogni $x$ in un opportuno intorno.
		
		Dobbiamo calcolare il rapporto incrementale $\frac{\Delta f}{\Delta x}$ e analizzare le due casistiche:
		
		\begin{itemize}
			
			\item Se $x > x_0$, il denominatore è positivo ($\Delta x > 0$) e il numeratore è sempre non positivo ($\Delta f \leq 0$) in quanto $x_0$ è massimo relativo. Dunque in totale abbiamo $\frac{\Delta f}{\Delta x} \leq 0$.
			
			\item Se $x < x_0$, il denominatore è negativo ($\Delta x < 0$) e il numeratore è sempre non positivo ($\Delta f \leq 0$) in quanto $x_0$ è massimo relativo. Dunque in totale abbiamo $\frac{\Delta f}{\Delta x} \geq 0$.
			
		\end{itemize}
		Combinando le due soluzioni al limite per $x \to x_0$, poichè la funzione è derivabile, si deduce che:
		\[
		f'(x_0) = 0
		\]
		Se $x_0$ fosse un minimo relativo, la dimostrazione sarebbe del tutto analoga.
		
	\end{dimostrazione}
	
	\begin{teorema}{Teorema di Rolle}{}
		
		Sia $f: [a, b] \to \R$ una funzione derivabile in $(a, b)$.
		Se $f(a) = f(b)$, allora $\exists x_0 \in (a, b)$ tale che $f'(x_0) = 0$.
		
	\end{teorema}
	
	Geometricamente, se dobbiamo partire da un punto e arrivare allo stesso livello, ci sarà per forza almeno un punto interno in cui la funzione cambia andamento (da crescente a decrescente o viceversa), ossia un punto a tangente orizzontale.
	
	\begin{center}
		
		\begin{tikzpicture}[scale=1]
			
			\draw[->] (-0.5, 0) -- (3.5, 0) node[right] {\small $x$};
			\draw[->] (0, -0.5) -- (0, 2) node[above] {\small $y$};
			
			\draw[blue, thick] plot[smooth, domain=0.5:3] (\x, {1 + 0.6*sin((\x-0.5)*180/2.5)});
			\fill[black] (0.5, 1) circle (1.5pt) node[below] {$a$};
			\fill[black] (3.0, 1) circle (1.5pt) node[below] {$b$};
			\draw[dashed] (0.5, 1) -- (3.0, 1);
			
			\draw[dashed, red] (1.75, 1.6) -- (1.75, 0) node[below] {$x_0$};
			\fill[red] (1.75, 1.6) circle (1.5pt);
			\draw[red, thick] (1.0, 1.6) -- (2.5, 1.6);
			
		\end{tikzpicture}
		
	\end{center}
	
	\begin{dimostrazione}{}{}
		
		Per il teorema di Weierstrass sappiamo che nell'immagine $f([a, b])$ esistono il minimo $m$ e il massimo $M$.
		
		\begin{description}
			
			\item [Caso $m = M$:] la funzione $f(x)$ è costante in $[a, b]$, quindi $f'(x) = 0, \, \forall x \in (a, b)$.
			
			\item [Caso $m < M$:]~
			\begin{itemize}
				
				\item Se $f(a) = f(b) = m$: secondo Weierstrass esistono minimo e massimo nell'insieme dei valori $f([a, b])$. Ed essendo $m < M$, si ha che il massimo $M$ deve essere assunto in un punto interno $x_0 \in (a, b)$. Dunque $\exists x_0 \in (a, b)$ tale che $f'(x_0) = 0$.
				
				\item Se $f(a) = f(b) = M$: la dimostrazione è analoga, in quanto il minimo $m$ sarà assunto in un punto interno dell'intervallo.
				
				\item Se invece $f(a)$ e $f(b)$ sono diversi sia da $m$ che da $M$: esistono $x_0, x_1 \in (a, b)$ tali che $f(x_0) = m$ e $f(x_1) = M$, e dunque $f'(x_0) = f'(x_1) = 0$.
			
			\end{itemize}
			
		\end{description}
		
	\end{dimostrazione}
	
	\begin{teorema}{Teorema di Lagrange}{}
		
		Sia $f: [a, b] \to \R$ una funzione continua in $[a, b]$ e derivabile in $(a, b)$.
		Allora $\exists x_0 \in (a, b)$ tale che:
		\[
		f'(x_0) = \frac{f(b) - f(a)}{b - a}
		\]
		
	\end{teorema}
	
	Graficamente, la quantità $\frac{f(b) - f(a)}{b - a}$ rappresenta la \textbf{pendenza della retta secante} passante per i punti $(a, f(a))$ e $(b, f(b))$. Il teorema ci assicura quindi che esiste almeno un punto $x_0$ la cui derivata (pendenza della retta tangente) è uguale alla pendenza della secante.
	
	\begin{dimostrazione}{}{}
		
		Consideriamo la funzione ausiliaria $g(x)$ definita come:
		\[
		g(x) := f(x) - \frac{f(b) - f(a)}{b - a} \cdot (x - a)
		\]
		Notiamo che $g \in C([a, b])$ ed è derivabile in $(a, b)$. Verifichiamo i valori della funzione agli estremi dell'intervallo:
		\begin{itemize}
			
			\item $g(a) = f(a) - 0 = f(a)$
			
			\item $g(b) = f(b) - \frac{f(b) - f(a)}{b - a} \cdot (b - a) = f(a)$
			
		\end{itemize}
		Poiché $g(a) = g(b)$, la funzione $g(x)$ soddisfa tutte le ipotesi del teorema di Rolle. Grazie a questo, sappiamo che $\exists x_0 \in (a, b)$ tale che $g'(x_0) = 0$.\\
		Calcolando la derivata prima di $g(x)$ otteniamo:
		\[
		g'(x) = f'(x) - \frac{f(b) - f(a)}{b - a}
		\]
		da cui segue che $g'(x_0) = 0 \iff f'(x_0) = \frac{f(b) - f(a)}{b - a}$.
		
	\end{dimostrazione}
	
	\begin{teorema}{Teorema di Cauchy}{}
		
		Siano $f, g \in C([a, b])$ derivabili in $(a, b)$, con $g'(x) \neq 0, \, \forall x \in (a, b)$.\\
		Allora $\exists x_0 \in (a, b)$ tale che:
		\[
		\frac{f(b) - f(a)}{g(b) - g(a)} = \frac{f'(x_0)}{g'(x_0)}
		\]
		
	\end{teorema}
	
	\begin{dimostrazione}{}{}
		
		Poniamo $g(a) \neq g(b)$ e definiamo $h \in C([a, b])$ come:
		\[
		h(x) = f(x) - \frac{f(b) - f(a)}{g(b) - g(a)} \cdot (g(x) - g(a))
		\]
		La funzione $h$ è derivabile in $(a, b)$ e soddisfa $h(a) = h(b) = f(a)$.\\
		Per il teorema di Rolle, $\exists x_0 \in (a, b)$ tale che:
		\[
		0 = h'(x_0) = f'(x_0) - \frac{f(b) - f(a)}{g(b) - g(a)} \cdot g'(x_0)
		\]
		Dividendo per $g'(x_0) \neq 0$, si ottiene infine:
		\[
		\frac{f'(x_0)}{g'(x_0)} = \frac{f(b) - f(a)}{g(b) - g(a)}
		\]
		
	\end{dimostrazione}
	
	\section{Monotonia e derivate}
	
	\begin{teorema}{Relazione tra monotonia e derivata}{}
		
		Sia $f: I \to \R$, con $I$ intervallo, una funzione derivabile in $I$.
		\begin{enumerate}
			
			\item[\textbf{A)}] Se $f$ è crescente, allora $f'(x) \geq 0, \, \forall x \in I$.
			
			\item[\textbf{B)}] Se $f'(x) \geq 0, \, \forall x \in I$, allora $f$ è crescente.
			
			\item[\textbf{C)}] Se $f'(x) > 0, \, \forall x \in I$, allora $f$ è strettamente crescente.
			
		\end{enumerate}
		
	\end{teorema}
	
	\begin{dimostrazione}{Dimostrazione del punto A}{}
		
		Sia $I = [a, b]$ e sia $x_0 \in (a, b)$.\\
		Se $f$ è crescente, allora si hanno due casi distinti:
		\begin{enumerate}
			
			\item Per $x < x_0 \implies f(x) \leq f(x_0)$
			
			\item Per $x > x_0 \implies f(x) \geq f(x_0)$
			
		\end{enumerate}
		
		Consideriamo il rapporto incrementale della funzione in $x_0$:
		\[
		\frac{f(x) - f(x_0)}{x - x_0}
		\]
		
		Analizziamo i due casi:
		\begin{enumerate}
			
			\item Per $x < x_0$, il numeratore è $\leq 0$ e il denominatore è $< 0$. Di conseguenza, i due termini hanno segno concorde, rendendo il rapporto incrementale $\geq 0$. Passando al limite per $x \to x_0^-$, la derivata sinistra soddisfa $f'_-(x_0) \geq 0$.
			
			\item Per $x > x_0$, il numeratore è $\geq 0$ e il denominatore è $> 0$. Anche in questo caso i due termini hanno segno concorde, quindi il rapporto incrementale è $\geq 0$. Passando al limite per $x \to x_0^+$, la derivata destra soddisfa $f'_+(x_0) \geq 0$.
		
		\end{enumerate}
		Essendo $f$ derivabile in $x_0$, i due limiti coincidono con $f'(x_0)$, dimostrando così che $f'(x_0) \geq 0$.
		
	\end{dimostrazione}
	% --- NEW PAGE ---
	
	\begin{dimostrazione}{Dimostrazione del punto B}{}
		
		Supponiamo che $f'(x) \geq 0$ per ogni $x \in I$. Possiamo applicare il teorema di Lagrange.
		
		Siano $a, b \in I$ con $a < b$, e $f$ derivabile in $[a, b]$. Allora esiste $c \in (a, b)$ tale che:
		\[
		f'(c) = \frac{f(b) - f(a)}{b - a} \implies f(b) - f(a) = \underbrace{f'(c)}_{\geq 0} \underbrace{(b - a)}_{> 0} \geq 0
		\]
		Da cui si ottiene:
		\[
		f(b) - f(a) \geq 0 \implies f(b) \geq f(a)
		\]
		Avendo dimostrato che per $b > a$ risulta $f(b) \geq f(a)$, possiamo concludere che la funzione $f$ è crescente.
		
	\end{dimostrazione}
	
	\begin{osservazione}{}{}
		
		Sia $f: I \to \R$ una funzione continua. 
		Se $f'(x) \geq 0$ per $x \leq x_0$ e $f'(x) \leq 0$ per $x \geq x_0$, la funzione $f$ cresce fino a $x_0$ per poi decrescere; il punto $x_0$ è dunque un \textbf{massimo locale}.
		
		Analogamente, se $f'(x) \leq 0$ a sinistra di $x_0$ e $f'(x) \geq 0$ a destra di $x_0$, allora $x_0$ rappresenta un \textbf{minimo locale}.
		
	\end{osservazione}
	
	\section{Funzioni lipschitziane}
	
	\begin{definizione}{Funzione lipschitziana}{}
		
		Sia $f: I \to \R$ una funzione definita su un intervallo $I \subseteq \R$.
		La funzione $f$ è detta \textbf{lipschitziana} se esiste una costante $L > 0$ tale che:
		\[
		|f(x) - f(y)| \leq L |x - y| \quad \forall x, y \in I
		\]
		
	\end{definizione}
	
	Un classico esempio di funzione lipschitziana è $f(x) = |x|$, con costante $L = 1$.
	In questo caso la condizione si riduce a verificare che:
	\[
	||x| - |y|| \leq |x - y|
	\]
	Dalla disuguaglianza triangolare $|a + b| \leq |a| + |b|$, riscriviamo $|x|$ come:
	\[
	|x| = |(x - y) + y| \leq |x - y| + |y| \implies |x| - |y| \leq |x - y|
	\]
	Poiché il valore assoluto è una funzione pari, scambiando i ruoli di $x$ e $y$ si ottiene anche $|y| - |x| \leq |y - x| = |x - y|$, da cui segue la tesi $||x| - |y|| \leq |x - y|$.
	
	\paragraph{Esempio di funzione non Lipschitziana}
	Consideriamo la funzione $f(x) = \sqrt{x}$, con $f: [0, a] \to \R$.
	Affinché $f$ sia Lipschitziana, dovrebbe esistere una costante $L \ge 0$ tale che $|\sqrt{x} - \sqrt{y}| \le L \cdot |x - y|$. In particolare, ponendo $y = 0$, la condizione diventa:
	\[
	|\sqrt{x}| \le L \cdot |x| \implies \frac{\sqrt{x}}{x} \le L \implies \frac{1}{\sqrt{x}} < L
	\]
	Tuttavia, per $x \to 0^+$, si ha che $\frac{1}{\sqrt{x}} \to +\infty$, dunque non potrà mai essere limitata da alcuna costante $L$. Abbiamo così dimostrato che nell'intervallo $[0, a]$ la funzione $f(x) = \sqrt{x}$ non è Lipschitziana.
	
	\begin{teorema}{Derivata limitata e Lipschitzianità}{}
		
		Sia $f: I \to \R$ derivabile in $I$. Se la derivata $f'(x)$ è limitata, cioè $\sup |f'(x)| = L < +\infty$, allora $f$ è Lipschitziana.
		
	\end{teorema}
	
	\begin{dimostrazione}{}{}
		
		Dati $x_1, x_2 \in I$ con $x_1 \neq x_2$, vogliamo dimostrare che $|f(x_1) - f(x_2)| \le L |x_1 - x_2|$.
		
		\smallbreak
		Siccome $f$ è derivabile nell'intervallo $[x_1, x_2]$, essendo questo una restrizione dell'intervallo $I$, possiamo applicare il teorema di Lagrange:
		\[
		\exists c \in (x_1, x_2) \quad \text{t.c.} \quad \frac{f(x_2) - f(x_1)}{x_2 - x_1} = f'(c)
		\]
		Da cui segue:
		\[
		f(x_2) - f(x_1) = f'(c)(x_2 - x_1) \implies |f(x_2) - f(x_1)| = |f'(c)| \cdot |x_2 - x_1|
		\]
		Sappiamo che $|f'(c)|$ è maggiorabile dalla costante $L$, per cui concludiamo:
		\[
		|f(x_1) - f(x_2)| \le L |x_1 - x_2|
		\]
		
	\end{dimostrazione}
	
	\begin{proposizione}{}{}
		
		Sia $f: I \to \R$ una funzione derivabile. Allora $f$ è costante se e solo se $f'(x) = 0, \, \forall x \in I$.
		
	\end{proposizione}
	
	\begin{dimostrazione}{}{}
		
		Supponiamo che $f'(x) = 0, \, \forall x \in I$. Siano $x_1, x_2 \in I$ con $x_1 < x_2$.
		
		\smallbreak
		Sempre per il teorema di Lagrange si ha:
		\[
		f(x_2) - f(x_1) = \underbrace{f'(c)}_{=0} (x_2 - x_1)
		\]
		Dunque $f(x_2) = f(x_1)$, il che è valido per ogni coppia di punti in $I$ se e solo se $f(x) = c$.
		
	\end{dimostrazione}
	
	\begin{teorema}{Teorema di de l'Hôpital}{}
		
		Siano $f, g: I_{x_0} \to \R$, con $I_{x_0} = \begin{cases} (a, b) \ni x_0 \\ (x_0, b) \\ (a, x_0) \end{cases}$.
		
		Se $\lim_{x \to x_0} f(x) = \lim_{x \to x_0} g(x) \in \{0, -\infty, +\infty\}$, $f$ e $g$ sono derivabili in $I_{x_0}$ (con la possibile eccezione di $x_0$) e supponendo che $\exists \lim_{x \to x_0} \frac{f'(x)}{g'(x)}$ con $g'(x) \neq 0$, allora:
		\[
		\lim_{x \to x_0} \frac{f(x)}{g(x)} = \lim_{x \to x_0} \frac{f'(x)}{g'(x)}
		\]
		
	\end{teorema}
	
	Come esempio applicativo, consideriamo il limite $\lim_{x \to 0} \frac{e^x - 1}{x}$, che si presenta nella forma indeterminata $\left[\frac{0}{0}\right]$. Sapendo che $(e^x)' = e^x$, $(1)' = 0$ e $(x)' = 1$, applicando il teorema di de l'Hôpital otteniamo:
	\[
	\lim_{x \to 0} \frac{e^x - 1}{x} = \lim_{x \to 0} \frac{e^x}{1} = e^0 = 1
	\]
	
	\begin{dimostrazione}{}{}
		
		Sia $\lim_{x \to x_0^+} f(x) = \lim_{x \to x_0^+} g(x) = 0$ e definiamo $f(x_0) = g(x_0) = 0$, estendendo le funzioni nel caso in cui $x_0$ sia un estremo dell'intervallo aperto.
		
		Consideriamo l'intervallo $[x_0, b]$ in cui $f$ e $g$ sono derivabili e con $g'(x) \neq 0$. Applichiamo il teorema di Cauchy:
		\[
			\frac{f(b) - f(x_0)}{g(b) - g(x_0)} = \frac{f'(c)}{g'(c)} \quad \text{con } c \in (x_0, b)
		\]
		Essendo $f(x_0) = g(x_0) = 0$, sostituendo si ha:
		\[
			\frac{f(b)}{g(b)} = \frac{f'(c)}{g'(c)}
		\]
		Passando al limite per $b \to x_0^+$, dato che $c \in (x_0, b)$, anche $c \to x_0^+$. Giungiamo dunque alla conclusione:
		\[
			\lim_{b \to x_0^+} \frac{f(b)}{g(b)} = \lim_{c \to x_0^+} \frac{f'(c)}{g'(c)} = \lim_{x \to x_0^+} \frac{f'(x)}{g'(x)}
		\]
		
		Ci accorgiamo che, per come sono definiti i valori e i rispettivi intervalli, si forma la catena $x_0 < c < b$ e possiamo dunque dire che $b \to x_0$ e $c \to x_0$.
		È dunque verificata la precedente uguaglianza che dimostra il teorema.
		
		Verifichiamo ora un caso particolare in cui $x_0 = +\infty$:
		\[
			l = \lim_{x \to +\infty} f(x) = \lim_{x \to +\infty} g(x) = 0
		\]
		
		Effettuando il cambio di variabile $z = \frac{1}{x}$, possiamo considerare il limite:
		\[
			\lim_{z \to 0^+} \frac{f\left(\frac{1}{z}\right)}{g\left(\frac{1}{z}\right)}
		\]
		A tale espressione possiamo applicare il teorema di De L'Hôpital, in quanto abbiamo già dimostrato la casistica per $x \to 0$:
		\[
			\lim_{z \to 0^+} \frac{\left(f\left(\frac{1}{z}\right)\right)'}{\left(g\left(\frac{1}{z}\right)\right)'} = \lim_{z \to 0^+} \frac{\left(\frac{1}{z}\right)' f'\left(\frac{1}{z}\right)}{\left(\frac{1}{z}\right)' g'\left(\frac{1}{z}\right)} = \lim_{z \to 0^+} \frac{f'\left(\frac{1}{z}\right)}{g'\left(\frac{1}{z}\right)}
		\]
		
		Sostituendo infine $x = \frac{1}{z}$, si giunge a:
		\[
			\lim_{x \to +\infty} \frac{f'(x)}{g'(x)}
		\]
		
	\end{dimostrazione}
	
	\begin{teorema}{Limite della derivata}{}
		
		Sia $f$ una funzione continua in $x_0$ e derivabile nell'insieme $(a, b) \setminus \{x_0\}$.\\
		Se esiste il limite $\lim_{x \to x_0} f'(x) = l$, allora $f$ è derivabile anche nel punto $x_0$ e vale $f'(x_0) = l$.
		
	\end{teorema}
	
	Verifichiamo la derivabilità usando il teorema appena enunciato:
	\[
		g'(x) = 
		\begin{cases}
			1 + 12x, & x > 0 \\
			\cos(5x) - 5x\sin(5x), & x < 0
		\end{cases}
	\]
	
	Calcoliamo il limite destro della derivata: $\lim_{x \to 0^+} g'(x) = \lim_{x \to 0^+} 1 + 12x = 1$.
	Analogamente, calcoliamo il limite sinistro: $\lim_{x \to 0^-} g'(x) = \lim_{x \to 0^-} \cos(5x) - 5x\sin(5x) = 1$.
	
	Siccome i due limiti coincidono, $g$ è derivabile anche in $x = 0$ e $g'(0) = 1$.
	
	Potrebbe anche esserci richiesto di stabilire se esiste $g''(0)$. E noi potremmo calcolarlo anche con il limite del rapporto incrementale, alla vecchia maniera.
	
	Calcoliamo il limite destro per il rapporto incrementale della derivata:
	\[
		\lim_{h \to 0^+} \frac{g'(x_0 + h) - g'(x_0)}{h}
	\]
	Siccome stiamo considerando $x_0 = 0$, abbiamo:
	\[
		\lim_{h \to 0^+} \frac{g'(h) - g'(0)}{h}
	\]
	Avendo definito $g'(0) = 1$, sostituiamo la definizione della funzione e otteniamo:
	\[
		\lim_{h \to 0^+} \frac{1 + 12h - 1}{h} = 12
	\]
	
	Similmente procediamo per il limite sinistro:
	\[
		\lim_{h \to 0^-} \frac{g'(x_0 + h) - g'(x_0)}{h} \sim \text{procedimento analogo} \sim \lim_{h \to 0^-} \frac{g'(h) - g'(0)}{h}
	\]
	Sostituendo la funzione per $x < 0$:
	\[
		\lim_{h \to 0^-} \frac{\cos(5h) - 5h\sin(5h) - 1}{h} = \lim_{h \to 0^-} \left( \frac{\cos(5h) - 1}{5h} \cdot 5 - \frac{5h\sin(5h)}{h} \right) = 0 \cdot 5 - 5 \cdot 0 = 0
	\]
	
	Siccome i due limiti differiscono, $g'$ non è derivabile in $x = 0$ e quindi la derivata seconda $g''(0)$ non esiste.
	
	\section{Retta tangente}
	
	Data una funzione $f(x)$ e consideratone il grafico $y = f(x)$, possiamo determinare l'equazione della retta tangente in un punto $(x_0, f(x_0))$ con la seguente formula:
	\[
		y = f(x_0) + f'(x_0)(x - x_0)
	\]
	ciò perché la derivata di una funzione $f$ ricordiamo essere il limite del coefficiente angolare della retta secante passante per $(x_0, f(x_0))$ e $(x_0+h, f(x_0+h))$.\\
	Infatti, per valori di $h$ molto piccoli (ossia tendenti a $0$), i due punti diventano praticamente coincidenti e la retta secante assume il comportamento della retta tangente.
	
	\chapter{Sviluppo di Taylor}
	
	Lo \textbf{sviluppo di Taylor} di una funzione, nell'intorno di un punto $x_0$, è la rappresentazione della funzione come somma di un polinomio e di un infinitesimo di ordine superiore al grado del polinomio stesso.
	
	Sia $f: I_{x_0} \to \R$:
	\begin{itemize}
		
		\item Possiamo approssimare $f$ con un polinomio di grado $0$ ($P_0$) scrivendo in questo modo:
		\[
			P_0(f, x_0) = f(x_0)
		\]
		dove $P_0(f, x_0)$ rappresenta il polinomio di Taylor di grado $0$ costruito in $x_0$, il quale risulta essere sempre pari alla costante $f(x_0)$.
		
	\end{itemize}
	
	Possiamo generalizzare, nel caso del polinomio $P_0$, nel seguente modo:
	\[
		P_0(f, x_0)(x) = f(x_0)
	\]
	Cioè il polinomio $P_0$, definito come approssimazione di $f$ per tutti i punti vicini ad $x_0$, per ogni $x$ dell'intorno di definizione della funzione è sempre uguale a $f(x_0)$.
	
	Possiamo dunque riscrivere la funzione $f$ confrontandola con il polinomio definito precedentemente:
	\[
		f(x) = P_0(f, x_0)(x) + o(1) \quad \longrightarrow \quad P_0(f, x_0)(x) - f(x) = o(1)
	\]
	dove $o(1)$ simboleggia l'errore di approssimazione.
	
	\section{Notazione di $o$-piccolo}
	
	Diciamo che $f$ è un $o$ ($o$-piccolo) di $g$ per $x \to x_0$ se:
	\[
	\lim_{x \to x_0} \frac{f(x)}{g(x)} = 0
	\]
	e si può indicare formalmente come $f \in o(g(x))$ oppure $f = o(g(x))$.
	
	Dunque, nel nostro caso:
	\begin{itemize}
		
		\item Sia $g(x) = o(1)$, allora $\lim_{x \to x_0} g(x) = 0$.
		
	\end{itemize}
	
	Sostituendo la definizione del polinomio $P_0(f, x_0)(x) = f(x_0)$ nel confronto tra $f$ e $P_0$, abbiamo che:
	\[
		f(x_0) - f(x) = o(1)
	\]
	ovvero la loro differenza tende a $0$. Ciò equivale a scrivere:
	\[
		\lim_{x \to x_0} (f(x_0) - f(x)) = 0 \quad \text{oppure} \quad \lim_{x \to x_0} \frac{P_0(f, x_0)(x) - f(x)}{1} = 0
	\]
	
	Andiamo adesso a vedere le casistiche per l'approssimazione a polinomi di grado superiore.
	
	\subsection{Approssimazione di grado 1}
	
	Se consideriamo $P_1(f, x_0)(x) = f(x_0) + f'(x_0)(x - x_0)$, allora:
	\[
		P_1(f, x_0)(x) - f(x) = o(x - x_0)
	\]
	
	Possiamo anche qui riscrivere la relazione sotto forma di limite:
	\[
		\lim_{x \to x_0} \frac{P_1(f, x_0)(x) - f(x)}{x - x_0} = 0 \implies \lim_{x \to x_0} \frac{f(x_0) + f'(x_0)(x - x_0) - f(x)}{x - x_0} = 0
	\]
	Separando i termini al numeratore:
	\[
		\lim_{x \to x_0} \left( \frac{f(x_0) - f(x)}{x - x_0} + \frac{f'(x_0)(x - x_0)}{x - x_0} \right) = 0
	\]
	Semplificando e riordinando il rapporto incrementale:
	\[
		\lim_{x \to x_0} \left( - \underbrace{\frac{f(x) - f(x_0)}{x - x_0}}_{\text{rapporto incrementale}} + f'(x_0) \right) = 0
	\]
	Poiché il limite del rapporto incrementale per $x \to x_0$ è pari a $f'(x_0)$, otteniamo:
	\[
		\lim_{x \to x_0} \left( -f'(x_0) + f'(x_0) \right) = 0
	\]
	
	\subsection{Approssimazione di grado 2}
	
	Se consideriamo il polinomio di secondo grado della forma:
	\[
		P_2(f, x_0)(x) = f(x_0) + f'(x_0)(x - x_0) + \sigma (x - x_0)^2
	\]
	vogliamo mostrare che il coefficiente incognito è $\sigma = \frac{f''(x_0)}{2}$, da cui segue:
	\[
		P_2(f, x_0)(x) = f(x_0) + f'(x_0)(x - x_0) + \frac{f''(x_0)}{2} (x - x_0)^2
	\]
	Valutiamo quindi la condizione sul resto $P_2(f, x_0)(x) - f(x) = o((x - x_0)^2)$, la quale equivale a:
	\[
		\lim_{x \to x_0} \frac{P_2(f, x_0)(x) - f(x)}{(x - x_0)^2} = 0
	\]
	Esplicitando il polinomio:
	\[
		\lim_{x \to x_0} \frac{f(x_0) + f'(x_0)(x - x_0) + \sigma (x - x_0)^2 - f(x)}{(x - x_0)^2} = 0
	\]
	Essendo $f(x_0)$ una costante, possiamo applicare il teorema di De L'Hôpital derivando i termini rispetto a $x$:
	\[
		\lim_{x \to x_0} \frac{f'(x_0) + 2\sigma(x - x_0) - f'(x)}{2(x - x_0)} = 0
	\]
	Semplificando il fattore $2(x - x_0)$ presente al numeratore e al denominatore:
	\[
		\lim_{x \to x_0} \frac{f'(x_0) + \sigma - f'(x)}{2(x - x_0)} = 0 \implies \sigma - \lim_{x \to x_0} \frac{f'(x) - f'(x_0)}{2(x - x_0)} = 0
	\]
	Ricordando la definizione di derivata seconda come limite del rapporto incrementale della derivata prima:
	\[
		\sigma - \frac{1}{2} \cdot \lim_{x \to x_0} \frac{f'(x) - f'(x_0)}{x - x_0} = 0 \implies \sigma - \frac{f''(x_0)}{2} = 0 \implies \sigma = \frac{f''(x_0)}{2}
	\]
	
	\subsection{Grado $k$ generico}
	
	Per un ordine $k$ generico, il polinomio di Taylor assume la forma:
	\[
		P_k(f, x_0)(x) = \sum_{i=0}^{k} \frac{f^{(i)}(x_0)}{i!} (x - x_0)^i
	\]
	In particolare, nel caso in cui il centro dello sviluppo sia $x_0 = 0$, il polinomio prende il nome di \textbf{polinomio di Maclaurin}.
	
	\begin{teorema}{Formula del resto di Peano}{}
		
		Se $f$ è derivabile $k$ volte nel punto $x_0$, allora sussiste la seguente relazione per il resto:
		\[
			P_k(f, x_0)(x) - f(x) = o((x - x_0)^k)
		\]
		
	\end{teorema}
	
	\newpage
	\section{Polinomi di Maclaurin delle funzioni elementari}
	
	Ricordiamo la struttura generale del polinomio di Maclaurin di ordine $n$:
	\[
		P_n(f, 0)(x) = \sum_{k=0}^{n} \frac{f^{(k)}(0)}{k!} x^k
	\]
	
	Vediamo la determinazione dello sviluppo per alcune funzioni elementari fondamentali:
	
	\begin{enumerate}
		
		\item \textbf{Funzione esponenziale} $f(x) = e^x$:\\
		Poiché $f'(x) = e^x$ e in generale $f^{(k)}(x) = e^x$, calcolando nell'origine otteniamo $f^{(k)}(0) = 1$ per ogni $k \in \N$.
		
		Il polinomio di Maclaurin è dunque $P_n(e^x, 0)(x) = \sum_{k=0}^{n} \frac{x^k}{k!}$, da cui lo sviluppo:
		\[
		e^x = \sum_{k=0}^{n} \frac{x^k}{k!} + o(x^n)
		\]
		
		\item \textbf{Funzione seno} $f(x) = \sin x$:\\
		Trattandosi di una funzione dispari, tutte le derivate di indice pari calcolate in $0$ si annullano. Calcoliamo le prime derivate:
		\begin{itemize}
			
			\item $f(0) = 0$
			
			\item $f'(x) = \cos x \implies f'(0) = 1$
			
			\item $f''(x) = -\sin x \implies f''(0) = 0$
			
			\item $f'''(x) = -\cos x \implies f'''(0) = -1$
			
			\item $f^{IV}(x) = \sin x \implies f^{IV}(0) = 0$
			
		\end{itemize}
		I valori delle derivate si ripetono periodicamente. Sopravvivono soltanto i termini di grado dispari, ottenendo lo sviluppo:
		\[
			\sin x = \sum_{k=0}^{n} (-1)^k \frac{x^{2k+1}}{(2k+1)!} + o(x^{2n+2})
		\]
		
		\item \textbf{Funzione coseno} $f(x) = \cos x$:\\
		Poiché la funzione è pari, le derivate di indice dispari calcolate in $0$ sono nulle. Le derivate soddisfano la relazione $f^{(4k+h)}(x) = f^{(h)}(x)$, ripetendosi ogni $4$ ordini di derivazione:
		\begin{itemize}
			
			\item $f(0) = 1$
			
			\item $f'(x) = -\sin x \implies f'(0) = 0$
			
			\item $f''(x) = -\cos x \implies f''(0) = -1$
			
			\item $f'''(x) = \sin x \implies f'''(0) = 0$
			
			\item $f^{IV}(x) = \cos x \implies f^{IV}(0) = 1$
			
		\end{itemize}
		Considerando unicamente i contributi delle derivate di indice pari, si ricava lo sviluppo:
		\[
		\cos x = \sum_{k=0}^{n} (-1)^k \frac{x^{2k}}{(2k)!} + o(x^{2n+1})
		\]
	\end{enumerate}
	
	\begin{osservazione}{}{}
		
		Un'osservazione opportuna riguarda il resto della differenza tra il polinomio di Taylor e la funzione approssimata.
		
		Si potrebbe pensare che, all'aumentare degli esponenti, aumenti anche l'errore. Tuttavia questo non è vero, come si può notare da un semplice esempio numerico:
		\begin{itemize}
			
			\item Sia $x - x_0 = 0.1$;
			
			\item $(0.1)^2 > (0.1)^3 > \dots > (0.1)^k$.
		
		\end{itemize}
		All'aumentare dell'esponente, il valore dell'errore diminuisce e il resto tende a zero più velocemente.
		
	\end{osservazione}
	
	\begin{teorema}{Teorema del resto di Lagrange}{}
		
		Poste le seguenti condizioni:
		\begin{itemize}
			
			\item $f$ è derivabile $k+1$ volte;
			
			\item $f, f', f'', \dots, f^{(k)}$ sono funzioni continue;
			
			\item $f^{(k+1)}$ esiste in un intorno di $x_0$ (con al più l'esclusione di $x_0$ stesso);
		
		\end{itemize}
		allora l'errore commesso nell'approssimazione è dato da:
		\[
		f(x) - P_k(f, x_0)(x) = \frac{1}{(k+1)!} \cdot f^{(k+1)}(\overline{x}) \cdot (x - x_0)^{k+1}
		\]
		con $\overline{x} \in (x, x_0)$.
		
	\end{teorema}
	
	\begin{proposizione}{}{}
		
		Il polinomio di MacLaurin di una funzione pari contiene solo termini di grado pari. Analogamente, il polinomio di MacLaurin di una funzione dispari contiene solo termini di grado dispari.
		
	\end{proposizione}
	
	\begin{dimostrazione}{}{}
		
		Se $f$ è una funzione pari, vale che $f(x) = f(-x)$. Derivando entrambi i membri si ottiene $f'(x) = -f'(-x)$, quindi la sua derivata prima $f'$ è una funzione dispari.
		
		In generale, tutte le derivate di ordine dispari $f^{(2k+1)}(x)$ risulteranno essere funzioni dispari.
		Sia $g$ una generica funzione dispari per cui vale la relazione $g(x) = -g(-x)$. Valutando tale funzione nel punto $x_0 = 0$, si ha:
		\[
		g(0) = -g(0) \implies 2g(0) = 0 \implies g(0) = 0
		\]
		Siccome per $x_0 = 0$ tutte le derivate di ordine dispari si annullano ($f^{(2k+1)}(0) = 0$), nel polinomio di MacLaurin scompaiono tutti i termini corrispondenti ai gradi dispari.
		
	\end{dimostrazione}
	
	Consideriamo la funzione $f(x) = (1 + x)^\alpha$:
	\begin{align*}
		f'(x) &= \alpha(1 + x)^{\alpha - 1} \\
		f''(x) &= \alpha(\alpha - 1)(1 + x)^{\alpha - 2} \\
		f^{(k)}(x) &= \alpha(\alpha - 1)\dots(\alpha - k + 1)(1 + x)^{\alpha - k}
	\end{align*}
	
	Nel punto $x_0 = 0$ le derivate successive portano allo sviluppo:
	\[
		(1 + x)^\alpha = \sum_{k=0}^n \underbrace{\frac{\alpha(\alpha - 1)\dots(\alpha - k + 1)}{k!}}_{\binom{\alpha}{k}} x^k + o(x^n)
	\]
	
	Nel caso particolare in cui $\alpha = -1$, si ottiene:
	\[
		(1 + x)^{-1} = \frac{1}{1 + x} = \sum_{k=0}^n (-1)^k x^k + o(x^n)
	\]
	
	\begin{teorema}{}{}
		
		Sia $f: (a, b) \to \R$, $x_0 \in (a, b)$ ed $f$ derivabile $n$ volte in $x_0$.
		Se $f(x) = P_n(x) + o((x - x_0)^n)$, con $P_n(x)$ generico polinomio di grado $\le n$, allora $P_n(x) = Q_n(x)$, dove $Q_n(x) = P_n(f, x_0)(x)$ è il polinomio di Taylor di grado $n$ di $f$ centrato in $x_0$.
		
		Inoltre:
		\[
			P_n(x) - Q_n(x) = \sum_{k=0}^n c_k (x - x_0)^k = c_0 + c_1(x - x_0)^1 + \dots + c_n(x - x_0)^n
		\]
		dove ogni termine $c_k$ rappresenta la differenza tra i coefficienti dei termini di $P_n(x)$ e $Q_n(x)$.
		
	\end{teorema}
	
	\begin{dimostrazione}{}{}
		
		Infatti, se esprimiamo $P_n(x)$ come generico polinomio di grado al più $n$:
		\[
			P_n(x) = \sum_{k=0}^n a_k(x - x_0)^k
		\]
		e ricordando la definizione del polinomio di Taylor $Q_n(x)$:
		\[
			Q_n(x) = \sum_{k=0}^n \frac{f^{(k)}(x_0)}{k!}(x - x_0)^k
		\]
		possiamo scrivere la differenza tra i due polinomi come:
		\[
			P_n(x) - Q_n(x) = \sum_{k=0}^n \underbrace{\left(a_k - \frac{f^{(k)}(x_0)}{k!}\right)}_{c_k} (x - x_0)^k
		\]
		
		Dunque il termine $c_k$ misura quanto i coefficienti di $P_n$ differiscono da quelli corretti.
		
		Adesso, sia $m$ l'indice più piccolo diverso da $0$ (cioè ad esempio $c_1 = 0, c_2 = 0, \dots, c_{m-1} = 0, c_m \neq 0$):
		\[
			\frac{P_n(x) - Q_n(x)}{(x-x_0)^m} = c_m + \sum_{k=m+1}^{n} c_k (x-x_0)^{k-m}
		\]
		ovvero, esplicitando la sommatoria:
		\[
			\frac{P_n(x) - Q_n(x)}{(x-x_0)^m} = c_m + c_{m+1}(x-x_0) + \dots + c_n(x-x_0)^{n-m}
		\]
		
		Ora passiamo al limite per $x \to x_0$:
		\begin{itemize}
			
			\item I termini $(x-x_0), (x-x_0)^2, \dots$ vanno a zero;
			
			\item Rimane solo:
			\[
				\lim_{x \to x_0} \frac{P_n(x) - Q_n(x)}{(x-x_0)^m} = c_m \neq 0
			\]
			
		\end{itemize}
		
		Tuttavia l'ipotesi iniziale ci diceva che:
		\begin{itemize}
			
			\item $P_n(x) - Q_n(x) = o((x-x_0)^n)$
			
			\item Dividendo per $(x-x_0)^m$:
			\[
				\frac{P_n(x) - Q_n(x)}{(x-x_0)^m} = o((x-x_0)^{n-m})
			\]
			
			\item Passando al limite:
			\[
				\lim_{x \to x_0} o((x-x_0)^{n-m}) = 0
			\]
			
		\end{itemize}
		
		Ma i due limiti non coincidono!
		
		Dunque la nostra ipotesi sul fatto che $\exists m \le n$ tale che $c_m \neq 0$ si è rivelata falsa, e quindi:
		\[
			c_0 = c_1 = c_2 = \dots = 0 \implies P_n(x) = Q_n(x)
		\]
		
		Questo dimostra l'\textbf{unicità del polinomio di Taylor/Maclaurin} per una funzione $f$.
		
	\end{dimostrazione}
	
	\section{Calcolo degli sviluppi mediante sostituzione e moltiplicazione}
	
	Vediamo come determinare lo sviluppo di MacLaurin di funzioni composte o prodotto tramite opportune sostituzioni e proprietà dell'$o$-piccolo.
	
	Consideriamo ad esempio la funzione $f(x) = e^{2x}$. Sostituendo $t = 2x$ nello sviluppo noto dell'esponenziale:
	\[
		e^t = \sum_{k=0}^n \frac{t^k}{k!} + o(t^n)
	\]
	otteniamo la seguente espressione:
	\[
		e^{2x} = \sum_{k=0}^n \frac{(2x)^k}{k!} + o((2x)^n)
	\]
	
	Notevole è il comportamento del resto secondo la notazione $o$-piccolo. Se $g(x) \in o((2x)^n)$, per definizione di limite si ha:
	\[
		\lim_{x \to 0} \frac{g(x)}{(2x)^n} = 0 \implies \lim_{x \to 0} \frac{g(x)}{x^n} = 0 \implies g(x) \in o(x^n)
	\]
	Possiamo dunque omettere il fattore costante $2^n$ all'interno dell'O-piccolo, in quanto non influisce sul comportamento asintotico per $x \to 0$. Lo sviluppo finale risulta pertanto:
	\[
		e^{2x} = \sum_{k=0}^n \frac{2^k x^k}{k!} + o(x^n)
	\]
	
	\begin{proposizione}{}{}
		
		In generale, se $f(x) \in o(x^n)$ per $x \to 0$, allora moltiplicando per un fattore $x^k$ si ottiene:
		\[
		x^k \cdot f(x) \in o(x^{n+k})
		\]
		
	\end{proposizione}
	
	Consideriamo il calcolo del polinomio di Taylor $P_6(f, 0)(x)$ per la funzione $f(x) = \sin(2x^2)$.
	
	Essendo il seno una funzione dispari, prendiamo in considerazione le derivate di indice dispari. In generale, lo sviluppo del sesto ordine per $\sin(t)$ attorno a $x_0$ si scrive come:
	\[
		P_6(\sin(t), x_0) = f'(x_0)(x - x_0) + \frac{f'''(x_0)}{6}(x - x_0)^3 + \frac{f^{(5)}(x_0)}{120}(x - x_0)^5 + o((x - x_0)^6)
	\]
	
	Nel nostro caso, ponendo $x_0 = 0$, otteniamo:
	\[
		P_6(\sin(t), 0) = f'(0) \cdot t + \frac{f'''(0)}{6} \cdot t^3 + \frac{f^{(5)}(0)}{120} \cdot t^5 + o(t^6)
	\]
	
	Sostituendo i valori delle derivate della funzione seno nel punto zero:
	\[
		\begin{aligned}
			P_6(\sin(t), 0) &= \cos(0) \cdot t - \frac{\cos(0)}{6} \cdot t^3 + \frac{\cos(0)}{120} \cdot t^5 + o(t^6) \\
			&= t - \frac{t^3}{6} + \frac{t^5}{120} + o(t^6)
		\end{aligned}
	\]
	
	Andiamo ora a sostituire la variabile $t = 2x^2$:
	\[
		2x^2 - \frac{(2x^2)^3}{6} + \frac{(2x^2)^5}{120} + o((2x^2)^6) = 2x^2 - \frac{8x^6}{6} + \frac{32x^{10}}{120} + o(64x^{12})
	\]
	
	L'espressione ottenuta risulta essere un polinomio di grado $12$; tuttavia, la richiesta iniziale era di arrestare lo sviluppo al grado $6$.
	
	Ci accorgiamo che per $x \to x_0 = 0$ vale $x^{10} \in o(x^6)$, in quanto:
	\[
		\lim_{x \to 0} \frac{x^{10}}{x^6} = \lim_{x \to 0} x^4 = 0
	\]
	
	Possiamo dunque "troncare" la successione dei termini al grado $6$ e aggiungere come resto $o(x^6)$.
	
	In conclusione, otteniamo:
	\[
		P_6(\sin(2x^2), 0) = 2x^2 - \frac{4}{3}x^6 + o(x^6)
	\]
	
	\section{Proprietà degli $o$-piccoli}
	
	\begin{enumerate}
		
		\item $o(x^n) \pm o(x^m) = o(x^p)$ dove $p = \min(n, m)$
		
		\begin{dimostrazione}{}{}
			
			Supponiamo $p = n < m$, $f(x) \in o(x^n)$ e $g(x) \in o(x^m)$. Si ha:
			\[
			\frac{f(x) \pm g(x)}{x^n} = \frac{f(x)}{x^n} \pm \frac{g(x)}{x^n} = \frac{f(x)}{x^n} \pm \frac{g(x)}{x^m} \cdot x^{m-n} \xrightarrow[x \to 0]{} 0
			\]
			in quanto $\frac{f(x)}{x^n} \to 0$ per definizione di $f(x) \in o(x^n)$ e $\frac{g(x)}{x^m} \to 0$ per definizione di $g(x) \in o(x^m)$.
			
		\end{dimostrazione}
		
		\item $x^m \in o(x^n) \iff m > n$
		
		\begin{dimostrazione}{}{}
			
			Se $m > n$, possiamo riscrivere il limite come:
			\[
			\lim_{x \to 0} \frac{x^m}{x^n} = \lim_{x \to 0} x^{m-n} = 0
			\]
			
		\end{dimostrazione}
		
		\item $x^m \cdot o(x^n) = o(x^{m+n}) \quad \text{e} \quad \frac{o(x^n)}{x^m} = o(x^{n-m})$
		
		\begin{dimostrazione}{}{}
			
			Sia $f(x) \in o(x^n)$, allora:
			\[
			\frac{x^m \cdot f(x)}{x^{m+n}} = \frac{f(x)}{x^n} \to 0
			\]
			il che è verificato per come abbiamo defin
			ito la funzione $f$.
		\end{dimostrazione}
		
		\item $o(x^n) \cdot o(x^m) = o(x^{m+n})$
		
		\begin{dimostrazione}{}{}
			
			Siano $f(x) \in o(x^n)$ e $g(x) \in o(x^m)$. Si ha:
			\[
				\frac{f(x) \cdot g(x)}{x^{n+m}} = \frac{f(x)}{x^n} \cdot \frac{g(x)}{x^m} \xrightarrow[x \to 0]{} 0
			\]
			
		\end{dimostrazione}
	\end{enumerate}
	
	\section{Altri sviluppi di Taylor}
	
	Se consideriamo la funzione $f(x) = \ln(1+x)$, calcolandone le prime derivate otteniamo:
	\[
		f'(x) = \frac{1}{1+x} = (1+x)^{-1}
	\]
	\[
		f''(x) = (-1)(1+x)^{-2}, \qquad f'''(x) = (-2)(-1)(1+x)^{-3}
	\]
	Dunque, la derivata $n$-esima generica ha la forma:
	\[
		f^{(n)}(x) = (n-1)! (-1)^{n-1} (1+x)^{-n}
	\]
	Allora, il polinomio di Maclaurin corrispondente di un generico grado $n$ sarà:
	\[
		T_n(f, 0) = \sum_{k=0}^n \frac{f^{(k)}(0)}{k!} \cdot x^k = \sum_{k=1}^n \frac{(k-1)!(-1)^{k-1}}{k!} \cdot x^k = \sum_{k=1}^n \frac{(-1)^{k-1}}{k} \cdot x^k
	\]
	
	\chapter{Numeri complessi}
	Si indicano con la lettera $\C$.
	
	Per introdurre i numeri complessi, consideriamo ad esempio la seguente equazione:
	\[
		x^2 + 1 = 0
	\]
	la quale sappiamo non ammettere soluzioni $x \in \R$.
	
	Definiamo dunque l'unità immaginaria $i$ tale che $i^2 = -1$, la quale costituisce uno zero del polinomio considerato. Definiamo poi l'insieme $\C$ come:
	\[
		\C = \{a + bi \mid a, b \in \R\}
	\]
	ovvero l'insieme di tutti i numeri espressi in tale forma.
	
	Definiamo di seguito le principali operazioni nell'insieme dei numeri complessi:
	\begin{description}
		
		\item [Somma:] $(a + ib) + (c + id) = (a + c) + i(b + d)$
		
		\item [Prodotto:] $(a + ib)(c + id) = ac + ibc + iad + i^2 bd = (ac - bd) + i(bc + ad)$
		
		\item [Opposto:] $-(a + ib) = -a - ib$
		
		\item [Inverso:]
		\[
			(a + ib)^{-1} = \frac{1}{a + ib} = \frac{1}{a + ib} \cdot \frac{a - ib}{a - ib} = \frac{a - ib}{a^2 - i^2 b^2} = \frac{a - ib}{a^2 + b^2} = \frac{a}{a^2 + b^2} - i \frac{b}{a^2 + b^2}
		\]
		
	\end{description}
	
	\begin{osservazione}{}{}
		
		L'operazione che consiste nel moltiplicare un numero complesso per se stesso cambiando il segno della parte immaginaria è detta \textbf{coniugio}. Il coniugato di un numero complesso $z = a + ib$ si indica comunemente con $\bar{z} = a - ib$.
		
	\end{osservazione}
	
	Ad esempio, dato $z = 2 + 3i$, il suo inverso è:
	\[
		z^{-1} = \frac{1}{z} \cdot \frac{\overline{z}}{\overline{z}} = \frac{1}{2+3i} \cdot \frac{2-3i}{2-3i} = \frac{2}{4+9} - i \frac{3}{4+9} = \frac{2}{13} - i \frac{3}{13}
	\]
	
	\begin{definizione}{Modulo e forma algebrica di un numero complesso}{}
		
		Sia $z = a + ib$ un numero complesso.
		\begin{itemize}
			
			\item $a = \operatorname{Re}(z)$ è la \textbf{parte reale} di $z$.
			
			\item $b = \operatorname{Im}(z)$ è la \textbf{parte immaginaria} di $z$.
			
			\item Il \textbf{modulo di $z$} è definito come la quantità reale non negativa:
			\[
				|z| := \sqrt{a^2 + b^2}
			\]
			
		\end{itemize}
		
	\end{definizione}
	
	\begin{osservazione}{}{}
		
		Nel caso in cui $\operatorname{Im}(z) = 0$, il modulo coincide col valore assoluto della parte reale, ossia $\sqrt{a^2} = |a|$.
		
	\end{osservazione}
	
	\section{Rappresentazione polare e trigonometrica}
	
	Ogni numero complesso $z = a + ib$ può essere rappresentato nel piano cartesiano (piano di Gauss) tramite il punto di coordinate $(a, b)$:
	
	\begin{center}
		
		\begin{tikzpicture}[scale=1.2]
			
			\draw[->, thick] (-0.3, 0) -- (2.5, 0) node[right] {\small $\operatorname{Re}$};
			\draw[->, thick] (0, -0.3) -- (0, 2) node[above] {\small $\operatorname{Im}$};
			
			\draw[dashed] (1.8, 1.5) -- (1.8, 0) node[below] {\small $a$};
			\draw[dashed] (1.8, 1.5) -- (0, 1.5) node[left] {\small $b$};
			
			\draw[blue, thick] (0, 0) -- (1.8, 1.5) node[midway, above left] {\small $|z|$};
			\fill[black] (1.8, 1.5) circle (1.5pt) node[right] {\small $z = (a, b)$};
			
			\draw[orange, thick, ->] (0.5, 0) arc (0:39.8:0.5);
			\node[orange] at (0.7, 0.2) {\small $\theta$};
		
		\end{tikzpicture}
		
	\end{center}
	
	In questa rappresentazione:
	\begin{itemize}
		
		\item $|z|$ rappresenta la \textbf{distanza dall'origine}, comunemente indicata con il raggio $r$.
	
		\item $\theta$ è l'\textbf{argomento} di $z$, dove $\operatorname{Arg}(z) \in (-\pi, \pi]$.
	
	\end{itemize}
	
	Per le relazioni trigonometriche fondamentali vale che:
	\[
		a = r \cdot \cos\theta, \qquad b = r \cdot \sin\theta
	\]
	
	Sostituendo le componenti nell'espressione $z = a + ib$, otteniamo la \textbf{rappresentazione polare o trigonometrica}:
	\[
		z = r(\cos\theta + i\sin\theta)
	\]
	
	Definiamo un'uguaglianza fondamentale nel dominio dei numeri complessi:
	\[
	e^{i\theta} = \cos\theta + i \cdot \sin\theta
	\]
	
	Notiamo dunque che per esprimere un generico numero complesso $z$ occorre moltiplicare per il raggio $r$:
	\[
		z = r \cdot e^{i\theta}
	\]
	
	In questo modo risulta più semplice effettuare calcoli tra numeri complessi distinti. Dati $z_1 = r_1 \cdot e^{i\theta_1}$ e $z_2 = r_2 \cdot e^{i\theta_2}$, si ha:
	\[
		z_1 \cdot z_2 = e^{i\theta_1} \cdot e^{i\theta_2} = e^{i(\theta_1 + \theta_2)}
	\]
	
	\chapter{Complessità asintotiche}
	
	Siano $f, g: \N \to \R$.
	
	\begin{definizione}{$O$-grande}{}
		
		Diciamo che $f \in O(g)$ se e solo se $\exists n_0 \in \N, \, c > 0$ tale che:
		\[
		0 \leq |f(n)| \leq c \cdot |g(n)| \quad \forall n \geq n_0
		\]
		
	\end{definizione}
	
	Sia ad esempio $f(n) = 3n^2 + 10n$. Vale che $f(n) \in O(n^2)$ poiché:
	\[
		3n^2 + 10n \leq 3n^2 + 10n^2 = 13n^2 \quad \forall n \geq 0
	\]
	In questo caso abbiamo scelto $c = 13$ e $n_0 = 0$. La costante $c$ è stata ricavata dalla maggiorazione scelta (e non è unica), mentre il valore $n_0$ è facilmente individuabile tramite semplici calcoli.
	
	Alcune delle classi $O$-grande più comuni sono:
	\begin{itemize}
	
		\item $O(\log n)$: complessità logaritmica
	
		\item $O(n)$: complessità lineare
	
		\item $O(n^k)$: crescita polinomiale (nel caso di un polinomio si prende il termine di grado maggiore)
	
		\item $O(\alpha^n)$: complessità esponenziale
	
	\end{itemize}
	
	\begin{definizione}{$\Omega$-grande}{}
		
		Siano $f, g: \mathbb{N} \to \mathbb{R}$. Diciamo che $f \in \Omega(g)$ se e solo se esistono $n_0 \in \mathbb{N}$ e $c > 0$ tali che:
		\[
		0 \le c|g(n)| \le |f(n)| \quad \forall n \ge n_0
		\]
		
	\end{definizione}
	
	Consideriamo ad esempio la funzione $f(n) = 3n^2 + 10n$. Si ha che $f(n) \in \Omega(n^2)$ poiché $3n^2 \le 3n^2 + 10n$ con $c = 3$ e $n_0 = 0$. Al contrario, si verifica che $n^2 \notin \Omega(n^3)$.
	
	\begin{definizione}{$\Theta$-grande}{}
		
		Siano $f, g: \mathbb{N} \to \mathbb{R}$. Diciamo che $f \in \Theta(g)$ se esistono $c_1, c_2 > 0$ e $n_0 \in \mathbb{N}$ tali che:
		\[
			c_1 \cdot |g(n)| \le |f(n)| \le c_2 \cdot |g(n)| \quad \forall n \ge n_0
		\]
		
		Equivalentemente, possiamo esprimere tale relazione come:
		\[
			f \in \Theta(g) \iff f \in O(g) \quad \text{e} \quad f \in \Omega(g)
		\]
		
	\end{definizione}
	
	\begin{definizione}{}{}
		
		Siano $f, g: \mathbb{N} \to \mathbb{R}^+ = (0, +\infty)$. Diciamo che $f \in \omega(g)$ se:
		\[
			\lim_{n \to +\infty} \frac{f(n)}{g(n)} = +\infty
		\]
		
	\end{definizione}
	
	\begin{teorema}{Relazione tra o-piccolo e O-grande}{}
		Se $f \in \omega(g)$ allora $f \in \Omega(g)$.
	\end{teorema}
	
	Non vale tuttavia il contrario. Consideriamo infatti $g(n): \N \to \R^+$ e $f(n) := 2 \cdot g(n)$.
	Se $f \in \Omega(g)$ allora $f(n) \ge c \cdot g(n) \implies 2 \cdot g(n) \ge c \cdot g(n)$ e ciò è chiaramente vero, con $n_0 = 0$ e $c = 2$.
	Tuttavia, $f \in \omega(g)$ vuol dire che $\lim_{n \to +\infty} \frac{f(n)}{g(n)} = +\infty$ ma, per come abbiamo definito $f$, $\lim_{n \to +\infty} \frac{2g(n)}{g(n)} = 2$ che è chiaramente diverso da $+\infty$.
	
	\section{Operazioni con gli O-grandi}
	Siano $f_1, f_2, g_1, g_2: \N \to \R^+$.
	Se $f_1 \in O(g_1)$ e $f_2 \in O(g_2)$ allora:
	\begin{enumerate}
		
		\item $f_1 + f_2 \in O(g_1 + g_2)$
		
		\item $f_1 \cdot f_2 \in O(g_1 \cdot g_2)$
		
		\item $(f_1)^k \in O({g_1}^k)$
	
	\end{enumerate}
	
	\begin{dimostrazione}{Dimostrazione 1}{}
		
		Dalle ipotesi sappiamo che:
		\begin{itemize}
		
			\item $f_1 \in O(g_1) \implies |f_1(n)| \le c_1 \cdot |g_1(n)| \quad \forall n \ge n_1$
		
			\item $f_2 \in O(g_2) \implies |f_2(n)| \le c_2 \cdot |g_2(n)| \quad \forall n \ge n_2$

		\end{itemize}
		Ponendo $n_0 = \max(n_1, n_2)$ e considerando $n \ge n_0$ possiamo maggiorare il tutto:
		\[
			f_1 + f_2 \le c_1 \cdot g_1(n) + c_2 \cdot g_2(n)
		\]
		Definendo $c = \max(c_1, c_2)$ otteniamo:
		\[
			f_1 + f_2 \le c \cdot g_1(n) + c \cdot g_2(n) \implies f_1 + f_2 \le c \cdot (g_1 + g_2)
		\]
		da cui la tesi.
		
	\end{dimostrazione}
	
	\begin{dimostrazione}{Dimostrazione 2}{}
		
		Analogamente a prima, sappiamo che:
		\[
			f_1 \cdot f_2 \le c_1 \cdot g_1(n) \cdot c_2 \cdot g_2(n) \implies f_1 \cdot f_2 \le (c_1 \cdot c_2)(g_1 \cdot g_2)
		\]
		Impostiamo $c = c_1 \cdot c_2$ e otteniamo il risultato sperato:
		\[
			f_1 \cdot f_2 \le c \cdot (g_1 \cdot g_2)
		\]
		
	\end{dimostrazione}
	
	\chapter{Serie numeriche}
	
	\begin{definizione}{Serie numerica}{}
		
		Data una successione di numeri reali $\{a_n\}$, si chiama \textbf{serie} dei termini $a_n$ la somma degli infiniti termini della successione.
		
	\end{definizione}
	
	Solitamente la notazione utilizzata è la seguente:
	\[
		\sum_{n=0}^{\infty} a_n
	\]
	che si legge "\textit{serie (o somma) per $n$ che va da $0$ a $\infty$ di $a_n$}".
	
	Tuttavia, sommare un numero infinito di termini è un po' complicato. Ci costruiamo dunque una successione intermedia, chiamata $S_n$ (successione delle somme parziali), che definiamo nel seguente modo:
	\[
		S_0 = a_0, \quad S_1 = a_0 + a_1, \quad \dots, \quad S_n = a_0 + a_1 + \dots + a_n = \sum_{k=0}^{n} a_k
	\]
	In questo modo, come si può notare, non abbiamo più una somma di infiniti termini.
	
	La nuova successione $\{s_n\}$ è detta \textbf{successione delle somme parziali della serie}.
	
	E per estendere il tutto ad un numero infinito:
	\[
		\sum_{n=0}^{\infty} a_n = \lim_{n \to +\infty} s_n = \lim_{n \to +\infty} \sum_{k=0}^{n} a_k
	\]
	e questo è il modo naturale e formale di definire una serie.
	
	Per questo limite abbiamo tre possibilità:
	\begin{itemize}
		
		\item $\lim_{n \to +\infty} \sum_{k=0}^{n} a_k = l \in \R$, si dice che la serie \textbf{converge}.
		
		\item $\lim_{n \to +\infty} \sum_{k=0}^{n} a_k = \pm\infty$, si dice che la serie \textbf{diverge}.
		
		\item $\nexists \lim_{n \to +\infty} \sum_{k=0}^{n} a_k$, si dice che la serie è \textbf{indeterminata}.
		
	\end{itemize}
	
	Solitamente, delle serie si studia il \textbf{comportamento}, cioè cosa sappiamo del limite e cosa possiamo dedurne in base al risultato.
	
	\begin{osservazione}{}{}
		
		Abbiamo visto le serie già nell'approssimazione di una funzione con il polinomio di Taylor. Potremmo usare le serie anche per scrivere un numero reale con sviluppo decimale infinito, ad esempio:
		\[
			2{,}3456\dots = 2 + \frac{3}{10} + \frac{4}{100} + \frac{5}{1000} + \frac{6}{10000} + \dots
		\]
		
	\end{osservazione}
	
	Inoltre, non è necessario iniziare la serie da $n=0$. Ad esempio:
	\[
		\sum_{n=1}^{\infty} \frac{1}{n} \qquad \text{oppure} \qquad \sum_{n=2}^{\infty} \frac{1}{n(n-1)}
	\]
	In quest'ultimo caso, iniziare da $n=0$ non sarebbe stato possibile poiché si sarebbe annullato il denominatore.
	
	Tuttavia, se una serie converge, iniziare da un indice più avanti \textbf{ne cambia il risultato}:
	\[
		\sum_{n=1}^{\infty} \frac{1}{n^3} = S, \qquad \sum_{n=7}^{\infty} \frac{1}{n^3} = \bar{S} = S - \sum_{n=1}^{6} \frac{1}{n^3}
	\]
	In questo modo andiamo a togliere i termini da $1$ a $6$ dal valore della somma totale $S$ trovato in precedenza.
	
	Dunque: \textbf{stesso comportamento, somma diversa}.
	
	\section{Introduzione alle serie numeriche}
	
	Consideriamo alcuni esempi fondamentali per comprendere il comportamento delle somme infinite (serie) e delle loro somme parziali.
	
	\begin{description}
		
		\item[Studio della successione $a_n = (-1)^n$:] 
		La successione assume valori alternati a seconda della parità dell'indice:
		\[
			a_n = (-1)^n = 
			\begin{cases} 
				1 & \text{se } n \text{ è pari} \\ 
				-1 & \text{se } n \text{ è dispari} 
			\end{cases}
		\]
		Proviamo a calcolare la sommatoria analizzando la successione delle somme parziali $S_n$:
		\begin{align*}
			S_0 &= (-1)^0 = 1 \\
			S_1 &= (-1)^0 + (-1)^1 = 1 - 1 = 0 \\
			S_2 &= (-1)^0 + (-1)^1 + (-1)^2 = 1 \\
			S_3 &= (-1)^0 + (-1)^1 + (-1)^2 + (-1)^3 = 0 
		\end{align*}
		
		In generale, notiamo che la successione delle somme parziali è data da:
		\[
			S_n = \sum_{k=0}^{n} (-1)^k = 
			\begin{cases} 
				0 & \text{se } n \text{ è dispari} \\ 
				1 & \text{se } n \text{ è pari} 
			\end{cases}
		\]
		Poiché la successione delle somme parziali oscilla continuamente tra $0$ e $1$, si ha che $\not\exists \lim_{n \to +\infty} S_n$. Di conseguenza, la serie è \textbf{indeterminata}.
		
		\item[Studio della successione costante $a_n = 1, \, \forall n \in \mathbb{N}$:] 
		Intuitivamente capiamo che la serie deve divergere a $+\infty$. Per dimostrarlo formalmente, calcoliamo la somma parziale $S_K$:
		\[
			S_K = \sum_{n=0}^{K} 1 = \underbrace{1 + 1 + 1 + \dots + 1}_{K+1 \text{ termini}} = K + 1
		\]
		Poiché il limite delle somme parziali è $\lim_{K \to +\infty} S_K = \lim_{K \to +\infty} (K + 1) = +\infty$, la serie \textbf{diverge} a $+\infty$.
		
	\end{description}
	
	\begin{osservazione}{}{}
		
		Più in generale, data una successione a termini costanti del tipo $a_n = c \in \mathbb{R}$, il comportamento della serie associata $\sum_{n=0}^{\infty} a_n$ è determinato unicamente dal valore di $c$:
		\[
			\sum_{n=0}^{\infty} c = 
			\begin{cases} 
				+\infty & \text{se } c > 0 \\ 
				-\infty & \text{se } c < 0 \\ 
				0 & \text{se } c = 0 
			\end{cases}
		\]
		
	\end{osservazione}
	
	\begin{description}
		
		\item[Studio della successione $a_n = \frac{n}{n+1}$:] 
		Contrariamente a quanto si potrebbe istintivamente pensare, la serie degli elementi di questa successione \textbf{diverge} a $+\infty$. 
		
		Analizzandone i primi termini, vediamo che questi sono:
		\[
			a_0 = 0, \quad a_1 = \frac{1}{2}, \quad a_2 = \frac{2}{3}, \quad a_3 = \frac{3}{4}, \ \dots
		\]
		A partire da $n = 1$, tutti i termini della successione sono maggiori o uguali a $\frac{1}{2}$. Di conseguenza, nella sommatoria stiamo accumulando infiniti termini stabilmente superiori a $\frac{1}{2}$, il che impedisce alla somma di convergere.
		
	\end{description}
	
	\begin{proposizione}{Condizione necessaria per la convergenza}{}
		
		Se la serie $\sum_{n=0}^{\infty} a_n$ converge, allora:
		\[
			\lim_{n\to\infty} a_n = 0
		\]
		Non possiamo affermare il contrario, ma possiamo dire che se $\lim_{n\to\infty} a_n \neq 0$ allora la serie non converge.
		
	\end{proposizione}
	
	\begin{osservazione}{}{}
		
		Ad esempio, si consideri la serie armonica $\sum_{n=1}^{\infty} \frac{1}{n} = +\infty$. La successione associata $a_n = \frac{1}{n}$ converge a $0$, ma la serie diverge a $+\infty$ poiché i suoi termini non tendono a $0$ abbastanza velocemente.
		
	\end{osservazione}
	
	\begin{dimostrazione}{}{}
		
		Dimostriamo la proposizione:
		\begin{itemize}
			
			\item Se la serie converge, allora esiste finito il limite delle somme parziali: $\exists \lim_{k\to\infty} S_k = S$.
			
			\item Riscriviamo le somme parziali come:
			\[
				S_k = \sum_{n=0}^{k} a_n = a_0 + \dots + a_k
			\]
			
			\item Di conseguenza, abbiamo anche:
			\[
				S_{k+1} = \sum_{n=0}^{k+1} a_n = a_0 + \dots + a_k + a_{k+1}
			\]
			
			\item Siccome $k$ è un indice generico, possiamo affermare che $\lim_{k\to\infty} S_{k+1} = S$.
			
			\item Infine, considerando che $S_{k+1} - S_k = a_{k+1}$ (ovvero l'unico elemento per cui le due somme parziali differiscono), si ha:
			\[
				\lim_{k\to\infty} a_{k+1} = \lim_{k\to\infty} (S_{k+1} - S_k) = S - S = 0
			\]
			
		\end{itemize}
		
	\end{dimostrazione}
	
	\section{Serie geometriche}
	
	Si dice \textbf{serie geometrica} una serie definita nel seguente modo:
	\[
		\sum_{n=0}^{\infty} x^n, \quad \text{con } x \in \mathbb{R}
	\]
	Per studiarne il comportamento dobbiamo verificare il valore della base $x$. A prescindere dal suo valore, possiamo riscrivere il termine generale come:
	\[
		x^n = \frac{1}{\frac{1}{x^n}}
	\]
	Dunque si presentano i seguenti casi:
	\begin{itemize}
		
		\item Se $x > 1$, allora $\lim_{n\to\infty} x^n = +\infty$.
		
		\item Se $x = 1$, allora $x^n = 1$, da cui $\lim_{n\to\infty} 1^n = 1$.
		
		\item Se $0 < x < 1$, riscrivendo $x^n = \frac{1}{(1/x)^n}$ con $\frac{1}{x} > 1$, si ottiene $\lim_{n\to\infty} x^n = 0$.
		
		\item Se $x \le -1$, il termine $x^n$ assume segni alterni, di conseguenza il limite non esiste.
	
	\end{itemize}
	
	Dunque l'unico modo per far convergere la serie è scegliere una base $-1 < x < 1$.
	
	Tuttavia, le serie geometriche possono essere riscritte come un valore calcolabile:
	\[
		S_k = \sum_{n=0}^{k} x^n = 1 + x + x^2 + \dots + x^k = \frac{1 - x^{k+1}}{1 - x} \quad \text{con } x \neq 1
	\]
	
	E dunque:
	\[
		\lim_{k \to +\infty} \frac{1 - x^{k+1}}{1 - x} = 
		\begin{cases}
			+\infty, & \text{se } x \ge 1 \\
			\dfrac{1}{1 - x}, & \text{se } -1 < x < 1 \\
			\nexists, & \text{se } x \le -1
		\end{cases}
	\]
	
	Che per la serie si traduce in:
	\[
		\sum_{n=0}^{\infty} x^n = 
		\begin{cases}
			\text{diverge a } +\infty, & \text{se } x \ge 1 \\
			\text{converge a } \dfrac{1}{1 - x}, & \text{se } -1 < x < 1 \\
			\text{indeterminata}, & \text{se } x \le -1
		\end{cases}
	\]
	
	Ad esempio:
	\[
		\sum_{n=0}^{\infty} \left(\frac{1}{2}\right)^n = \frac{1}{1 - \frac{1}{2}} = 2 \quad \text{siccome } -1 < \frac{1}{2} < 1
	\]
	
	In generale, per serie geometriche che non partono da $0$ ma da un generico $p$:
	\[
		\sum_{n=p}^{\infty} x^n = \frac{x^p}{1 - x}
	\]
	
	\section{Serie telescopiche}
	
	\begin{definizione}{Serie telescopica}{}
		
		Si dicono \textbf{serie telescopiche} quelle serie di successioni $a_n$ in cui il termine generale può essere riscritto nella forma:
		\[
			a_n = b_{n+1} - b_n
		\]
		
	\end{definizione}
	
	Sapendo ciò, se consideriamo la somma parziale $S_k = a_0 + a_1 + \dots + a_k$, possiamo riscriverla come:
	\[
		S_k = (b_1 - b_0) + (b_2 - b_1) + \dots + (b_k - b_{k-1}) + (b_{k+1} - b_k)
	\]
	
	Semplificando i termini opposti, si ottiene:
	\[
		S_k = b_{k+1} - b_0
	\]
	
	Di conseguenza, il comportamento della serie è determinato dal limite della successione $b_k$:
	\[
		\sum_{n=0}^{\infty} a_n := \lim_{k \to +\infty} S_k = \lim_{k \to +\infty} (b_{k+1} - b_0) = 
		\begin{cases}
			\pm\infty, & \text{se } b_k \to \pm\infty \\
			b - b_0, & \text{se } b_k \to b \in \mathbb{R} \\
			\text{indeterminata}, & \text{se } \nexists \lim_{k \to +\infty} b_k
		\end{cases}
	\]
	
	Partendo invece da un indice generico $p$, si sostituisce $b_p$ ad ogni istanza di $b_0$.
	
	Consideriamo la serie:
	\[
		\sum_{n=1}^{\infty} \ln\left(1 + \frac{1}{n}\right)
	\]
	Riscrivendo in modo conveniente il logaritmo si ha:
	\[
		\ln\left(1 + \frac{1}{n}\right) = \ln\left(\frac{n+1}{n}\right) = \ln(n+1) - \ln(n)
	\]
	Abbiamo così trovato $b_n = \ln(n)$ tale che $a_n = b_{n+1} - b_n$. Dunque la somma parziale $K$-esima risulta:
	\[
		\sum_{n=1}^{K} \ln\left(1 + \frac{1}{n}\right) = \ln(K+1) - \ln(1) = \ln(K+1)
	\]
	Pertanto, il comportamento della serie è dato da:
	\[
		\sum_{n=1}^{\infty} \ln\left(1 + \frac{1}{n}\right) = \lim_{K \to +\infty} \ln(K+1) = +\infty
	\]
	
	\begin{proposizione}{Proprietà delle serie}{}
		
		Siano $\sum a_n$ e $\sum b_n$ due serie. Valgono le seguenti proprietà:
		\begin{itemize}
			
			\item $\sum_{n=0}^{\infty} \lambda \cdot a_n = \lambda \cdot \sum_{n=0}^{\infty} a_n \quad \forall \lambda \in \mathbb{R}$
			
			\item $\sum_{n=0}^{\infty} (a_n + b_n) = \sum_{n=0}^{\infty} a_n + \sum_{n=0}^{\infty} b_n$, tranne nel caso in cui una serie diverga a $+\infty$ e l'altra a $-\infty$ (o viceversa).
		
		\end{itemize}
		
	\end{proposizione}
	
	\section{Serie a termini positivi}
	
	Sia $a_n \geq 0$. Allora la serie dei termini $a_n$ non può essere indeterminata. Questo è molto intuitivo, in quanto se tutti i termini sono positivi non possono esserci oscillazioni o cambi di segno.
	
	Dunque, le serie a termini positivi possono essere unicamente:
	\begin{itemize}
		
		\item \textbf{Convergenti}
		
		\item \textbf{Divergenti} a $+\infty$
		
	\end{itemize}
	
	Consideriamo ad esempio la serie $\sum_{n=0}^{\infty} \frac{n}{n+1}$:
	\begin{itemize}
		
		\item I termini sono positivi, quindi la serie non è indeterminata.
		
		\item Poiché il termine generale tende a $1$ (e quindi non tende a $0$), la serie non soddisfa la condizione necessaria di convergenza.
		
		\item Di conseguenza, la serie \textbf{diverge}.
	
	\end{itemize}
	
	\begin{teorema}{Teorema del confronto per le serie}{}
		
		Siano $a_n, b_n \geq 0$ tali che $a_n \leq b_n$ per ogni $n \geq n_0$.
		\begin{itemize}
		
			\item Se $\sum_{n=0}^{\infty} b_n < +\infty \implies \sum_{n=0}^{\infty} a_n < +\infty$
		
			\item Se $\sum_{n=0}^{\infty} a_n = +\infty \implies \sum_{n=0}^{\infty} b_n = +\infty$
		
		\end{itemize}
	\end{teorema}
	
	\begin{dimostrazione}{}{}
		
		Dimostriamo la prima proposizione, in quanto la seconda non è altro che la contronominale della prima.
		
		Per ipotesi, la serie $\sum_{n=0}^{\infty} b_n$ converge, dunque la successione delle sue somme parziali $S_k$ ammette limite finito:
		\[
			\lim_{k \to +\infty} S_k = S
		\]
		con $S_k \leq S$ (poiché la successione delle somme parziali di una serie a termini positivi è monotonamente crescente e quindi limitata superiormente dal suo estremo superiore, che coincide con il limite).
		
		Sapendo che $a_n \leq b_n$, per le rispettive somme parziali $T_k = \sum_{n=0}^{k} a_n$ vale la relazione:
		\[
			T_k \leq S_k
		\]
		e di conseguenza:
		\[
			\lim_{k \to +\infty} T_k = T \leq S_k \leq S
		\]
		Essendo la successione $T_k$ limitata superiormente da $S$, essa deve convergere a un limite finito $T \leq S$.
		
	\end{dimostrazione}
	
	Il teorema precedentemente annunciato risulta particolarmente utile quando dobbiamo studiare il comportamento di una serie, confrontandola con altre serie di cui conosciamo il comportamento.
	
	Per esempio, usiamo la serie telescopica $\sum_{n=1}^{\infty} \ln\left(1 + \frac{1}{n}\right) = +\infty$ per verificare il comportamento della serie armonica $\sum_{n=1}^{\infty} \frac{1}{n}$.\\
	Sappiamo che il logaritmo è minore di ogni polinomio, da cui vale la relazione $\ln(1+x) \leq x$.
	Siccome $a_n = \ln(1 + 1/n) \leq b_n = 1/n$ e la serie di $a_n$ diverge, per il teorema del confronto anche la serie di $b_n$ divergerà.
	
	\section{Serie armonica generalizzata}
	
	Chiamiamo \textbf{serie armonica generalizzata} la serie del tipo:
	\[
		\sum_{n=1}^{\infty} \frac{1}{n^\alpha}
	\]
	
	Sappiamo già che per $\alpha = 1$ la serie diverge, mentre per $\alpha = 2$ la serie converge. Verifichiamo come si comporta la serie per altri valori di $\alpha$:
	\begin{itemize}
		
		\item Per $\alpha > 2$, si ha $\frac{1}{n^\alpha} < \frac{1}{n^2} \iff n^2 < n^\alpha$. Per il teorema del confronto, siccome $\sum \frac{1}{n^2}$ converge, anche $\sum \frac{1}{n^\alpha}$ converge.
		
		\item Per $\alpha < 1$, si ha $\frac{1}{n^\alpha} > \frac{1}{n} \iff n^\alpha < n$. Anche in questo caso, per confronto, siccome $\sum \frac{1}{n}$ diverge, anche $\sum \frac{1}{n^\alpha}$ diverge.
		
		\item Per $1 < \alpha < 2$, si tratta di un caso più complesso, dimostrabile tramite il confronto integrale, ma anche in questo caso la serie converge.
		
	\end{itemize}
	
	In sintesi, possiamo riassumere il comportamento della serie come segue:
	\[
		\sum_{n=1}^{\infty} \frac{1}{n^\alpha} = 
		\begin{cases}
			+\infty & \text{se } \alpha \leq 1 \\
			< +\infty & \text{se } \alpha > 1
		\end{cases}
	\]
	
	\section{Studio del carattere di una serie}
	
	\begin{teorema}{Criterio del confronto asintotico}{}
		
		Siano $a_n, b_n \geq 0$ con $b_n > 0$.\\
		Se esiste il limite:
		\[
		\lim_{n \to +\infty} \frac{a_n}{b_n} = \ell \in (0, +\infty)
		\]
		allora la serie $\sum a_n$ converge se e solo se la serie $\sum b_n$ converge (e analogamente, la serie $\sum a_n$ diverge se e solo se la serie $\sum b_n$ diverge).
		
	\end{teorema}
	
	Consideriamo l'applicazione del criterio del confronto asintotico alla seguente serie:
	\[
		\sum_{n=1}^{\infty} \sin\left(\frac{1}{n}\right)
	\]
	Posto $a_n = \sin(1/n)$, scegliamo come successione di confronto $b_n = 1/n$. Sfruttando il limite notevole del seno, si ottiene:
	\[
		\lim_{n \to +\infty} \frac{\sin(1/n)}{1/n} = 1
	\]
	Poiché il limite è un valore finito e diverso da zero, le due serie hanno lo stesso comportamento. Siccome la serie armonica $\sum \frac{1}{n}$ diverge, anche la serie di partenza $\sum \sin(1/n)$ risulta divergente.
	
	\begin{teorema}{Variante del confronto asintotico}{}
		
		Siano $(a_n)$ e $(b_n)$ due successioni a termini positivi, con $b_n > 0$ per ogni $n \in \N$.
		Se esiste il limite:
		\[
		\lim_{n \to +\infty} \frac{a_n}{b_n} = 0
		\]
		allora, se la serie $\sum b_n$ converge, anche la serie $\sum a_n$ converge.
		
	\end{teorema}
	
	Ad esempio, consideriamo la serie:
	\[
		\sum_{n=2}^{\infty} \frac{1}{n^2 \ln(n)}
	\]
	Scegliendo come successione di confronto $b_n = \frac{1}{n^2}$, andiamo a verificare il limite del rapporto tra i termini generali:
	\[
		\lim_{n \to +\infty} \frac{\frac{1}{n^2 \ln(n)}}{\frac{1}{n^2}} = \lim_{n \to +\infty} \frac{1}{\ln(n)} = 0
	\]
	Poiché il limite è nullo e la serie associata a $b_n = \frac{1}{n^2}$ è una serie armonica generalizzata convergente, per la variante del confronto asintotico anche la serie di partenza converge.
	
	Per l'applicazione dei criteri di confronto basati sulle radici, è opportuno richiamare alcuni limiti notevoli fondamentali:
	
	\begin{itemize}
		
		\item Dato $l > 0$, si ha:
		\[
		\lim_{n \to +\infty} \sqrt[n]{l} = 1
		\]
		
		\item Per la successione identica, si ha:
		\[
		\lim_{n \to +\infty} \sqrt[n]{n} = 1
		\]
		
		\item Per il rapporto di due qualsiasi polinomi $P(n)$ e $Q(n)$, indipendentemente dal loro grado, si ha:
		\[
		\lim_{n \to +\infty} \sqrt[n]{\frac{P(n)}{Q(n)}} = 1
		\]
		
	\end{itemize}
	
	\begin{teorema}{Criterio della radice}{}
		
		Sia $a_n \geq 0$. Supponiamo che esista il limite:
		\[
		\lim_{n \to +\infty} \sqrt[n]{a_n} = l \in \R \cup \{+\infty\}
		\]
		Allora:
		\begin{itemize}
			
			\item Se $l < 1$, la serie $\sum a_n$ \textbf{converge};
			
			\item Se $l > 1$, la serie $\sum a_n$ \textbf{diverge};
			
			\item Se $l = 1$, il criterio non fornisce informazioni.
		
		\end{itemize}
		
	\end{teorema}
	
	Ad esempio, per la serie $\sum_{n=1}^{\infty} \frac{2^n+n^2+1}{3^n+n+4}$, calcolando il limite della radice si ottiene:
	\[
		\lim_{n \to +\infty} \sqrt[n]{a_n} = \frac{2}{3} < 1
	\]
	e di conseguenza la serie converge.
	
	\begin{teorema}{Criterio del rapporto}{}
		
		Sia $a_n > 0$. Al fine di studiare il comportamento della serie, consideriamo il limite:
		\[
		\lim_{n \to +\infty} \frac{a_{n+1}}{a_n} = l
		\]
		Allora:
		\begin{itemize}
			
			\item Se $l < 1$, la serie \textbf{converge};
			
			\item Se $l > 1$, la serie \textbf{diverge};
			
			\item Se $l = 1$, il criterio non fornisce informazioni.
		
		\end{itemize}
		
	\end{teorema}
	
	Questo criterio risulta particolarmente utile quando nella successione compare un fattoriale. Consideriamo ad esempio la serie $\sum_{n=0}^{\infty} \frac{2^n}{n!}$, in cui $a_n = \frac{2^n}{n!}$ e $a_{n+1} = \frac{2^{n+1}}{(n+1)!}$:
	\[
		\lim_{n \to +\infty} \frac{2^{n+1}}{(n+1)!} \cdot \frac{n!}{2^n} = \lim_{n \to +\infty} \frac{2 \cdot 2^n}{n!(n+1)} \cdot \frac{n!}{2^n} = \lim_{n \to +\infty} \frac{2}{n+1} = 0 < 1
	\]
	Essendo il limite minore di $1$, la serie converge.
	
	\section{Serie non a termini positivi}
	
	Per questa tipologia di serie non possiamo limitare a priori il comportamento a convergente o divergente, ma dobbiamo considerare anche la possibilità che la serie sia indeterminata. Di conseguenza, per studiarne il comportamento dovremo utilizzare tecniche differenti, tra cui:
	
	\begin{definizione}{Convergenza assoluta}{}
		
		Una serie $\sum_{n=0}^{\infty} a_n$ si dice \textbf{assolutamente convergente} se converge la serie dei suoi valori assoluti, ovvero se converge la serie:
		\[
			\sum_{n=0}^{\infty} |a_n|
		\]
		
	\end{definizione}
	
	Vale inoltre un importante risultato che lega la convergenza assoluta a quella ordinaria: se una serie converge assolutamente, allora la serie \textbf{converge} anche semplicemente.
	
	\subsection{Convergenza assoluta e Criterio di Leibniz}
	
	Consideriamo il seguente esempio per comprendere l'utilità della convergenza assoluta:
	\[
		\sum_{n=0}^{\infty} \frac{\cos(n^2+n+3)}{n^3+1} \implies \sum_{n=0}^{\infty} \left| \frac{\cos(n^2+n+3)}{n^3+1} \right|
	\]
	Ricordando che $|\cos x| \le 1$, possiamo applicare il criterio del confronto:
	\[
		\left| \frac{\cos(n^2+n+3)}{n^3+1} \right| \le \frac{1}{n^3+1} < \frac{1}{n^3}
	\]
	Siamo riusciti a maggiorare la serie con una serie convergente, dunque anche la serie dei valori assoluti $\sum |a_n|$ converge. Per il teorema della convergenza assoluta, se la serie dei valori assoluti converge, allora anche la serie di partenza $\sum a_n$ converge.
	
	Ci accorgiamo di come questo sia un criterio molto potente, tuttavia presenta dei limiti. Se consideriamo ad esempio la serie:
	\[
		\sum_{n=1}^{\infty} \frac{\sin(n)}{n} \implies \left| \frac{\sin(n)}{n} \right| \le \frac{1}{n}
	\]
	In questo caso non possiamo concludere nulla poiché la serie armonica $\sum \frac{1}{n}$ diverge. Anche se avessimo dimostrato che la serie dei valori assoluti diverge, non avremmo potuto trarre conclusioni sulla serie di partenza, in quanto la convergenza assoluta non fornisce informazioni in caso di divergenza della serie dei moduli.
	
	\begin{teorema}{Criterio di Leibniz}{}
		
		Si consideri una serie a segno alterno della forma:
		\[
			\sum_{n=0}^{\infty} (-1)^n a_n \quad \text{con } a_n \ge 0
		\]
		Se la successione $a_n$ rispetta i seguenti requisiti:
		\begin{enumerate}
			
			\item $a_n \ge 0$
			
			\item $a_n$ è monotona decrescente
			
			\item $\lim_{n \to +\infty} a_n = 0$
		
		\end{enumerate}
		Allora la serie \textbf{converge}.
		
	\end{teorema}
	
	\begin{osservazione}{}{}
		
		Consideriamo la serie armonica a segno alterno:
		\[
			\sum_{n=1}^{\infty} \frac{(-1)^n}{n}
		\]
		In questo caso $a_n = \frac{1}{n}$. Poiché la successione è a termini positivi, decrescente e infinitesima per $n \to +\infty$, per il criterio di Leibniz la serie \textbf{converge}.
		
	\end{osservazione}
	
	\chapter{Serie di potenze}
	
	Facciamo un esempio con una categoria di serie già vista:
	\[
		\sum_{n=0}^{\infty} x^n = \frac{1}{1-x} \quad \forall x \in (-1,1)
	\]
	Come vediamo, questa non è un'approssimazione, ma una vera e propria identità. Questo vale però solo se proseguiamo fino all'infinito: come abbiamo visto per i polinomi di Taylor, se ci fermiamo ad un certo punto otteniamo uno ``scarto'', un resto.
	
	Definendo ora più formalmente quanto visto sopra...
	
	\begin{definizione}{Serie di potenze}{}
		
		Data una successione $a_n$ e $x_0 \in \R$ fissato detto "centro", chiamiamo \textbf{serie di potenze} la serie costruita nel seguente modo:
		\[
			\sum_{n=0}^{\infty} a_n (x - x_0)^n
		\]
		dove $x$ è un'incognita, dunque il tutto viene espresso in funzione di $x$.\\
		Il punto delle serie di potenze è proprio di poterle riscrivere come funzioni definite in determinati intervalli.
		
	\end{definizione}
	
	Per una serie di potenze ci si pone le seguenti domande:
	\begin{itemize}
		
		\item Per quali $x$ la serie converge?
		
		\item Per quali $x$ si studia la funzione somma:
		\[
			S(x) = \sum_{n=0}^{\infty} a_n (x - x_0)^n
		\]
		
	\end{itemize}
	
	Ad esempio, consideriamo la serie $\sum_{n=0}^{\infty} \frac{x^n}{n!}$, con centro $x_0 = 0$ e coefficienti $a_n = \frac{1}{n!}$.
	
	\begin{osservazione}{}{}
		
		Ogni serie converge nel centro (banale).
		
		Poniamo infatti $x = x_0$:
		\[
		\sum_{n=0}^{\infty} a_n (x - x_0)^n = \sum_{n=0}^{\infty} a_n (x_0 - x_0)^n = \sum_{n=0}^{\infty} a_n (0)^n = a_0
		\]
		Poiché $0^0 = 1$, l'unico valore che si salva è $a_0 \cdot 1$.
		
	\end{osservazione}
	
	\begin{osservazione}{}{}
		
		Presa una serie di potenze, possiamo effettuare il cambio di variabile $y = x - x_0$ e esprimere tutto in funzione di $y$:
		\[
			\sum_{n=0}^{\infty} a_n (x - x_0)^n = \sum_{n=0}^{\infty} a_n y^n
		\]
		In questo modo il centro sarà $y_0 = 0$ e lo studio risulterà più semplice.
		
	\end{osservazione}
	
	Supponiamo $x_0 = 0$ e studiamo la serie $\sum_{n=0}^{\infty} a_n x^n$.
	
	\begin{proposizione}{}{}
		
		Se esiste il limite $\lim_{n \to \infty} \sqrt[n]{|a_n|} = l \in (0, +\infty)$, allora la serie converge $\forall x \in \left(-\frac{1}{l}, \frac{1}{l}\right)$ e non converge per il resto dei valori di $x$.
		
	\end{proposizione}
	
	\begin{dimostrazione}{}{}
		
		Proviamo a verificare la convergenza assoluta della serie:
		\[
			\sum_{n=0}^{\infty} |a_n \cdot x^n|
		\]
		Applichiamo il criterio della radice:
		\[
			\lim_{n \to \infty} \sqrt[n]{|a_n \cdot x^n|} = |x| \cdot \lim_{n \to \infty} \sqrt[n]{|a_n|} = |x| \cdot l
		\]
		Per il criterio della radice, la serie converge assolutamente se il limite è minore di $1$:
		\[
			|x| \cdot l < 1 \implies |x| < \frac{1}{l} \implies -\frac{1}{l} < x < \frac{1}{l}
		\]
		Otteniamo così lo stesso intervallo enunciato in precedenza.
		
	\end{dimostrazione}
	
	Supponiamo di avere una generica serie di potenze con centro in $0$ ($x_0 = 0$):
	\[
		\sum_{n=0}^{\infty} a_n \cdot x^n
	\]
	Sappiamo già che una serie converge sicuramente nel suo centro.
	
	Per identificare l'intervallo di convergenza di una serie dobbiamo definire il cosiddetto \textbf{raggio di convergenza $R$}, ossia quel valore $R \in [0, +\infty]$ tale che:
	\begin{itemize}
		
		\item la serie converge $\forall x$ tale che $|x - x_0| < R$;
		
		\item la serie non converge $\forall x$ tale che $|x - x_0| > R$.
	
	\end{itemize}
	
	Più formalmente, se consideriamo $E$ come l'insieme dei punti $x$ in cui la serie converge, il raggio di convergenza $R$ è definito come:
	\[
		R = \sup \{ |x - x_0| \mid x \in E \}
	\]
	dove $|x - x_0|$ rappresenta la distanza di un punto $x$ dal centro $x_0$. Il valore $R$ rappresenta dunque la distanza massima dal centro per cui la serie risulta convergente.
	
	Schematizzando, in base al valore di $R$ si presentano tre diverse casistiche:
	\[
		R = \begin{cases} 
			0, & \text{la serie converge solo nel centro } x_0 \\ 
			+\infty, & \text{la serie converge } \forall x \in \R \\ 
			(0, +\infty), & \text{la serie converge nell'intervallo } (x_0 - R, x_0 + R) 
		\end{cases}
	\]
	
	Per il calcolo del raggio di convergenza si utilizzano principalmente due metodi:
	
	\begin{description}
		
		\item [Criterio della radice] Se esiste il limite $\lim_{n \to +\infty} \sqrt[n]{|a_n|} = \ell \in [0, +\infty]$, allora il raggio $R$ è dato da:
		\[
			R = \begin{cases} 
				\frac{1}{\ell}, & \text{se } \ell \in (0, +\infty) \\ 
				0, & \text{se } \ell = +\infty \\ 
				+\infty, & \text{se } \ell = 0 
			\end{cases}
		\]
		
		\item [Criterio del rapporto] Se $a_n \neq 0$ e se esiste il limite $\lim_{n \to +\infty} \left| \frac{a_{n+1}}{a_n} \right| = \ell$, il raggio $R$ si calcola analogamente a quanto visto sopra. In particolare, vale la relazione diretta:
		\[
			R = \lim_{n \to +\infty} \left| \frac{a_n}{a_{n+1}} \right|
		\]
		
	\end{description}
	
	\begin{description}
		
		\item [Esempio] Consideriamo la serie $\sum_{n=0}^{\infty} \frac{x^n}{n!}$, con $a_n = \frac{1}{n!}$ e centro $x_0 = 0$. Calcoliamo direttamente il raggio di convergenza $R$:
		\[
			R = \lim_{n \to +\infty} \left| \frac{a_n}{a_{n+1}} \right| = \lim_{n \to +\infty} \frac{\frac{1}{n!}}{\frac{1}{(n+1)!}} = \lim_{n \to +\infty} (n+1) = +\infty
		\]
		Essendo $R = +\infty$, la serie converge per ogni $x \in \R$.
	
	\end{description}
	
	Vediamo un esempio dove $x_0 \neq 0$.
	Consideriamo la serie $\sum_{n=0}^{\infty} \frac{n^2+4n+1}{n^2+7} (x-2)^n$, dove $a_n = \frac{n^2+4n+1}{n^2+7}$ e $x_0 = 2$.
	
	Proviamo a studiarla con il criterio della radice:
	\[
		\lim_{n \to +\infty} \sqrt[n]{\frac{n^2+4n+1}{n^2+7}} = 1 = l, \quad R = \frac{1}{l} = \frac{1}{1} = 1
	\]
	Ricordiamo che la radice $n$-esima di un qualsiasi polinomio, non importa il grado, è sempre $1$.
	
	Siccome $R=1$, sappiamo che la serie converge in:
	\[
		(x_0 - 1, x_0 + 1) = (2-1, 2+1) = (1, 3)
	\]
	
	Ora bisogna calcolare la serie agli estremi e stabilirne il comportamento:
	
	\begin{itemize}
		
		\item Per $x = 1$:
		\[
			\sum_{n=0}^{\infty} \frac{n^2+4n+1}{n^2+7} (1-2)^n = \sum_{n=0}^{\infty} \frac{n^2+4n+1}{n^2+7} (-1)^n
		\]
		In $x = 1$ la serie è indeterminata, in quanto è a segni alterni ma la successione non rispetta i criteri di Leibniz.
		
		\item Per $x = 3$:
		\[
			\sum_{n=0}^{\infty} \frac{n^2+4n+1}{n^2+7} (3-2)^n = \sum_{n=0}^{\infty} \frac{n^2+4n+1}{n^2+7} (1)^n = \sum_{n=0}^{\infty} \frac{n^2+4n+1}{n^2+7}
		\]
		In $x = 3$ otteniamo una serie a termini positivi (dunque non può essere indeterminata) la cui successione però non tende a $0$, e dunque non può convergere.
	
	\end{itemize}
	
	Infine, possiamo dire che la serie $\sum_{n=0}^{\infty} \frac{n^2+4n+1}{n^2+7} (x-2)^n$ converge $\forall x \in (1, 3)$ e non converge per $x \le 1 \lor x \ge 3$.
	
	Come detto nell'introduzione alle serie di potenze, nell'intervallo in cui queste convergono possiamo definire una funzione.
	Formalizzando:
	\begin{itemize}
		
		\item Se $R = 0$, la serie converge in un solo punto, dunque è banale.
		
		\item Se $R > 0$:
		\[
			S(x) = \sum_{n=0}^{\infty} a_n \cdot (x - x_0)^n \quad \forall x \in (x_0 - R, x_0 + R)
		\]
		Per includere o meno gli estremi, bisogna calcolare la corrispondente serie numerica in quei punti.
		
	\end{itemize}
	
	\begin{teorema}{Continuità della funzione somma}{}
		
		$S(x)$ è continua in $(x_0 - R, x_0 + R)$.
		Inoltre, il \textbf{teorema di Abel} ci dice che se la serie converge agli estremi, allora è continua fino all'estremo in cui converge.
		
	\end{teorema}
	
	\begin{teorema}{Derivabilità della funzione somma}{}
		
		$S(x)$ è derivabile infinite volte in $(x_0 - R, x_0 + R)$.
		La derivata si indica come:
		\[
			S'(x) = \sum_{n=1}^{\infty} n \cdot a_n \cdot (x - x_0)^{n-1}
		\]
		
	\end{teorema}
	
	Sviluppando per esteso la serie e la sua derivata:
	\[
		\begin{aligned}
			S(x) &= a_0 + a_1(x - x_0) + a_2(x - x_0)^2 + a_3(x - x_0)^3 + \dots \\
			S'(x) &= 0 + a_1 + 2a_2(x - x_0) + 3a_3(x - x_0)^2 + \dots
		\end{aligned}
	\]
	
	Essendo $S(x)$ derivabile infinite volte, possiamo scrivere anche la derivata seconda, terza, quarta, $\dots$, $n$-esima, $\dots$
	Ad esempio:
	\[
		S''(x) = \sum_{n=2}^{\infty} n(n-1) \cdot a_n \cdot (x - x_0)^{n-2} = 2a_2 + 6a_3(x - x_0) + \dots
	\]
	
	Notiamo inoltre una cosa particolare: se proviamo a calcolare $S(x)$ e tutte le sue derivate nel punto $x = x_0$, otteniamo ogni volta i termini che contengono il fattore $(x - x_0)$ si annullano e scompaiono:
	\[
		S(x_0) = a_0, \quad S'(x_0) = a_1, \quad S''(x_0) = 2a_2, \quad S'''(x_0) = 6a_3, \dots
	\]
	
	Individuiamo un pattern: rimangono solo i termini della successione il cui indice è lo stesso del grado della derivata, e il coefficiente è il fattoriale del grado stesso. Formalizzando:
	\[
		S^k(x_0) = k! \cdot a_k \implies a_k = \frac{S^k(x_0)}{k!}
	\]
	
	Così facendo, possiamo scrivere un'identità che stabilisce il legame tra $S(x)$, le sue derivate e i suoi coefficienti:
	\[
		S(x) = \sum_{n=0}^{+\infty} \frac{S^n(x_0)}{n!} (x - x_0)^n
	\]
	
	Ci teniamo ora a fare un appunto sulla potenza della serie geometrica:
	\[
		\sum_{n=0}^{+\infty} x^n = \frac{1}{1-x}
	\]
	per $x \in (-1, 1)$, che è chiaramente una serie di potenze con $a_n = 1$ e $x_0 = 0$. Possiamo dunque definire $f(x) = \frac{1}{1-x}$ per $x \in (-1, 1)$.
	
	Se invece volessimo fare il percorso al contrario e partire da $g(x) = \frac{1}{1+x}$ e ricavarne la serie?
	\[
		\frac{1}{1+x} = \frac{1}{1 - (-x)} = \sum_{n=0}^{+\infty} (-x)^n = \sum_{n=0}^{+\infty} (-1 \cdot x)^n = \sum_{n=0}^{+\infty} (-1)^n \cdot x^n
	\]
	dove $a_n = (-1)^n$ e $x_0 = 0$, definita per $x \in (-1, 1)$.
	
	Adesso vogliamo di nuovo partire da una serie...
	\[
		T(x) = \sum_{n=0}^{+\infty} \frac{(-1)^n \cdot x^{n+1}}{n+1}
	\]
	dove $a_n = \frac{(-1)^n}{n+1}$ e $x_0 = 0$.
	
	Studiamoci il raggio di convergenza:
	\[
		l = \lim_{n \to +\infty} \sqrt[n]{|a_n|} = \lim_{n \to +\infty} \sqrt[n]{\frac{1}{n+1}} = 1 \implies R = \frac{1}{l} = 1
	\]
	
	Senza perderci in calcoli otteniamo che la serie converge in $(-1, 1)$ e, studiando gli estremi, converge anche in $x = 1$.
	
	Se proviamo ora a fare la derivata di questa serie, otterremo una cosa interessante:
	\[
		T'(x) = \sum_{n=0}^{\infty} \frac{(-1)^n (n+1) x^n}{n+1} = \sum_{n=0}^{\infty} (-1)^n x^n = \frac{1}{1+x}
	\]
	Come vediamo, $T(x)$ è una funzione definita in $(-1, 1]$ che, se derivata, ci restituisce $1 / (1+x)$.
	
	Questa funzione è proprio $\ln(1+x)$, dunque risulta vera la seguente identità:
	\[
		\ln(1+x) = \sum_{n=0}^{\infty} \frac{(-1)^n x^{n+1}}{n+1} = x - \frac{x^2}{2} + \frac{x^3}{3} - \frac{x^4}{4} \dots
	\]
	ma solo per $x \in (-1, 1]$.
	
	E se invece, presa $\frac{1}{1-y}$, ponessimo $y = x^2$?
	La corrispondente serie sarebbe $\sum_{n=0}^{\infty} y^n = \sum_{n=0}^{\infty} x^{2n}$ che, scritta per estesa, risulta $1 + x^2 + x^4 + x^8 + \dots$
	
	Se prendiamo dunque $f(x) = \frac{1}{1-y}$, per $x \in (-1, 1]$, abbiamo che $y = (-x^2)$ e la corrispondente serie sarà:
	\[
		\sum_{n=0}^{\infty} (-1)^n x^{2n} = 1 - x^2 + x^4 - x^6 \dots
	\]
	
	Proviamo a fare il discorso di prima:
	\[
		S(x) = \sum_{n=0}^{\infty} \frac{(-1)^n x^{2n+1}}{2n+1} \implies S'(x) = \sum_{n=0}^{\infty} \frac{(-1)^n (2n+1) x^{2n}}{2n+1} = \sum_{n=0}^{\infty} (-1)^n x^{2n}
	\]
	Dunque $S(x)$ è quella funzione che, se derivata, è proprio uguale a $1 / (1+x^2)$: è l'arcotangente.
	\[
		\operatorname{arctg}(x) = \sum_{n=0}^{\infty} \frac{(-1)^n x^{2n+1}}{2n+1} = x - \frac{x^3}{3} + \frac{x^5}{5} - \frac{x^7}{7} \dots
	\]
	ma solo per $x \in (-1, 1]$.
	
	Introduciamo ora lo sviluppo per le funzioni seno e coseno:
	\begin{itemize}
		
		\item $\sin(x) = \sum_{n=0}^{\infty} \frac{(-1)^n \cdot x^{2n+1}}{(2n+1)!} = x - \frac{x^3}{3!} + \frac{x^5}{5!} - \frac{x^7}{7!} + \dots$
		
		\item $\cos(x) = \sum_{n=0}^{\infty} \frac{(-1)^n \cdot x^{2n}}{(2n)!} = 1 - \frac{x^2}{2!} + \frac{x^4}{4!} - \frac{x^6}{6!} + \dots$
		
	\end{itemize}
	
	Notiamo così anche il legame che intercorre tra il seno e la sua derivata, il coseno (la cui verifica viene omessa ma risulta facilmente dimostrabile).
	
	Un'ultima importante nota: se volessimo calcolare $\sin(2x)$, $\cos(x^2)$, $\ln(1 + \sqrt{x})$ (in poche parole, con un argomento che non è il solito $x$), non serve ogni volta ricalcolare le derivate da capo, ma basta sostituire il nuovo argomento nello sviluppo noto.
	Ad esempio:
	\[
		\sin(2x) = 2x - \frac{(2x)^3}{3!} + \frac{(2x)^5}{5!} - \dots
	\]
	
	Attenzione però! In tal caso cambia anche il dominio di definizione della funzione analitica. Ad esempio, per $\ln(1 + 2x)$, sapendo che l'argomento deve essere compreso in $(-1, 1]$, avremo $2x \in (-1, 1] \implies x \in (-\frac{1}{2}, \frac{1}{2}]$.
	
	Ovviamente possiamo anche percorrere il cammino inverso:
	\[
		f(x) = \frac{1}{2 - 5x} = \frac{1}{2} \cdot \frac{1}{1 - \frac{5}{2}x}
	\]
	Riconosciamo nella seconda parte una serie geometrica con ragione $y = \frac{5}{2}x$. Riscriviamo dunque come:
	\[
		f(x) = \frac{1}{2} \cdot \sum_{n=0}^{\infty} \left(\frac{5}{2}x\right)^n = \frac{1}{2} \cdot \sum_{n=0}^{\infty} \left(\frac{5}{2}\right)^n \cdot x^n
	\]
	In questo caso la successione dei coefficienti è $a_n = \frac{1}{2} \cdot \left(\frac{5}{2}\right)^n$ e il centro è $x_0 = 0$.
	
	Cambia chiaramente anche il dominio di convergenza della serie:
	\[
		\frac{5}{2}x \in (-1, 1) \implies x \in \left(-\frac{2}{5}, \frac{2}{5}\right)
	\]
	
	Con queste informazioni possiamo calcolare la derivata di qualsiasi ordine sfruttando i coefficienti della serie di potenze.
	
	\chapter{Funzioni analitiche}
	
	Avendo scritto $\arctan(x)$ ed $e^x$ come polinomi, possiamo anche riscrivere approssimazioni di $\pi$ ed $e$ sotto forma di serie:
	\begin{itemize}
		
		\item $\arctan(x) = x - \frac{x^3}{3} + \frac{x^5}{5} - \frac{x^7}{7} \dots$
		
		$\arctan(1) = \frac{\pi}{4} = 1 - 1/3 + 1/5 - 1/7 \dots$
		
		Dunque $\pi = 4(1 - 1/3 + 1/5 - 1/7 \dots)$
		
		\item $e^x = 1 + x + \frac{x^2}{2!} + \frac{x^3}{3!} + \dots$
		
		$e^1 = e = 1 + 1 + 1/2 + 1/3! + 1/4! + \dots$
		
	\end{itemize}
	
	Presa $f \in C^\infty(I)$ (significa "derivabile infinite volte nell'intervallo $I \subseteq \R$").
	
	Preso $x_0 \in I$, possiamo riscrivere $f(x_0)$ come:
	\[
		f(x) = f(x_0) + f'(x_0)(x - x_0) + \frac{f''(x_0)(x - x_0)^2}{2!} + \frac{f'''(x_0)(x - x_0)^3}{3!} + \dots
	\]
	\[
		+ \frac{f^{(k)}(x_0)(x - x_0)^k}{k!} + R_k(x_0, x)
	\]
	dove $R_k(x_0, x)$ è un resto dipendente dal grado al quale ci siamo fermati.
	
	Se invece volessimo ottenere l'identità, potremmo usare ciò che abbiamo imparato con le serie di potenze:
	\[
		f(x) = \sum_{n=0}^{\infty} \frac{f^{(n)}(x_0)}{n!} (x - x_0)^n
	\]
	
	Se \textbf{vale} quest'identità, diciamo che $f$ è \textbf{analitica} in $I$.
	Per gli esempi visti prima, diciamo che $\ln(1+x)$ e $\arctan(x)$ sono entrambe analitiche in $(-1, 1]$.
	
	\chapter{Funzioni nel campo complesso}
	Di solito, con "funzioni nel campo complesso" si intendono funzioni definite come $f: \C \to \C$, tuttavia possiamo considerare quelle definite come $f: \R \to \C$.
	
	Queste ultime sono formalmente definite come segue:
	\[
		f(t) := f_1(t) + i \cdot f_2(t)
	\]
	dove $f_1(t)$ è la parte reale e $f_2(t)$ è la parte immaginaria.
	
	Per linearità dell'operazione di derivazione abbiamo:
	\[
		f'(t) := f_1'(t) + i \cdot f_2'(t)
	\]
	
	Siamo adesso curiosi di riscrivere $e^z$, con $z \in \C$, come serie di potenze:
	\[
		e^z = \sum_{n=0}^{\infty} \frac{z^n}{n!} \quad \forall z \in \C
	\]
	il che è valido, in quanto la potenza è una serie di moltiplicazioni (operazione definita anche in $\C$) e il fattoriale è definito per i numeri naturali di cui $n$ fa parte.
	
	Anche nei complessi vale la proprietà per cui:
	\[
		e^{z_1} \cdot e^{z_2} = e^{z_1 + z_2} = \sum_{n=0}^{\infty} \frac{(z_1 + z_2)^n}{n!}
	\]
	
	Possiamo dunque generalizzare dicendo $e^z = e^{a+ib} = e^a \cdot e^{ib}$. 
	Notiamo che $e^a$ è un numero reale, dunque possiamo in un certo senso "trascurarlo", mentre per $e^{ib}$ proviamo a scriverlo come serie di potenze:
	\[
		e^{ib} = \sum_{n=0}^{\infty} \frac{(ib)^n}{n!} = \sum_{n=0}^{\infty} \frac{(i)^n \cdot b^n}{n!}
	\]
	vediamo che $(i)^n$ ciclicamente diventa:
	\[
		i^0 = 1, \quad i^1 = i, \quad i^2 = -1, \quad i^3 = -i, \quad i^4 = 1, \quad i^5 = i
	\]
	
	Se sviluppiamo la sommatoria otteniamo dunque:
	\[
		\sum_{n=0}^{\infty} \frac{(i)^n \cdot b^n}{n!} = \frac{i^0 b^0}{0!} + \frac{i^1 b^1}{1!} + \frac{i^2 b^2}{2!} + \frac{i^3 b^3}{3!} + \frac{i^4 b^4}{4!} + \dots = 1 + ib - \frac{b^2}{2!} - \frac{ib^3}{3!} + \frac{b^4}{4!} + \dots
	\]
	
	Se continuassimo all'infinito e raggruppassimo tra loro tutte le parti reali e complesse:
	\[
		\underbrace{\left(1 - \frac{b^2}{2!} + \frac{b^4}{4!} - \frac{b^6}{6!} \dots\right)}_{\text{sviluppo in serie di } \cos(b)} + i \cdot \underbrace{\left(b - \frac{b^3}{3!} + \frac{b^5}{5!} - \frac{b^7}{7!} \dots\right)}_{\text{sviluppo in serie di } \sin(b)}
	\]
	
	Ciò vuol dire che $e^{ib} = \cos(b) + i \cdot \sin(b)$. Questa viene chiamata \textbf{formula di Eulero}, e può essere riutilizzata per ogni possibile valore di $b$.
	
	Tornando a monte, possiamo dire:
	\[
		e^z = e^{a+ib} = e^a \cdot e^{ib} = e^a (\cos(b) + i \cdot \sin(b))
	\]
	
	Dimostriamo ora che anche la derivata rimane la stessa anche nel campo complesso:
	\begin{itemize}
		\item Sia $\C \ni \lambda = \lambda_1 + \lambda_2$
		\item $e^{\lambda t} = \sum_{n=0}^{\infty} \frac{(\lambda t)^n}{n!} = \sum_{n=0}^{\infty} \frac{\lambda^n \cdot t^n}{n!}$
	\end{itemize}
	\[
	(e^{\lambda t})' = \sum_{n=1}^{\infty} \frac{\lambda^n \cdot n \cdot t^{n-1}}{n \cdot (n-1)!} = \sum_{n=1}^{\infty} \frac{\lambda \cdot \lambda^{n-1} \cdot t^{n-1}}{(n-1)!} = \lambda \cdot \sum_{n=1}^{\infty} \frac{\lambda^{n-1} \cdot t^{n-1}}{(n-1)!}
	\]
	
	Se effettuiamo un cambio di variabile e poniamo $k = n-1$:
	\[
	(e^{\lambda t})' = \lambda \cdot \sum_{k=0}^{\infty} \frac{\lambda^k \cdot t^k}{k!}
	\]
	% --- NEW PAGE ---
	
	\chapter{Integrali secondo Riemann}
	
	L'idea dell'integrale nasce dal bisogno di calcolare aree irregolari.
	Per far ciò, da cosa partiamo? Ovviamente da calcoli di aree regolari che sappiamo già fare.
	\begin{itemize}
		
		\item Sia $R$ un rettangolo, $A(R) = b \cdot h$.
		
		\item $R_1 \cup R_2$, con $(\mathring{R_1} \cap \mathring{R_2}) = \emptyset$ (insieme dei punti in comune tra i due rettangoli).
		$A(R_1 \cup R_2) = A(R_1) + A(R_2)$.
		
	\end{itemize}
	
	Generalizzando...
	\begin{itemize}
		
		\item Sia $f: [a, b] \to \R$, $f \geq 0$ limitata.
		\begin{definizione}{Sottografico di una funzione $f$}{}
			
			Definiamo il \textbf{sottografico di $f$} come l'insieme dei punti sottostanti o coincidenti al grafico di $f$:
			\[
				R_f := \{(x, y) \mid a \leq x \leq b, \, 0 \leq y \leq f(x)\}
			\]
			
		\end{definizione}
		
	\end{itemize}
	
	Denotiamo $S_f$ come l'area di $f$. Se invece volessimo calcolare $A(R_f)$?
	\begin{itemize}
		
		\item Se $f(x) = K$, $\forall x \in [a, b]$:
		\[
			S_f = \text{area del rettangolo } [a, b] \times [0, K] = (b - a) \cdot K
		\]
		
		\item Se $f(x) = \begin{cases} \alpha_1, & \text{se } x \in [a, c) \\ \alpha_2, & \text{se } x \in [c, b] \end{cases}$:
		\[
			S_f = \text{somma aree dei due rettangoli} = (c - a) \cdot \alpha_1 + (b - c) \cdot \alpha_2
		\]
		
	\end{itemize}
	
	Per ottenere l'area di una superficie irregolare, l'idea è di approssimare la regione come somma delle aree dei vari rettangoli costruiti.
	
	\begin{definizione}{Partizione di un intervallo}{}
		
		Una partizione $P$ di un intervallo $[a, b]$ è un insieme $P = \{x_0, x_1, \dots, x_n\}$ tale che:
		\[
			a = x_0 < x_1 < x_2 < \dots < x_n = b
		\]
		Definiamo inoltre l'ampiezza di $P$ come la distanza massima che intercorre tra due punti:
		\[
			|P| := \max\{x_k - x_{k-1}, \, k = 1, \dots, n\}
		\]
		
	\end{definizione}
	
	\begin{osservazione}{}{}
		
		Consideriamo il caso particolare in cui suddividiamo l'intervallo in parti uguali:
		\[
			h = \frac{b-a}{n}, \quad x_0 = a
		\]
		dunque avremo che $x_1 = a + h,\, x_2 = a + 2h, \dots,\, x_n = b = a + nh$.
		
	\end{osservazione}
	
	\begin{definizione}{Funzione costante a tratti}{}
		
		Una funzione $\varphi_p: [a, b] \to \R$ è costante a tratti se esiste una partizione $P$ e un insieme $\{a_1, \dots, a_n\}$ tale che $\varphi_p(x) = a_k$ se $x \in [x_{k-1}, x_k], \, k = 1, \dots, n$.
		
	\end{definizione}
	
	Per una funzione costante a tratti, l'area sottesa si calcola come somma delle aree dei rettangoli:
	\[
		S(\varphi_p) = \sum_{k=1}^n a_k (x_k - x_{k-1})
	\]
	
	Osserviamo che la definizione di funzione costante a tratti è valida anche per $a_k < 0$. Tuttavia in tal caso avremo somme tra elementi anche negativi, e dunque non ha molto senso parlare di area nel senso euclideo del termine, bensì parliamo di \textbf{area con segno}.
	
	Vogliamo ora dare un'approssimazione dell'area di un sottografico.
	\begin{itemize}
		
		\item Definiamo $f: [a, b] \to \mathbb{R}$ limitata.
		
		\item Fissiamo una partizione $P = \{x_0, \dots, x_n\}$.
		
		\item Definiamo $\varphi_P: [a, b] \to \mathbb{R}$ come:
		\[
			\varphi_P(x) = \inf(f(x)) \quad \text{per } x \in [x_{k-1}, x_k]
		\]
		che chiamiamo elemento $m_k$.
		
		\item Definiamo $\varphi^P: [a, b] \to \mathbb{R}$ come:
		\[
			\varphi^P(x) = \sup(f(x)) \quad \text{per } x \in [x_{k-1}, x_k]
		\]
		che chiamiamo elemento $M_k$.
		
	\end{itemize}
	
	Si nota come valga sempre la relazione $\varphi_P \le f \le \varphi^P$.
	
	\begin{osservazione}{}{}
		
		Geometricamente, i valori $m_k$ ed $M_k$ rappresentano rispettivamente l'altezza del rettangolo inferiore e del rettangolo superiore costruiti sull'intervallo $[x_{k-1}, x_k]$.
		
	\end{osservazione}
	
	Introduciamo ora la seguente notazione per le somme delle aree:
	\begin{itemize}
		
		\item $S_{\varphi_P} = S_P f$, la somma dei rettangolini "sotto" alla funzione.
		
		\item $S^{\varphi^P} = S^P f$, la somma dei rettangolini "sopra" alla funzione.
		
	\end{itemize}
	
	Abbiamo dunque che $S_P f \le S f \le S^P f$.
	
	Nonostante non abbiamo trovato un valore preciso, abbiamo trovato un'approssimazione per l'area del sottografico della funzione.
	
	Per migliorare questa approssimazione basta infittire la partizione, aumentando il numero di punti e di conseguenza diminuendo la distanza tra questi.
	Possiamo dunque definire:
	\begin{itemize}
		
		\item $s_P f = \sum_{k=1}^n m_k (x_k - x_{k-1})$ le somme integrali per difetto.
		
		\item $S^P f = \sum_{k=1}^n M_k (x_k - x_{k-1})$ le somme integrali per eccesso.
		
	\end{itemize}
	
	Se $Q$ è un raffinamento di $P$, contiene più punti di $P$, cioè $P \subseteq Q$.
	In questo modo abbiamo $s_P f \le s_Q f$ e $S^Q f \le S^P f$.
	
	Strategia per il calcolo del valore dell'area del sottografico di $f$:
	\begin{enumerate}
		
		\item Scelgo una successione di partizioni $P_n$ t.c. al crescere di $n$ la partizione si infittisce. Dunque, per $n \to +\infty$, la distanza tende a $0$.
		
		\item Calcolo dunque $\lim_{n \to +\infty} S^P f$ e $\lim_{n \to +\infty} s_P f$.
		
	\end{enumerate}
	Se i limiti coincidono, abbiamo trovato esattamente il valore dell'area del sottografico di $f$.
	
	\begin{osservazione}{}{}
		La funzione presa in considerazione è $f(x) = e^x$ in $[a, b]$.
		Come partizione prendo quella definita nell'esempio:
		\[
		P_n = \{x_k = a + kh,\, k = \{0 \dots n\}\}, \quad h := |P_n| = \frac{b-a}{n} \xrightarrow{n \to +\infty} 0
		\]
		Calcoliamo le somme integrali:
		\begin{itemize}
			\item $m_k = \inf e^x = e^{x_{k-1}} = e^{a + kh - 1}$ per $x \in [x_k, x_{k-1}]$
			\item $M_k = \sup e^x = e^{x_k} = e^{a + kh}$ per $x \in [x_k, x_{k-1}]$
		\end{itemize}
		
		\[
			S_f^\mathcal{P} = \sum_{k=1}^{n} e^{a + h(k-1)} \cdot h = e^a \cdot h \cdot \sum_{k=1}^{n} (e^h)^{k-1}
		\]
		Riconosciamo una serie geometrica, con base $e^h$ ed esponente $k-1$:
		\[
			= e^a \cdot h \cdot \sum_{j=0}^{n} (e^h)^e = e^a \cdot h \cdot \frac{1 - (e^h)^n}{1 - e^h} = \frac{h}{1 - e^h} \cdot e^a \left(1 - (e^{\frac{b-a}{n}})^n\right) = \frac{h}{1 - e^h} \cdot e^a (1 - e^{b-a})
		\]
		\[
			= \frac{h}{1 - e^h} \cdot e^a - e^b = \frac{h}{1 - e^h} \cdot (-1) (e^b - e^a) = \frac{h}{e^h - 1} (e^b - e^a)
		\]
		\[
			\xrightarrow[n \to +\infty \text{ e di conseguenza per } h \to 0]{} e^b - e^a
		\]
		
	\end{osservazione}
	
	Omettiamo la dimostrazione delle somme per eccesso che, ed è facilmente verificabile, daranno lo stesso risultato.
	Potremmo dunque dire che $S e^x = e^b - e^a$, per un opportuno intervallo $[a, b]$.
	
	\begin{definizione}{Funzione integrabile}{}
		
		Sia $f: [a, b] \to \R$ limitata.
		Se per una successione di partizioni $\mathcal{P}_n$ tale che $|\mathcal{P}_n| \to 0$ esistono finiti e coincidono i seguenti limiti:
		\[
			\lim_{n \to +\infty} S^{\mathcal{P}_n} f = \lim_{n \to +\infty} S_{\mathcal{P}_n} f = l
		\]
		allora $f$ è integrabile in $[a, b]$ e denotiamo come $\int_{a}^{b} f(x) \, dx$ l'integrale di $f$ tra $a$ e $b$.
		
	\end{definizione}
	
	\begin{osservazione}{}{}
		
		Non importa quale sia il "nome" della variabile per cui scegliamo di integrare, l'importante è che questo sia condiviso dalla funzione e dal differenziale.
		\[
			\int_{a}^{b} f(x) \, dx = \int_{a}^{b} f(y) \, dy = \int_{a}^{b} f(u) \, du = \dots
		\]
		
	\end{osservazione}
	
	Abbiamo già visto come l'esponenziale sia integrabile in un generico intervallo $[a, b]$.
	Tuttavia, anche le \textbf{funzioni costanti a tratti sono integrabili}: per studiarne le somme parziali basta scegliere una partizione che contenga gli estremi degli intervalli.
	
	\begin{teorema}{Versione semplice del Teorema Fondamentale del Calcolo Integrale}{}
		
		Siano $f, F$ due funzioni tali che:
		\begin{itemize}
			
			\item $f$ è lipschitziana in $[a, b]$ (se questo requisito può spaventare, può essere sostituito da "$f$ derivabile in $[a, b]$ con derivata continua").
			
			\item $F$ è derivabile con $F' = f$.
		
		\end{itemize}
		Allora $f$ è integrabile in $[a, b]$ e vale:
		\[
			\int_{a}^{b} f(x) \, dx = F(b) - F(a)
		\]
		
	\end{teorema}
	
	Abbiamo chiamato "semplice" questa versione del teorema in quanto il requisito per $f$ è molto restrittivo.
	
	\begin{dimostrazione}{}{}
		Calciamoci le somme per difetto, usando la formula che abbiamo visto prima: $S_{P_n} = \sum_{k=1}^{n} m_k (x_k - x_{k-1})$. Siccome per ipotesi $f$ è continua, $m_k$ non è più solo l'estremo inferiore dell'intervallo $[x_{k-1}, x_k]$, ma è proprio il minimo.
		Per contraddistinguerlo, lo indichiamo come segue:
		\[
			f(\xi_k) := \min(f(x_k), f(x_{k-1}))
		\]
		L'uguaglianza diventa dunque:
		\[
			S_{P_n} = \sum_{k=1}^{n} f(\xi_k) (x_k - x_{k-1})
		\]
		Siccome per ipotesi abbiamo $F$ derivabile, possiamo applicare il teorema di Lagrange nel seguente modo:
		\[
			F(x_k) - F(x_{k-1}) = F'(\eta_k)(x_k - x_{k-1})
		\]
		Ricordiamo per ipotesi che $F' = f$, dunque $F'(\eta_k) = f(\eta_k)$. Riscriviamo l'uguaglianza aggiungendo e togliendo $f(\eta_k)$:
		\[
			S_{P_n} = \sum_{k=1}^{n} \left(f(\xi_k) + f(\eta_k) - f(\eta_k)\right)(x_k - x_{k-1})
		\]
		Spezziamo la sommatoria:
		\[
			S_{P_n} = \sum_{k=1}^{n} (f(\xi_k) - f(\eta_k))(x_k - x_{k-1}) + \sum_{k=1}^{n} f(\eta_k)(x_k - x_{k-1})
		\]
		\begin{center}
			\small
			\textit{(siccome $f$ è Lipschitziana, allora $f(\xi_k) - f(\eta_k)$ è una quantità che tende a $0$, dunque l'intera somma tende a $0$)} $\qquad$ \textit{(il contenuto corrisponde a $F(x_k) - F(x_{k-1})$)}
		\end{center}
		
		Fatte le dovute considerazioni, espandiamo la sommatoria per osservare se troviamo un pattern:
		\[
			S_{P_n} = F(x_1) - F(x_0) + F(x_2) - F(x_1) + \dots + F(x_n) - F(x_{n-1}) = 
		\]
		\[
			= F(x_n) - F(x_0) = F(b) - F(a)
		\]
		La dimostrazione è analoga per le somme superiori.
	
	\end{dimostrazione}
	
	Alcuni esempi:
	\begin{itemize}
		
		\item $F(x) = \frac{x^2}{2}$, $F' = f = x \implies \int_{a}^{b} x \, dx = \frac{b^2}{2} - \frac{a^2}{2}$
		
		\item $F(x) = \frac{x^{n+1}}{n+1}$, $F' = f = x^n \implies \int_{a}^{b} x^n \, dx = \frac{b^{n+1}}{n+1} - \frac{a^{n+1}}{n+1}$, regola generale per i polinomiali.
		
		\item $F(x) = -\cos x$, $F' = f = \sin x \implies \int_{a}^{b} \sin x \, dx = -\cos(b) + \cos(a)$
		
		\item $F(x) = \sin x$, $F' = f = \cos x \implies \int_{a}^{b} \cos x \, dx = \sin(b) - \sin(a)$
		
		\item $F(x) = \arctan x$, $F' = f = \frac{1}{1+x^2} \implies \int_{a}^{b} \frac{1}{1+x^2} \, dx = \arctan(b) - \arctan(a)$
		
		\item $F(x) = e^x$, $F' = f = e^x \implies \int_{a}^{b} e^x \, dx = e^b - e^a$
		
	\end{itemize}
	
	\begin{osservazione}{}{}
		Per indicare l'integrale calcolato nell'intervallo $[a, b]$ scriviamo nel seguente modo:
		\[
			\int_{a}^{b} f(x) \, dx = [F(x)]_{x=a}^{x=b} = [F(x)]_{a}^{b}
		\]
		e chiamiamo questo integrale con gli estremi \textbf{integrale definito}.
	\end{osservazione}
	
	\section{Proprietà degli integrali}
	
	\begin{description}
	
		\item [Linearità:] Siano $f, g$ integrabili in $[a, b]$. 
		Allora anche $\lambda f + \mu g$ è integrabile in $[a, b]$ $\forall \lambda, \mu \in \R$:
		\[
			\int_{a}^{b} (\lambda f(x) + \mu g(x)) \, dx = \lambda \int_{a}^{b} f(x) \, dx + \mu \int_{a}^{b} g(x) \, dx
		\]
		
		\item [Additività:] Sia $f$ integrabile in $[a, b]$. Sia $c \in (a, b)$. Allora $f$ è integrabile in $[a, c]$ e $[c, b]$. Inoltre, è possibile stabilire la seguente relazione:
		\[
			\int_{a}^{b} f(x) \, dx = \int_{a}^{c} f(x) \, dx + \int_{c}^{b} f(x) \, dx
		\]
	
	\end{description}
	
	Graficamente:
	\begin{center}
		
		\begin{tikzpicture}[scale=1.1]
			
			\draw[->, thick] (-0.5, 0) -- (4, 0);
			\draw[->, thick] (0, -0.5) -- (0, 3);
			
			\draw[domain=0.5:3.5, smooth, variable=\x, blue, thick] plot ({\x}, {1.5 + 0.5*sin(\x*120)});
			
			\draw[dashed, thick] (1, 0) node[below] {$a$} -- (1, 1.83);
			\draw[dashed, thick] (2, 0) node[below] {$c$} -- (2, 1.25);
			\draw[dashed, thick] (3, 0) node[below] {$b$} -- (3, 1.93);
			
			\fill[orange, opacity=0.3] (1, 0) -- plot[domain=1:2, smooth, variable=\x] ({\x}, {1.5 + 0.5*sin(\x*120)}) -- (2, 0) -- cycle;
			\fill[cyan, opacity=0.3] (2, 0) -- plot[domain=2:3, smooth, variable=\x] ({\x}, {1.5 + 0.5*sin(\x*120)}) -- (3, 0) -- cycle;
			
		\end{tikzpicture}
		
	\end{center}
	
	Abbiamo "spezzato" l'area del sottografo di $f$, calcolando due parti separate e poi sommandole.
	
	\begin{osservazione}{}{}
		
		Se poniamo $c = a$, la formula diventa:
		\[
			\int_{a}^{b} f(x) \, dx = \int_{a}^{a} f(x) \, dx + \int_{a}^{b} f(x) \, dx
		\]
		Risolvendo l'equazione (cioè eliminando i termini uguali a dx. e sx.) otteniamo:
		\[
			\int_{a}^{a} f(x) \, dx = 0
		\]
		Il che diventa una vera e propria definizione, in particolare quella dell'integrale calcolato in un punto solo.
		
	\end{osservazione}
	
	\begin{osservazione}{}{}
		
		Se poniamo $b = a$, la formula diventa:
		\[
			\int_a^a f(x) \, dx = \int_a^c f(x) \, dx + \int_c^a f(x) \, dx
		\]
		dove il primo membro è per definizione pari a $0$. Riordinando abbiamo che:
		\[
			\int_a^c f(x) \, dx = - \int_c^a f(x) \, dx
		\]
		Questo ci dice che se invertiamo gli estremi di integrazione, si inverte anche il segno dell'integrale.
		\smallbreak
		Ad esempio: $\int_1^2 e^x \, dx = - \int_2^1 e^x \, dx$.
		
	\end{osservazione}
	
	\begin{proposizione}{}{}
		
		Siano $f, g$ integrabili in $[a, b]$ e sia $f(x) \leq g(x)$ per ogni $x \in [a, b]$.
		Allora risulta:
		\[
			\int_a^b f(x) \, dx \leq \int_a^b g(x) \, dx
		\]
		
	\end{proposizione}
	
	Geometricamente, se $f$ è sempre sotto $g$, anche l'area del sottografico di $f$ sarà minore o uguale a quella del sottografico di $g$.
	
	Una conseguenza importante della monotonia è la seguente disuguaglianza per il valore assoluto dell'integrale:
	\[
		\left| \int_a^b f(x) \, dx \right| \leq \int_a^b |f(x)| \, dx
	\]
	
	\begin{dimostrazione}{}{}
		
		In generale sappiamo che $-|z| \leq z \leq |z|$. Se come $z$ prendiamo $f(x)$ abbiamo:
		\[
			-|f(x)| \leq f(x) \leq |f(x)|
		\]
		Applicando linearità e monotonia dell'integrale otteniamo:
		\[
			-\int_a^b |f(x)| \, dx \leq \int_a^b f(x) \, dx \leq \int_a^b |f(x)| \, dx
		\]
		che è proprio quanto si voleva dimostrare.
		
	\end{dimostrazione}
	
	Altre categorie di \textbf{funzioni integrabili} sono le \textbf{funzioni monotone}, cioè quelle che mantengono un determinato comportamento per ogni $x$ del dominio. In particolare:
	\begin{itemize}
		
		\item \textbf{Monotona crescente}: $x \leq y$ implica $f(x) \leq f(y)$ $\forall x, y \in [a, b]$.
		
		\item \textbf{Monotona decrescente}: $x \leq y$ implica $f(x) \geq f(y)$ $\forall x, y \in [a, b]$.
		
	\end{itemize}
	
	\begin{teorema}{Teorema di integrabilità delle funzioni monotone}{}
		
		Se $f: [a, b] \to \mathbb{R}$ è monotona, allora è integrabile in $[a, b]$.
		
	\end{teorema}
	
	\begin{dimostrazione}{}{}
		
		Fisso una partizione $P_n$, con $|P_n| = \frac{b-a}{n} = h$.
		
		Calciamoci $S^{P_n} - S_{P_n}$: se il risultato tende a $0$, allora le due somme convergono allo stesso valore e $f$ sarà dunque integrabile.
		\begin{align*}
			S^{P_n} &= \sum_{k=1}^{n} \sup_{x \in [x_{k-1}, x_k]} (f(x)) (x_k - x_{k-1}) \\
			S_{P_n} &= \sum_{k=1}^{n} \inf_{x \in [x_{k-1}, x_k]} (f(x)) (x_k - x_{k-1})
		\end{align*}
		
		Abbiamo detto che $f$ è monotona. Supponiamo sia crescente.
		Osserviamo che, siccome è monotona crescente, se abbiamo un intervallino con soli due elementi, $x_{k-1}$ e $x_k$, con $x_{k-1} \leq x_k$, allora anche $f(x_{k-1}) \leq f(x_k)$.
		
		Stabilito questo, notiamo che tra i due, l'estremo inferiore è proprio $f(x_{k-1})$ e analogamente per il superiore, $f(x_k)$.
		
		Riscriviamo dunque le somme integrali:
		\[
			S^{P_n} = \sum_{k=1}^{n} f(x_k) (x_k - x_{k-1}), \quad S_{P_n} = \sum_{k=1}^{n} f(x_{k-1}) (x_k - x_{k-1})
		\]
		Sottraendoli otteniamo $\sum_{k=1}^{n} (f(x_k) - f(x_{k-1})) (x_k - x_{k-1})$.
		
		Riconosciamo $f(x_k) - f(x_{k-1})$ come somma telescopica, che avrà dunque come risultato $f(x_n) - f(x_0) = f(b) - f(a)$.
		
		Infine, notiamo come $(x_k - x_{k-1})$ rappresenti la distanza tra due punti, cioè $h$. In quanto scritto così è un valore che non dipende dall'indice $k$, lo portiamo fuori.
		
		Restiamo con $h \cdot (f(b) - f(a))$. Per $n \to +\infty$, $h \to 0$, e dunque tutta la somma tende a $0$.
		
	\end{dimostrazione}
	
	\begin{osservazione}{}{}
		
		La combinazione lineare di più funzioni monotone (e dunque integrabili) è integrabile.
		
	\end{osservazione}
	
	\begin{proposizione}{}{}
		
		Le funzioni con un numero finito di cambi di monotonia sono integrabili.
		
	\end{proposizione}
	
	Prendiamo ad esempio la funzione $f(x) = |x|$ in $[-2, 3]$.
	Vediamo che ha un cambio solo. In certi casi, per calcolare l'integrale conviene spezzare per additività e calcolare l'integrale dei due comportamenti separatamente.
	
	\[
		\int_{-2}^{3} |x| \, dx = \int_{-2}^{0} |x| \, dx + \int_{0}^{3} |x| \, dx
	\]
	ricordiamo che, per definizione, $|x| = -x$ se $x < 0$ e $|x| = x$ se $x \geq 0$.
	
	Ci sbarazziamo dunque del modulo scrivendo:
	\[
		\int_{-2}^{3} |x| \, dx = -\int_{-2}^{0} x \, dx + \int_{0}^{3} x \, dx = \left[ -\frac{x^2}{2} \right]_{-2}^{0} + \left[ \frac{x^2}{2} \right]_{0}^{3}
	\]
	
	\begin{teorema}{Integrabilità delle funzioni continue}{}
		
		Le funzioni continue in $[a, b]$ sono integrabili.
		
	\end{teorema}
	
	\begin{teorema}{Integrabilità delle funzioni discontinue}{}
		
		Le funzioni limitate in $[a, b]$ con un numero finito di punti di discontinuità sono integrabili.
		
	\end{teorema}
	
	Consideriamo un esempio grafico:
	
	\begin{center}
		
		\begin{tabular}{cc}
			
			\begin{tikzpicture}[scale=0.8]
				
				\draw[->, thick] (-0.5, 0) -- (4, 0);
				\draw[->, thick] (0, -0.5) -- (0, 3);
				
				\draw[thick] (0.5, 1) to[out=45, in=225] (1.5, 2.5) to[out=45, in=180] (2.2, 1.2) to[out=0, in=135] (3.5, 2.8);
				
				\draw[dashed] (0.5, 0) node[below] {$a$} -- (0.5, 1);
				\draw[dashed] (1.5, 0) node[below] {$c$} -- (1.5, 2.5);
				\draw[dashed] (2.2, 0) node[below] {$d$} -- (2.2, 1.2);
				\draw[dashed] (3.5, 0) node[below] {$b$} -- (3.5, 2.8);
				
				\fill[blue, opacity=0.3] (0.5,0) -- (0.5,1) to[out=45, in=225] (1.5,2.5) -- (1.5,0) -- cycle;
				\fill[blue, opacity=0.3] (1.5,0) -- (1.5,2.5) to[out=45, in=180] (2.2,1.2) -- (2.2,0) -- cycle;
				\fill[blue, opacity=0.3] (2.2,0) -- (2.2,1.2) to[out=0, in=135] (3.5,2.8) -- (3.5,0) -- cycle;
				
			\end{tikzpicture}
			&
			\begin{tikzpicture}[scale=0.8]
				
				\draw[->, thick] (-0.5, 0) -- (4, 0);
				\draw[->, thick] (0, -0.5) -- (0, 3);
				
				\draw[thick] (0.5, 1) to[out=10, in=270] (1.8, 2.8);
				\draw[thick] (2.2, 1.5) -- (3.5, 1.5);
				
				\draw[dashed] (0.5, 0) node[below] {$a$} -- (0.5, 1);
				\draw[dashed] (1.8, 0) node[below] {$c$} -- (1.8, 2.8);
				\draw[dashed] (3.5, 0) node[below] {$b$} -- (3.5, 1.5);
				
			\end{tikzpicture}
			\\
			\parbox{5cm}{\centering È integrabile e l'integrale totale è la somma dell'integrale calcolato in ciascun intervallino.}
			&
			\parbox{5cm}{\centering Non è integrabile, in quanto non è limitata per colpa dell'asintoto verticale in $x = c$.}
			
		\end{tabular}
		
	\end{center}
	
	\section{Funzioni non integrabili}
	
	Finora abbiamo parlato delle funzioni che si possono integrare. Per quanto riguarda quelle non integrabili, abbiamo le funzioni non limitate. Tuttavia, per curiosità, vediamo una funzione limitata ma non integrabile:
	\[
		f(x) = \begin{cases}
		0 & \text{se } x \in [a, b] \cap \Q \\
		1 & \text{se } x \in [a, b] \setminus \Q
	\end{cases}
	\]
	
	Questa si chiama \textbf{funzione di Dirichlet}. Vediamo perché non è integrabile calcolando le somme integrali:
	\[
		\mathcal{S}^{P_n} = \sum_{k=1}^n \sup_{x \in [x_{k-1}, x_k]} (f(x)) (x_k - x_{k-1})
	\]
	\[
		\mathcal{S}_{P_n} = \sum_{k=1}^n \inf_{x \in [x_{k-1}, x_k]} (f(x)) (x_k - x_{k-1})
	\]
	
	Stavolta vediamo semplicemente come supremum e infimum sono rispettivamente $1$ e $0$. In questo modo le somme per difetto saranno $0$, e chiaramente sono diverse dalle somme per eccesso.
	
	\section{Conseguenze della monotonia}
	Vediamo ora un'altra conseguenza della monotonia. Sia $f$ integrabile in $[a, b]$, $f$ è dunque limitata. Formalmente possiamo scrivere:
	\[
		m \le f(x) \le M \quad \forall x \in [a, b]
	\]
	Siano, per esempio, $m = \inf(f(x))$, $M = \sup(f(x))$. Avendo racchiuso $f$ in quel modo, per monotonia abbiamo:
	\[
		\int_a^b m \, dx \le \int_a^b f(x) \, dx \le \int_a^b M \, dx \quad \text{cioè}
	\]
	\[
		m(b-a) \le \int_a^b f(x) \, dx \le M(b-a)
	\]
	Questa disuguaglianza è utile ad esempio per quando volessimo stimare il valore di un integrale che non sappiamo calcolare.
	
	Adesso proviamo a stimare il valore di un integrale che sappiamo calcolare, per vedere se il metodo funziona:
	\[
		\int_{0}^{1} \frac{4}{1 + x^2} \, dx
	\]
	
	\begin{itemize}
		
		\item Osserviamo che $f$ è continua nell'intervallo scelto. Dunque, per il teorema di Weierstrass, ammette minimo e massimo.
		
		\item Siccome $f$ è decrescente, il minimo è il valore corrispondente ad $x = 1$ e il massimo è il valore corrispondente ad $x = 0$.
		\[
			\min_{x \in [0, 1]} f(x) = f(1) = 2 \qquad \max_{x \in [0, 1]} f(x) = f(0) = 4
		\]
		
		\item Abbiamo dunque:
		\[
			\int_{0}^{1} 2 \, dx \leq \int_{0}^{1} \frac{4}{1 + x^2} \, dx \leq \int_{0}^{1} 4 \, dx \implies 2 \leq \int_{0}^{1} \frac{4}{1 + x^2} \, dx \leq 4
		\]
		
		\item Abbiamo stimato il nostro integrale. Se ora andassimo a calcolarlo, otterremmo che questo vale $\pi$, e dunque con:
		\[
			2 \leq \pi \leq 4
		\]
		avremmo un'approssimazione molto rozza di $\pi$.
		
	\end{itemize}
	
	\begin{teorema}{Teorema della media integrale}{}
		
		Sia $f$ continua in $[a, b]$. Allora $\exists x_0 \in [a, b]$ tale che:
		\[
			f(x_0) = \frac{\int_{a}^{b} f(x) \, dx}{b - a}
		\]
		
	\end{teorema}
	
	Il valore alla destra dell'uguale è detto \textbf{media integrale di $f$} e si indica con $\langle f \rangle$.
	Questo teorema è un corollario del teorema dei valori intermedi, in quanto $f$ è continua e si ha che:
	\[
		m(b - a) \leq \langle f \rangle \leq M(b - a)
	\]
	
	\begin{osservazione}{}{}
		
		Se $f$ è integrabile in $[a, b]$, allora è integrabile in $I \subseteq [a, b]$.
		
		Se $I$ è un generico intervallo $[x_0, x]$, con $x_0 \ge a$ fissato e $x$ variabile, purché rimanga $\le b$.
		La scrittura $\int_{x_0}^x f(t) \, dt$ ha dunque senso, e anzi ci permette di definire $G(x) := \int_{x_0}^x f(t) \, dt$ che viene infatti chiamata \textbf{funzione definita da integrale}.
		
	\end{osservazione}
	
	\begin{osservazione}{}{}
		
		Se $f \in C^1([a, b])$ e $F' = f$, allora $G(x) = F(b) - F(a)$.
	
	\end{osservazione}
	
	Viene spontaneo chiederci: cosa possiamo dire sulla funzione $G(x)$? Continua? Derivabile?
	
	Consideriamo come esempio la seguente funzione:
	\[
		G(x) = \int_0^x \operatorname{sgn}(t) \, dt, \text{ con } \operatorname{sgn}(t) = \begin{cases} 1 & \text{se } x > 0 \\ 0 & \text{se } x = 0 \\ -1 & \text{se } x < 1 \end{cases}
	\]
	
	Vediamo come si comporta $G$ per i vari valori di $x$:
	\begin{itemize}
		
		\item $x = 0 \colon G(0) = \int_0^0 \operatorname{sgn}(t) \, dt = 0$
		
		\item $x > 0 \colon G(x) = \int_0^x 1 \, dt = [t]_0^x = x$
		
		\item $x < 0 \colon G(x) = \int_0^x (-1) \, dt = -\int_x^0 (-1) \, dt = -[-t]_x^0 = -x$
		
	\end{itemize}
	
	Abbiamo dunque $G(x) = \begin{cases} x & \text{se } x > 0 \\ 0 & \text{se } x = 0 \\ -x & \text{se } x < 0 \end{cases} = |x|$
	
	Non è detto che sia una funzione "a caso", spesso descrivono funzioni che già conosciamo.
	
	In questo caso, $G$ è continua ma non derivabile.
	
	\section{Funzione integrale}
	
	Consideriamo una funzione $G(x)$ definita come:
	\[
		G(x) = \int_{0}^{x} f(t) \, dt
	\]
	dove $f(t)$ è definita a tratti come:
	\[
		f(t) = \begin{cases}
		t, & \text{se } 0 \leq t \leq 1 \\
		1, & \text{se } 1 < t \leq 2
		\end{cases}
	\]
	
	Vogliamo calcolare $G(x)$ nei diversi intervalli:
	\begin{itemize}
		
		\item Per $x \in [0, 1]$: $G(x) = \int_{0}^{x} t \, dt = \left[ \frac{t^2}{2} \right]_{0}^{x} = \frac{x^2}{2}$
		
		\item Per $x \in (1, 2]$: $G(x) = \int_{0}^{1} t \, dt + \int_{1}^{x} 1 \, dt = \dots = x - \frac{1}{2}$
		
	\end{itemize}
	
	Per la proprietà additiva dell'integrale, spezziamo l'intervallo di integrazione in quanto la funzione è definita a tratti: partiamo da $0$ ad $1$ (di cui conosciamo già la formula) e sommiamo il restante tratto da $1$ ad $x$.
	
	\begin{proposizione}{}{}
		
		Sia $f$ integrabile in $[a, b]$. Allora $G(x)$ è \textbf{lipschitziana} (dunque è \textbf{continua}).
	
	\end{proposizione}
	
	\begin{dimostrazione}{}{}
		
		Scriviamo la differenza $G(x) - G(y)$ sfruttando le proprietà degli integrali:
		\[
			G(x) - G(y) = \int_{x_0}^{x} f(t) \, dt - \int_{x_0}^{y} f(t) \, dt
		\]
		Cambiamo segno invertendo gli estremi del secondo integrale:
		\[
			= \int_{x_0}^{x} f(t) \, dt + \int_{y}^{x_0} f(t) \, dt = \int_{y}^{x} f(t) \, dt
		\]
		Passiamo ora al modulo sfruttando la monotonia dell'integrale:
		\[
			\left| \int_{y}^{x} f(t) \, dt \right| \leq \int_{y}^{x} |f(t)| \, dt
		\]
		Siccome $f$ è integrabile, allora è limitata (maggiorata superiormente da $M = \sup f$ in $[a, b]$). Dunque $|f(t)| \leq M$ e, per monotonia:
		\[
			\left| \int_{y}^{x} f(t) \, dt \right| \leq \int_{y}^{x} M \, dt = M(x - y)
		\]
		che è proprio la condizione necessaria affinché una funzione sia lipschitziana.
		
	\end{dimostrazione}
	
	\begin{proposizione}{}{}
		
		Se $f(t) \geq 0$, allora $G(x)$ è monotona crescente.
		
	\end{proposizione}
	
	\begin{proposizione}{}{}
		
		Se $f(t)$ è continua, allora $G(x)$ è derivabile. Inoltre si ha che $G' = f$.
	
	\end{proposizione}
	
	\begin{dimostrazione}{}{}
		
		Scriviamo il rapporto incrementale:
		\[
			(G(x+h) - G(x)) \cdot \frac{1}{h} = \frac{1}{h} \left( \int_{x_0}^{x+h} f(t) \, dt - \int_{x_0}^{x} f(t) \, dt \right) =
		\]
		\[
			= \frac{1}{h} \left( \int_{x_0}^{x+h} f(t) \, dt + \int_{x}^{x_0} f(t) \, dt \right) = \frac{1}{h} \left( \int_{x}^{x+h} f(t) \, dt \right) = \frac{\int_{x}^{x+h} f(t) \, dt}{h}
		\]
		Siccome $h$ è l'ampiezza della partizione questo calcolo è proprio la formula della media integrale.
		Dunque $\exists c \in [x, x+h]$ tale che $f(c) = \frac{\int_{x}^{x+h} f(t) \, dt}{h}$.
		
		Passiamo ora al limite: $\lim_{h \to 0^+} f(c)$.
		Siccome $c \in [x, x+h]$, man mano che $h \to 0$, $x+h \to x$ e di conseguenza anche $c \to x$.
		Il risultato del rapporto incrementale per $G(x)$ dunque è proprio $f(x)$.
		
	\end{dimostrazione}
	
	\begin{teorema}{Teorema fondamentale del calcolo integrale}{}
		
		Sia $f$ continua in $[a, b]$. $F$ è una primitiva di $f$ definita come $F' = f$ se e solo se:
		\[
		F(x) = \int_{x_0}^{x} f(t) \, dt + c
		\]
		con $x_0 \in [a, b]$, $c \in \mathbb{R}$.
		
	\end{teorema}
	
	Implementiamo il \textit{TFCI} per gli integrali definiti:
	\begin{itemize}
		
		\item Caso $x = b$: $\quad F(b) = \int_{x_0}^{b} f(t) \, dt + c$
		
		\item Caso $x = a$: $\quad F(a) = \int_{x_0}^{a} f(t) \, dt + c$
	
	\end{itemize}
	
	Allora:
	\[
		F(b) - F(a) = \int_{x_0}^{b} f(t) \, dt + c - \left(\int_{x_0}^{a} f(t) \, dt + c\right) = \int_{x_0}^{b} f(t) \, dt + \int_{a}^{x_0} f(t) \, dt = \int_{a}^{b} f(t) \, dt
	\]
	dove nel secondo passaggio abbiamo applicato il cambio degli estremi e nel terzo l'additività dell'integrale.
	
	Che è lo stesso risultato ottenuto per il teorema degli integrali definiti.
	Ciò significa che è rispettata l'ipotesi di lipschitzianità di $f$, dunque $F(x)$, la primitiva, è una funzione continua.
	
	\section{Integrale indefinito}
	
	\begin{definizione}{Integrale indefinito}{}
		
		L'insieme delle primitive di una funzione $f$ si chiama \textbf{integrale indefinito}:
		\[
			\int f(t) \, dt := \{F(t) + c, \, c \in \mathbb{R}\}
		\]
	
	\end{definizione}
	
	Valgono tutte le formule viste precedentemente, ma anziché usare la notazione $F(b) - F(a)$ scriviamo $F(x) + c$. Ad esempio:
	\[
		\int x^n \, dx = \frac{x^{n+1}}{n+1} + c
	\]
	
	Riportiamo di seguito alcuni integrali notevoli e casi particolari:
	\begin{itemize}
		
		\item $\int \frac{1}{\cos^2 x} \, dx = \operatorname{tg}(x) + c$
		
		\item $\int \frac{1}{\sin^2 x} \, dx = \operatorname{cotan}(x) + c$
		
		\item $\int \frac{1}{\sqrt{1-x^2}} \, dx = \arcsin(x) + c$
		
		\item $\int \frac{-1}{\sqrt{1-x^2}} \, dx = \arccos(x) + c$
		
		\item $\int a^x \, dx = \frac{a^x}{\ln a} + c$
		
		\item $\int e^{ax} \, dx = \frac{e^{ax}}{a} + c$
		
		\item $\int \cos(ax) \, dx = \frac{\sin(ax)}{a} + c$
		
		\item $\int \sin(ax) \, dx = -\frac{\cos(ax)}{a} + c$
	
	\end{itemize}
	
	\section{Formule di integrazione}
	
	\subsection{Integrazione per sostituzione}
	
	In alcuni degli esempi proposti precedentemente abbiamo visto come gli argomenti di alcune funzioni non fossero semplicemente $x$, bensì funzioni più o meno complesse: in questo caso usiamo la \textbf{formula di integrale di funzioni composte}.
	
	Questa si può applicare quando abbiamo un integrale fatto in questo modo:
	\[
		\int f(g(x)) \cdot g'(x) \, dx = F(g(x)) + c
	\]
	
	Questo è in realtà un caso molto particolare di una formula più generale chiamata \textbf{formula di integrale per sostituzione}.
	
	In generale abbiamo $\int f(g(x)) \cdot g'(x) \, dx = \int f(t) \, dt$, facendo la sostituzione $t = g(x)$.
	
	Come procediamo?
	\begin{itemize}
		
		\item Sostituisco $t = g(x)$
		
		\item Calcolo $dt = g'(x) \, dx$
		
		\item L'integrale diventa $\int f(t) \, dt = F(t) + c = F(g(x)) + c$
	
	\end{itemize}
	
	Se stiamo lavorando con gli integrali definiti, quando applichiamo la sostituzione cambiano anche gli estremi di integrazione:
	\[
		\int_{a}^{b} f(g(x)) \cdot g'(x) \, dx = \int_{g(a)}^{g(b)} f(t) \, dt
	\]
	con la sostituzione $t = g(x)$.
	
	\begin{proposizione}{}{}
		
		Sia $f$ continua in $[-a, a]$, con $a > 0$.
		
		\begin{itemize}
		
			\item Se $f$ è pari, allora $\int_{-a}^{a} f(x) \, dx = 2 \int_{0}^{a} f(x) \, dx$.
		
			\item Se $f$ è dispari, allora $\int_{-a}^{a} f(x) \, dx = 0$.
		
		\end{itemize}
		
	\end{proposizione}
	
	\subsection{Integrazione per parti}
	
	Vediamo adesso una formula importante che possiamo ricavare proprio dal Teorema Fondamentale del Calcolo Integrale.
	Siano $f, g$ derivabili in $[a, b]$. Ricordiamo dunque che $(f \cdot g)' = f' \cdot g + f \cdot g'$.
	Dunque, se integriamo $(f \cdot g)'$, otteniamo:
	\[
		\int_{a}^{b} (f \cdot g)' \, dx = \int_{a}^{b} (f' \cdot g + f \cdot g') \, dx
	\]
	\[
		[f(x) \cdot g(x)]_{a}^{b} = \int_{a}^{b} f' \cdot g \, dx + \int_{a}^{b} f \cdot g' \, dx
	\]
	Isolando quel termine, otteniamo la cosiddetta \textbf{formula di integrazione per parti}, che vale sia per gli integrali definiti che indefiniti.
	
	\begin{itemize}
		
		\item \textbf{Definiti}: $\int_{a}^{b} f'(x) \cdot g(x) \, dx = [f(x) \cdot g(x)]_{a}^{b} - \int_{a}^{b} f(x) \cdot g'(x) \, dx$
		
		\item \textbf{Indefiniti}: $\int f'(x) \cdot g(x) \, dx = f(x) \cdot g(x) - \int f(x) \cdot g'(x) \, dx$
		
	\end{itemize}
	
	\begin{osservazione}{Nota per gli esercizi}{}
		
		Se dovesse capitare di dover integrare per parti più volte di seguito, scegliere, se possibile, per la stessa funzione lo stesso "ruolo". Ad esempio entrambe le volte $f'$ oppure entrambe le volte $g$.
		
	\end{osservazione}
	
	\section{Formulario di integrazione}
	
	Presentiamo di seguito un utile formulario per affrontare alcune tipologie ricorrenti di integrali indefiniti, particolarmente utili quando si applicano i metodi di integrazione per parti o per sostituzione.
	
	Se abbiamo un integrale della forma:
	\begin{itemize}
		
		\item $\int P(x) \cdot \begin{cases} e^{ax} \\ \cos(ax) \\ \sin(ax) \end{cases} dx \longmapsto \int P(x) \cdot \begin{cases} (\frac{e^{ax}}{a})' \\ (\frac{\sin(ax)}{a})' \\ (\frac{-\cos(ax)}{a})' \end{cases} dx$
		
		\item $\int P(x) \cdot \begin{cases} \ln(x) \\ \operatorname{arctg}(x) \end{cases} dx \longmapsto \int Q'(x) \cdot \begin{cases} \ln(x) \\ \operatorname{arctg}(x) \end{cases} dx$
		
		\item $\int e^{ax} \cdot \begin{cases} \sin(bx) \\ \cos(bx) \end{cases} dx$.\\
		In questo caso riscriviamo il fattore esponenziale come $(e^{ax}/a)'$, dunque $g'$. Dovremo applicare due volte la formula di integrazione per parti, ottenendo un'equazione in cui l'integrale di partenza (chiamato $I$) ricompare al secondo membro:
		\[
			I = A - I \implies I = \frac{A}{2}
		\]
		
		\item $\int \begin{cases} \sin^k(x) \\ \cos^k(x) \end{cases} dx$\\
		Per il calcolo di tali integrali si procede distinguendo la parità dell'esponente $k$:
		\begin{itemize}
			
			\item $k$ pari: si procede per parti.
			
			\item $k$ dispari: si procede per sostituzione.
			
		\end{itemize}
		
	\end{itemize}
	
	\section{Integrali di funzioni razionali}
	Siano $P, Q$ polinomi, consideriamo l'integrale:
	\[
		\int \frac{P(x)}{Q(x)} \, dx
	\]
	Vediamo i vari casi in base ai gradi di $P$ e $Q$:
	\begin{description}
		
		\item $\deg P = 0,\, \deg Q = 1$:
		\[
			\int \frac{1}{ax+b} \, dx = \text{per sostituzione } t = ax+b = \frac{1}{a} \ln|ax+b| + c
		\]
		
		\item $\deg P = \deg Q = 1$:
		\[
			\int \frac{cx+d}{ax+b} \, dx = \text{vogliamo riscrivere la frazione come } A + \frac{B}{ax+b}
		\]
		In generale, la formula per trovare $A, B$ è:
		\[
			A(ax+b) + B = cx+d \implies Aax + Ab + B = cx+d \implies \begin{cases} Aa = c \\ Ab + B = d \end{cases}
		\]
		
		\item $\deg P = 0,\, \deg Q = 2$:
		\[
			\int \frac{1}{Ax^2 + Bx + C} \, dx = \frac{1}{A} \int \frac{1}{x^2 + \frac{B}{A}x + \frac{C}{A}} \, dx
		\]
		e andiamo a calcolarci le radici del denominatore.
		
		Per comodità rinominiamo $b = \frac{B}{A}$ e $c = \frac{C}{A}$
		Distinguiamo in casi in base al discriminante:
		
		\begin{itemize}
			
			\item \textbf{Caso $\Delta = 0$}\\
			Avremo $x_{1,2} = -\frac{b}{2}$, e possiamo riscrivere $x^2 + bx + c = \left(x + \frac{b}{2}\right)^2$.\\
			L'integrale diventa dunque:
			\[
				\int \frac{1}{\left(x + \frac{b}{2}\right)^2} \, dx = \int \left(x + \frac{b}{2}\right)^{-2} \, dx = -\frac{1}{x + \frac{b}{2}} + c
			\]
			
			\item \textbf{Caso $\Delta > 0$}\\
			In tal caso abbiamo due radici $x_1 \neq x_2$, dunque $x^2 + bx + c = (x - x_1)(x - x_2)$ e la funzione diventa:
			\[
				\frac{1}{(x - x_1)(x - x_2)} = \text{vogliamo scriverla come} = \frac{A}{x - x_1} + \frac{B}{x - x_2}
			\]
			Con $A, B$ ricavabili dalla seguente formula:
			\[
				\begin{cases}
					A + B = 0 \\
					-A - B = 1
				\end{cases}
				\quad \text{oppure} \quad
				A = \frac{1}{x_1 - x_2}, \, B = \frac{1}{x_2 - x_1}
			\]
			L'integrale diventa dunque:
			\[
				\int \left( \frac{A}{x - x_1} + \frac{B}{x - x_2} \right) dx = A \left( \ln\left|\frac{x - x_1}{x - x_2}\right| \right) + c
			\]
			
			\item \textbf{Caso $\Delta < 0$}\\
			Il denominatore non avrà radici, ma possiamo comunque scrivere $x^2 + bx + c = (x - \alpha)^2 + \beta^2$.
			\[
				\begin{aligned}
					\int \frac{1}{(x - \alpha)^2 + \beta^2} \, dx &= \substack{\text{sostituzione} \\ t = x - \alpha \\ dt = dx} = \int \frac{1}{t^2 + \beta^2} \, dt = \frac{1}{\beta^2} \int \frac{1}{\frac{t^2}{\beta^2} + 1} \, dt = \\
					&= \substack{\text{sostituisco} \\ u = t/\beta \\ u \cdot \beta = t \\ \beta \, du = dt} = \frac{1}{\beta^2} \int \frac{\beta}{u^2 + 1} \, du = \frac{1}{\beta} \operatorname{arctg}(u) + c = \\
					&= \frac{1}{\beta} \operatorname{arctg}(t/\beta) + c = \frac{1}{\beta} \operatorname{arctg}\left(\frac{x - \alpha}{\beta}\right) + c
				\end{aligned}
			\]
			
			Per quanto riguarda il caso $\Delta < 0$, purtroppo non esiste una formula fissa per $\alpha$ e $\beta$, ma serve il colpo d'occhio. Ecco alcuni esempi:
			
			\begin{itemize}
				
				\item $x^2 + 2x + 2 = (x^2 + 2x + 1) + 1 = (x + 1)^2 + 1$\\
				Qui abbiamo che $\alpha = -1, \beta = 1$.
				
				\item $x^2 + 4x + 5 = (x^2 + 4x + 4) + 1 = (x + 2)^2 + 1$\\
				Qui $\alpha = -2, \beta = 1$.
				
			\end{itemize}
			
		\end{itemize}
		
		\item $\deg P = 1,\, \deg Q = 2$:
		\[
			\int \frac{dx + e}{x^2 + bx + c} \, dx
		\]
		
		Anche stavolta dobbiamo distinguere in base al segno del discriminante $\Delta$:
		
		\begin{itemize}
			
			\item [Caso $\Delta > 0$:] $x^2 + bx + c = (x - x_1)(x - x_2)$, dunque vorremmo riscrivere il polinomio come:
			\[
				\int \frac{dx + e}{x^2 + bx + c} \, dx = \int \left( \frac{A}{x - x_1} + \frac{B}{x - x_2} \right) dx
			\]
			trovando $A$ e $B$ tramite il sistema:
			\[
				\begin{cases}
					A + B = d \\
					A x_2 + B x_1 = \dots
				\end{cases}
			\]
			portando alla soluzione $A \ln|x - x_1| + B \ln|x - x_2| + c$.
			
			\item [Caso $\Delta \le 0$:] Vogliamo riscrivere la funzione integranda come:
			\[
				\frac{dx + e}{x^2 + bx + c} = A \cdot \frac{2x + b}{x^2 + bx + c} + \frac{B}{x^2 + bx + c}
			\]
			Notiamo che al primo addendo avremo che il numeratore sarà la derivata del denominatore, permettendoci di ottenere un logaritmo come risultato. 
			Qui ricaviamo $A$ e $B$ dal sistema:
			\[
				\begin{cases}
					2A = d \\
					Ab + B = e
				\end{cases}
			\]
			e così possiamo scrivere la soluzione che sarà:
			\[
				\int \frac{dx + e}{x^2 + bx + c} \, dx = A \ln|x^2 + bx + c| + B \int \frac{1}{x^2 + bx + c} \, dx
			\]
			dove l'ultimo integrale lo sappiamo risolvere con le formule viste in precedenza.
			
		\end{itemize}
		
	\end{description}
	
	In generale, ci sono casi in cui potremmo avere $\deg P \ge \deg Q$, e in quel caso occorre scomporre oppure fare la divisione tra polinomi.
	Cerchiamo dunque di riscrivere:
	\[
		\frac{P(x)}{Q(x)} = \text{risultato} + \frac{\text{resto}}{Q(x)}
	\]
	
	Vediamo un esempio con il seguente integrale: $\int \frac{x^3}{x^2 + 2x + 2} \, dx$. Eseguendo la divisione tra polinomi:
	
	\begin{center}
		
		\begin{tabular}{r|l}
			
			$x^3 + 0x^2 + 0x + 0$ & $x^2 + 2x + 2$ \\
			\cline{2-2}
			$- x^3 - 2x^2 - 2x +0$ & $\multirow{2}{*}{x - 2}$ \\
			\cline{1-1}
			// $\;\; -2x^2 - 2x \phantom{+ 0}$ & \\
			$+ 2x^2 + 4x + 4$ & \\
			\cline{1-1}
			// $\;\;\quad 2x + 4$ & 
			
		\end{tabular}
	\end{center}
	
	Dunque possiamo riscrivere l'integrale come:
	\[
		\int \frac{x^3}{x^2 + 2x + 2} \, dx = \int \left( x - 2 + \frac{2x + 4}{x^2 + 2x + 2} \right) dx
	\]
	che risolviamo poi con i metodi che già conosciamo.
	
	\chapter{Integrali impropri}
	
	Negli integrali secondo Riemann abbiamo visto come, per ipotesi, $f$ dovesse essere continua su $[a, b]$, un intervallo chiuso e limitato.
	
	Ma che succede nel caso $f$ non fosse limitata in quell'intervallo?
	Un esempio può essere $f(x) = \frac{1}{x}$ in $[-1, 0)$ ad esempio: dato che in $x_0 = 0$ la funzione non è definita, l'intervallo rimane aperto.
	
	In casi come questo si parla di \textbf{integrali impropri di seconda specie}.
	
	\begin{definizione}{Integrale improprio di seconda specie}{}
		
		Sia $f: [a, x_0) \to \R$ tale che:
		\begin{enumerate}
			
			\item $f$ è continua in $[a, x_0)$
			
			\item $\lim_{x \to x_0^-} f(x) = \pm\infty$
			
		\end{enumerate}
		
	\end{definizione}
	
	Per la proprietà (1) allora $f$ è integrabile in senso proprio (cioè "normalmente" come abbiamo visto finora) in qualsiasi intervallo del tipo $[a, x_0 - \epsilon]$, con $\epsilon > 0$.
	
	Se dunque calcoliamo $\lim_{\epsilon \to 0^+} \int_{a}^{x_0 - \epsilon} f(x) \, dx$ e questo limite risulta in un numero finito $l \in \R$, allora diciamo che $f$ è integrabile in senso improprio in $[a, x_0)$ e:
	\[
		\int_{a}^{x_0} f(x) \, dx = l
	\]
	
	Il discorso è analogo per intervalli del tipo $(x_0, b]$.
	
	Se $\lim_{\epsilon \to 0^+} \int_{x_0 + \epsilon}^{b} f(x) \, dx = l \in \R$ allora $\int_{x_0}^{b} f(x) \, dx = l$.
	
	Consideriamo la funzione $f(x) = \frac{1}{x^\alpha}$, con $\alpha > 0$ nell'intervallo $(0, 1]$. 
	Ci chiediamo: siccome $f$ non è limitata in $x_0 = 0$, è integrabile in senso improprio in $(0, 1]$?
	
	Siccome siamo nel caso $(x_0, b]$, dobbiamo studiare l'integrale:
	\[
		\int_{x_0 + \varepsilon}^{b} f(x) \, dx = \int_{\varepsilon}^{1} \frac{1}{x^\alpha} \, dx = 
		\begin{cases}
			\alpha = 1 \implies [\ln|x|]_\varepsilon^1 = -\ln\varepsilon \\
			\alpha \neq 1 \implies \left[ \frac{x^{-\alpha + 1}}{-\alpha + 1} \right]_\varepsilon^1 = \frac{1 - \varepsilon^{1-\alpha}}{1 - \alpha}
		\end{cases}
	\]
	
	Calcoliamo ora il limite per $\varepsilon \to 0^+$:
	\[
		\lim_{\varepsilon \to 0^+} \int_{\varepsilon}^{1} \frac{1}{x^\alpha} \, dx = 
		\begin{cases}
			\alpha = 1 \implies \lim_{\varepsilon \to 0^+} -\ln\varepsilon = +\infty \\
			\alpha \neq 1 \implies 
			\begin{cases}
				\text{se } 1 - \alpha > 0 \implies \alpha < 1 : \frac{1}{1-\alpha} \\
				\text{se } 1 - \alpha < 0 \implies \alpha > 1 : +\infty
			\end{cases}
		\end{cases}
	\]
	
	Possiamo dunque dire che $\frac{1}{x^\alpha}$ in $(0, 1]$ è integrabile in senso improprio se e solo se $\alpha < 1$.
	Analogamente alle serie, possiamo dire che $f$ converge solo per $\alpha < 1$ e diverge se $\alpha \ge 1$.
	
	In generale, se ci troviamo a calcolare $\int_{0}^{A} \frac{1}{x^\alpha} \, dx$ ci conviene spezzarlo in:
	\[
		\int_{0}^{1} \frac{1}{x^\alpha} \, dx + \int_{1}^{A} \frac{1}{x^\alpha} \, dx
	\]
	che sappiamo entrambi come trattare.
	Il discorso vale in generale per tutte le funzioni del tipo $f(x) = \frac{1}{(x_0 - x)^\alpha}$ negli intervalli $[a, x_0)$ dove dobbiamo studiare $\lim_{\varepsilon \to 0^+} \int_{a}^{x_0 - \varepsilon} \frac{1}{(x_0 - x)^\alpha} \, dx$ che esce finito se e solo se $\alpha < 1$.
	
	Come per le serie, possiamo definire:
	
	\begin{itemize}
		
		\item \textbf{Funzioni a segno costante}: se $f(x) \geq 0$ nell'intervallo considerato, allora l'integrale improprio di $f$ o è finito, o vale $+\infty$.
		
		\item \textbf{Criterio del confronto}: siano $f, g$ continue in $[a, x_0)$, t.c. $0 \leq f \leq g$. Allora $0 \leq \int_{a}^{x_0} f(x) \, dx \leq \int_{a}^{x_0} g(x) \, dx$ e dunque:
		\begin{itemize}
			
			\item Se $g$ è integrabile i.s.i. (in senso improprio), allora lo è anche $f$.
			
			\item Se l'integrale improprio di $f = +\infty$, allora anche quello di $g$.
			
		\end{itemize}
		
		Ad esempio, sia $f(x) = \frac{\sqrt{7-x}}{\sqrt{3-x}}$ in $[1, 3)$. Vogliamo maggiorarla con una funzione: possiamo sceglierne una con lo stesso denominatore ma con numeratore maggiore:
		\[
			1 \leq x < 3 \rightsquigarrow -3 < -x \leq -1 \rightsquigarrow 4 < 7-x \leq 6
		\]
		\[
			\frac{\sqrt{7-x}}{\sqrt{3-x}} \leq \frac{\sqrt{6}}{\sqrt{3-x}} \rightsquigarrow \text{questa è una versione generalizzata di } \frac{1}{x^\alpha} \text{ con } \alpha = \frac{1}{2} < 1
		\]
		Siccome $g$ è integrabile i.s.i., per confronto lo è anche $f$.
		
		\item \textbf{Criterio del confronto asintotico}: sia $f$ continua in $[a, x_0)$, $f \geq 0$. Se $\lim_{x \to x_0^-} (x_0 - x)^\alpha \cdot f(x) = l \in (0, +\infty)$, allora significa che asintoticamente $f$ si comporta come $\frac{1}{(x_0 - x)^\alpha}$, dunque:
		\begin{itemize}
			
			\item $\alpha < 1$: integrabile i.s.i.
			
			\item $\alpha \geq 1$: non integrabile i.s.i.
		
		\end{itemize}
		
		Ad esempio, $\int_{0}^{\pi/2} \frac{\sqrt{x}}{\sin x} \, dx$. Scegliamo $\alpha = 1/2$ in modo da ottenere il limite notevole $\frac{x}{\sin x} = 1$. Dunque, dato che $\alpha = 1/2$, la funzione è integrabile i.s.i.
		
	\end{itemize}
	
	E se invece avessimo che uno dei due estremi di integrazione è $\pm\infty$?
	In tal caso si parlerebbe di \textbf{integrali impropri di prima specie}, e ovviamente il discorso vale per $f$ continue in $[a, +\infty)$ oppure $(-\infty, b]$.
	
	\begin{definizione}{Integrale impropri di prima specie}{}
		
		Se $\exists \lim_{c \to +\infty} \int_{a}^{c} f(x) \, dx \in \R$ si dice che $f$ è integrabile in senso improprio in $[a, +\infty)$ e vale che:
		\[
		\int_{a}^{+\infty} f(x) \, dx = \lim_{c \to +\infty} \int_{a}^{c} f(x) \, dx = l
		\]
		
	\end{definizione}
	
	\begin{osservazione}{}{}
		Calcoliamo il seguente integrale improprio noto:
		\[
			\int_{1}^{+\infty} \frac{1}{x^\alpha} \, dx = \lim_{c \to +\infty} \int_{1}^{c} \frac{1}{x^\alpha} \, dx = 
			\begin{cases}
				\alpha = 1 \implies +\infty \\
				\alpha > 1 \implies \frac{1}{1-\alpha}\\
				\alpha < 1 \implies +\infty
			\end{cases}
		\]
		
	\end{osservazione}
	
	Anche per gli integrali impropri di prima specie valgono, seppur con qualche accorgimento, il criterio del confronto e confronto asintotico visti precedentemente.
	
	\begin{itemize}
		
		\item \textbf{Criterio del confronto:} siano $f, g$ continue in $[a, +\infty)$, tali che $0 \le f \le g$.
		Allora $0 \le \int_{a}^{+\infty} f(x) \, dx \le \int_{a}^{+\infty} g(x) \, dx$ e dunque:
		\begin{itemize}
			
			\item se $\int_{a}^{+\infty} g(x) \, dx < +\infty$ allora $\int_{a}^{+\infty} f(x) \, dx < +\infty$
			
			\item se $\int_{a}^{+\infty} f(x) \, dx = +\infty$ allora $\int_{a}^{+\infty} g(x) \, dx = +\infty$
			
		\end{itemize}
		
		\begin{osservazione}{}{}
			
			Prendiamo ad esempio l'integrale $\int_{1}^{+\infty} \frac{\arctan x}{x^2} \, dx$. 
			Siccome $\arctan x$ è limitata e il suo massimo vale $\pi/2$, possiamo maggiorare la funzione:
			\[
			\int_{1}^{+\infty} \frac{\arctan x}{x^2} \, dx \le \int_{1}^{+\infty} \frac{\pi/2}{x^2} \, dx
			\]
			che è una versione generalizzata di $\frac{1}{x^\alpha}$ con $\alpha > 1$, dunque converge.
			
			Consideriamo invece $\int_{1}^{+\infty} \frac{\arctan x}{x} \, dx$. Analogamente possiamo scrivere:
			\[
			\int_{1}^{+\infty} \frac{\arctan x}{x} \, dx \le \int_{1}^{+\infty} \frac{\pi/2}{x} \, dx
			\]
			che però stavolta diverge.
			
		\end{osservazione}
		
		\item \textbf{Criterio del confronto asintotico:} sia $f$ continua in $[a, +\infty)$ con $f \ge 0$.
		Se $\lim_{x \to +\infty} (x)^\alpha \cdot f(x) = l \in (0, +\infty)$ allora $f(x) \sim \frac{1}{x^\alpha}$ e dunque:
		\begin{itemize}
			\item Se $\alpha > 1$: $f$ è integrabile in senso improprio.
			\item Se $\alpha \le 1$: $f$ non è integrabile in $[a, +\infty)$.
		\end{itemize}
		
		Per gli esercizi, come per le serie, scegliamo i termini di grado maggiore a numeratore e denominatore.
		
		\begin{osservazione}{}{}
			Consideriamo l'integrale improprio:
			\[
			\int_{1}^{+\infty} \frac{x + \cos x}{x^3 + \sin x} \, dx
			\]
			Siccome $\cos$ e $\sin$ sono limitate, sono asintoticamente trascurabili, dunque la funzione si comporta come $\frac{x}{x^3} = \frac{1}{x^2}$.
			Essendo $\alpha = 2 > 1$, la funzione è integrabile in senso improprio.
		\end{osservazione}
		
	\end{itemize}
	
	\chapter{Equazioni differenziali}
	
	Un'equazione differenziale è un'equazione dove:
	\begin{itemize}
		
		\item L'incognita è una funzione, che indichiamo come $y(t): I \subseteq \R \to \R$.
		
		\item Il grado dell'equazione è rappresentato dall'ordine più alto di derivata in cui appare $y(t)$.
		
	\end{itemize}
	Ad esempio, $y'(t) = \cos(t)$ è un'equazione differenziale di primo grado, mentre $y''(t) = y'(t)$ è di secondo grado.
	
	Le equazioni differenziali nella forma
	\[
		y^{(n)}(t) = f(t, y(t), y'(t), \dots, y^{(n-1)}(t))
	\]
	si dicono \textbf{in forma normale}.
	
	Nonostante possa sembrare strano a dirsi, gli esempi visti prima sono entrambi in forma normale: per $y' = \cos(t)$ abbiamo $n=1$, dunque $f$ ha come variabile solo $t$. Un'equazione come $[y'(t)]^2 = \sin(y(t)) + \cos(y(t)) - t^2 + 1$ \textbf{non} è invece in forma normale, dato che compare un elevamento a potenza.
	
	\section{Equazioni differenziali semplici}
	Sono la categoria di equazioni differenziali più semplice perché si presentano nella forma:
	\[
		y'(t) = f(t)
	\]
	Sostanzialmente ci viene posta una domanda a cui sappiamo già rispondere: qual è quella funzione la cui derivata è $f(t)$?
	Semplice! È come calcolare una primitiva di $f$, dunque le soluzioni dell'equazione saranno del tipo:
	\[
		y(t) = \int f(t) \, dt
	\]
	\begin{osservazione}{}{}
		
		In casi come questo abbiamo infinite soluzioni, dato che l'insieme delle primitive è infinito.
		
	\end{osservazione}
	
	Se invece volessimo un'unicità, cioè una sola soluzione, dovremmo impostare l'equazione in questo modo:
	\[
		\begin{cases}
			y'(t) = g(t)\\
			y(t_0) = y_0
		\end{cases}
	\]
	che viene detta \textbf{condizione iniziale}. Tale sistema è detto \textbf{problema di Cauchy}.
	
	Vogliamo ora definire un altro prototipo di equazione differenziale, cioè di funzioni $y(t)$ che aumentano o diminuiscono in modo proporzionale a se stesse in un certo intervallo di tempo $\Delta t$. Formalmente:
	\[
		y(t + \Delta t) - y(t) = \lambda \cdot y(t) \, \Delta t
	\]
	Dividendo tutto per $\Delta t$:
	\[
		\frac{y(t + \Delta t) - y(t)}{\Delta t} = \lambda \cdot y(t)
	\]
	Se consideriamo una variazione temporale molto piccola, dunque $\Delta t \to 0$, avremo:
	\[
		\lim_{\Delta t \to 0} \frac{y(t + \Delta t) - y(t)}{\Delta t} = \lambda \cdot y(t) \implies y'(t) = \lambda \cdot y(t)
	\]
	in quanto il limite del rapporto incrementale definisce proprio $y'$.
	
	Ragionando, è facilmente intuibile come la funzione per cui accade questo è $e^{\lambda t}$, dunque le soluzioni saranno del tipo:
	\[
		y(t) = K \cdot e^{\lambda t}
	\]
	
	Vediamo un esempio pratico per entrare meglio nel meccanismo:
	\[
		\begin{cases}
			y'(t) = 5y(t) \\
			y(1) = 2
		\end{cases}
	\]
	Questa equazione è della forma $y' = \lambda y$, dunque le soluzioni sono della forma $y(t) = K e^{\lambda t}$, con $K$ da determinare per ottenere una soluzione unica.
	
	Sostituendo i dati noti, otteniamo $y(t) = K \cdot e^{5t}$, da cui imponendo la condizione iniziale $y(1) = K e^5 = 2$, ricaviamo $K = 2e^{-5}$.
	L'unica soluzione di questo problema di Cauchy è quindi $y(t) = 2e^{-5} \cdot e^{5t} = 2e^{5(t-1)}$.
	
	Una tipologia analoga a questa sono le equazioni nella forma:
	\[
		y'(t) = g(t) \cdot y(t)
	\]
	Le cui soluzioni possono essere scritte nella forma:
	\[
		y(t) = K \cdot e^{G(t)}, \quad \text{con } G(t) = \int g(t) \, dt
	\]
	
	Vediamo un esempio di questa tipologia:
	\[
		\begin{cases}
			y'(t) = t \cdot y(t) \\
			y(0) = 1
		\end{cases}
	\]
	Riconosciamo subito la forma $y' = g \cdot y$, dunque la soluzione sarà $y = e^G$, con $G$ primitiva di $g$.
	
	Calcoliamo la primitiva:
	\[
		G(t) = \int t \, dt = \frac{t^2}{2}
	\]
	L'unica soluzione del problema è pertanto $y(t) = e^{\frac{t^2}{2}}$.
	
	\section{Equazioni differenziali lineari}
	
	Finora abbiamo visto equazioni del tipo $y'(t) = \lambda(t) \cdot y(t)$, che si dicono \textbf{equazioni lineari omogenee}.
	Se avessimo invece un'equazione come $y'(t) = \lambda(t)y(t) + b(t)$ questa sarebbe un'\textbf{equazione lineare non omogenea}.
	
	Le non omogenee hanno una soluzione del tipo:
	\[
		y(t) = K \cdot e^{\Lambda(t)} + \bar{y}(t)
	\]
	dove $\bar{y}(t)$ si dice \textit{soluzione particolare}.
	
	\subsection{Metodo di variazione delle costanti}
	
	Per trovare concretamente $\bar{y}(t)$, utilizziamo il \textbf{metodo di variazione delle costanti}, cioè cercando un $\bar{y}(t) = B(t) \cdot e^{\Lambda(t)}$.
	In seguito a numerosi calcoli, otteniamo che:
	\[
		B(t) = \int b(t) \cdot e^{-\Lambda(t)} \, dt
	\]
	
	Se la soluzione particolare era $\bar{y}(t) = B(t) \cdot e^{\Lambda(t)}$ e $B(t)$ è quanto visto sopra, abbiamo che la soluzione $y(t)$ è la seguente:
	\[
		y(t) = K \cdot e^{\Lambda(t)} + B(t) \cdot e^{\Lambda(t)} = e^{\Lambda(t)}[K + B(t)]
	\]
	
	Vediamo di seguito un esempio applicativo per comprendere il metodo.
	
	\begin{osservazione}{}{}
		
		Data l'equazione con condizione iniziale:
		\[
			\begin{cases}
				y'(t) = \frac{1}{t} y(t) + t^2 \\
				y(1) = 2
			\end{cases}
			\quad \implies \quad
			\begin{aligned}
				\lambda(t) &= \frac{1}{t} \\
				b(t) &= t^2
			\end{aligned}
		\]
		
		Procediamo con i passi risolutivi:
		\begin{enumerate}
			
			\item Troviamo $\Lambda(t) = \int \lambda(t) \, dt = \int \frac{dt}{t} = \ln(t)$.
			
			\item Calcoliamo $e^{\Lambda(t)} = e^{\ln(t)} = t$.
			
			\item Troviamo $B(t) = \int b(t) \cdot e^{-\Lambda(t)} \, dt = \int t^2 \cdot \frac{1}{t} \, dt = \frac{t^2}{2}$.
			
			\item Mettiamo insieme i pezzi per la soluzione generica:
			\[
			y(t) = t \left[ K + \frac{t^2}{2} \right]
			\]
			
			\item Applichiamo la condizione iniziale per trovare $K$ e dunque l'unica soluzione:
			\[
				y(1) = 1 \left( K + \frac{1}{2} \right) = 2 \quad \implies \quad K = \frac{3}{2}
			\]
			
		\end{enumerate}
		
	\end{osservazione}
	
	\subsection{Metodo di somiglianza}
	Un altro metodo per trovare $y(t)$ senza calcolare integrali è detto \textbf{metodo di somiglianza}.
	Questo si applica ad un'equazione lineare non omogenea che ha:
	\begin{itemize}
		
		\item $\lambda(t) \equiv \lambda$, cioè una costante.
		
		\item $b(t)$ è un polinomio, un'esponenziale, seno oppure coseno.
		
	\end{itemize}
	
	Abbiamo tre casi principali. Vediamo il primo:
	\begin{description}
		
		\item [Caso 1:] $\lambda \neq 0$ e $b(t)$ un polinomio di grado $n$. Cerchiamo una soluzione particolare $\bar{y}(t)$ che sia un polinomio di grado $n$.
		Vediamo un esempio applicativo per comprendere il procedimento.
		\begin{osservazione}{}{}
			
			Data l'equazione con condizione iniziale:
			\[
			\begin{cases}
				y' = 3y + t \\
				y(0) = 1
			\end{cases}
			\]
			Notiamo che $\lambda(t) = 3$ e $b(t) = t$, che risulta essere un polinomio di primo grado.
			Cerchiamo dunque una soluzione particolare del tipo $\bar{y}(t) = c_1 t + c_2$, anch'essa un polinomio di primo grado.
			
			Per trovare le costanti $c_1$ e $c_2$, essendo $\bar{y}(t)$ una soluzione particolare, possiamo sostituirla ad ogni occorrenza di $y$ nell'equazione originale.
			Calcoliamo la derivata: $\bar{y}'(t) = c_1$. Sostituiamo $\bar{y}'(t)$ a $y'(t)$ e $\bar{y}(t)$ a $y(t)$ nell'equazione:
			\[
			c_1 = 3(c_1 t + c_2) + t \implies 3c_1 t + 3c_2 - c_1 + t = 0
			\]
			Raccogliendo i termini in $t$ e i termini noti otteniamo:
			\[
			t(3c_1 + 1) + 3c_2 - c_1 = 0
			\]
			Per il principio di identità dei polinomi, affinché questa equazione sia sempre verificata, i coefficienti di $t$ e il termine noto devono essere contemporaneamente nulli. Impostiamo ora il sistema:
			\[
			\begin{cases}
				3c_1 + 1 = 0 \\
				3c_2 - c_1 = 0
			\end{cases}
			\]
			Risolvendo il sistema otteniamo i valori delle costanti: $c_1 = -\frac{1}{3}$ e $c_2 = -\frac{1}{9}$.
			
			Con questi valori la soluzione particolare è $\bar{y}(t) = -\frac{1}{3}t - \frac{1}{9}$, e possiamo scrivere la soluzione generale dell'equazione differenziale come:
			\[
			y(t) = K e^{\lambda(t)} + \bar{y}(t) = K e^{3t} - \frac{1}{3}t - \frac{1}{9}
			\]
			Non resta che imporre la condizione iniziale $y(0) = 1$ per trovare il valore della costante $K$ come al solito.
			
		\end{osservazione}
		
		\item[\textbf{Caso 2:}] Si verifica quando $b(t) = L \cdot e^{mt}$, dunque $y'(t) = \lambda \cdot y(t) + L \cdot e^{mt}$.
		Se $\lambda \neq m$, $\bar{y}(t) = c_1 \cdot e^{mt}$. La procedura per trovare il coefficiente $c_1$ è analoga a prima.
		
		\item[\textbf{Caso 3:}] Si verifica se $b(t) = L \cdot \cos(\alpha t)$ oppure $L \cdot \sin(\alpha t)$ e dunque l'equazione diventa $y'(t) = \lambda y(t) + L \cdot \cos/\sin(\alpha t)$.
		In tal caso $\bar{y}(t) = c_1 \cdot \cos(\alpha t) + c_2 \cdot \sin(\alpha t)$. Procedura analoga.
		
	\end{description}
	
	\newpage
	\section{Equazioni lineari del secondo ordine}
	Sono equazioni nella seguente forma:
	\[
		y''(t) + a(t)y'(t) + b(t)y(t) = g(t)
	\]
	Se $g(t) \equiv 0$ si dice omogenea.
	Se $a, b, g \in C(I)$
	
	\begin{description}
		
		\item [Caso omogeneo:] $g(t) \equiv 0$. In questo caso si trovano $y_1$ e $y_2$ che soddisfano l'equazione omogenea e non sono proporzionali (ad esempio non dobbiamo avere $y_1 = \alpha y_2$).
		
		Tutte le soluzioni dell'equazione omogenea possono essere espresse come combinazione lineare delle due soluzioni:
		\[
			y(t) = c_1 y_1(t) + c_2 y_2(t)
		\]
		
		\item[Caso non omogeneo] $g(t) \neq 0$: Bisogna cercare $\bar{y}(t)$ soluzione particolare dell'equazione completa.
		In questo caso tutte le soluzioni sono del tipo:
		\[
			y(t) = c_1 y_1(t) + c_2 y_2(t) + \bar{y}(t)
		\]
		Dove $y_1$ e $y_2$ sono le stesse che risolverebbero l'equazione omogenea.
		
	\end{description}
	
	Vogliamo adesso verificare se, come per le EDO del primo ordine, esiste un modo algoritmico per determinare le due soluzioni grazie alle quali potremo trovare tutte le altre.
	
	Per ora ci limitiamo solo allo studio delle omogenee, dunque nella forma
	\[
		y'' + a(t)y' + b(t)y = 0
	\]
	
	Osserviamo un parallelismo con le EDO del primo ordine:
	\begin{itemize}
		
		\item Se in una EDO del primo ordine della forma $ay' + by = 0$ poniamo $y = e^{\lambda t}$ come possibile soluzione, abbiamo che, sostituendo $e^{\lambda t}$ nell'equazione:
		\[
			a \cdot \lambda e^{\lambda t} + b \cdot e^{\lambda t} = 0 \implies e^{\lambda t}(a\lambda + b) = 0
		\]
		e per forza di cose abbiamo che $\lambda = -b/a$ è soluzione dell'equazione di primo grado tra parentesi.
		
		\item Se avessimo invece una EDO del secondo ordine $ay'' + by' + cy = 0$ e provassimo a fare la stessa cosa con l'esponenziale, otterremmo:
		\[
			a \cdot \lambda^2 \cdot e^{\lambda t} + b \cdot \lambda \cdot e^{\lambda t} + c \cdot e^{\lambda t} = e^{\lambda t}(a\lambda^2 + b\lambda + c) = 0
		\]
		essendo l'equazione risultante, detta equazione caratteristica, di secondo grado, non possiamo garantire solo soluzioni reali come nel caso delle EDO del primo ordine.
		
	\end{itemize}
	
	È dunque conveniente ricordare che, se $\lambda \in \mathbb{C}$, allora $(e^{\lambda t})' = \lambda e^{\lambda t}$, dunque funziona esattamente come per le variabili reali.
	
	Ci si è convinti facilmente che conviene definire come soluzione particolare una funzione $z(t) \colon I \subseteq \mathbb{R} \to \mathbb{C}$. Questa, come tutte le funzioni complesse, è strutturata in modo che anch'essa è data da una parte reale e una immaginaria:
	\[
		z(t) = z_1(t) + i \cdot z_2(t)
	\]
	
	Dunque, se sostituiamo questa funzione e le sue derivate alla generica EDO del secondo ordine omogenea $y'' + ay' + by = 0$, otteniamo:
	\[
		z'' + az' + bz = z''_1 + i \cdot z''_2 + a(z'_1 + i \cdot z'_2) + b(z_1 + i \cdot z_2) = 0
	\]
	separando i termini con $i$ e senza $i$:
	\[
		(z''_1 + az'_1 + bz_1) + i(z''_2 + az'_2 + bz_2) = 0
	\]
	
	Sappiamo invece risolvere quelle in cui $a(t) \equiv a$ e $b(t) \equiv b$, cioè sono entrambe costanti.
	Se poniamo come soluzione particolare $z(t) = e^{\lambda t}$ con $\lambda \in \C$ e sostituiamo, otteniamo che:
	\[
		y'' + ay + b = 0 \implies \lambda^2 e^{\lambda t} + a\lambda e^{\lambda t} + b e^{\lambda t} = 0 \implies e^{\lambda t}(\lambda^2 + a\lambda + b) = 0
	\]
	dove $\lambda^2 + a\lambda + b = 0$ viene detta \textbf{equazione caratteristica} dell'equazione differenziale.
	
	Si tratta ora di risolvere una banale equazione di secondo grado, dove se $\Delta \ge 0$ allora $\lambda_{1,2} \in \R$, mentre se $\Delta < 0$ allora $\lambda_{1,2} \in \C$.
	
	Per trovare algoritmicamente delle soluzioni, dobbiamo distinguere in casi in base al valore di $\Delta$:
	\begin{description}
		
		\item [Caso $\Delta > 0$:] $\implies \lambda_1 \neq \lambda_2$, scegliamo dunque come soluzioni distinte:
		\[
			y_1 = e^{\lambda_1 t}, \quad y_2 = e^{\lambda_2 t} \xrightarrow{\text{per il teorema delle soluzioni}} y(t) = c_1 \cdot e^{\lambda_1 t} + c_2 \cdot e^{\lambda_2 t}
		\]
		
		\item [Caso $\Delta = 0$] $\implies \lambda_1 = \lambda_2 = -a/2$, e dato che non possiamo scegliere due soluzioni uguali dato che devono essere linearmente indipendenti, prendiamo:
		\[
			y_1 = e^{\lambda_1 t}, \quad y_2 = t \cdot e^{\lambda_1 t} \implies y(t) = c_1 \cdot e^{\lambda_1 t} + c_2 \cdot t \cdot e^{\lambda_1 t}
		\]
		
		\item [Caso $\Delta < 0$] $\implies \lambda_1 \neq \lambda_2$ ma $\lambda_{1,2} \in \C$, dunque possiamo scriverli come $\lambda_{1,2} = p + i \cdot q$.
		Proprio per questa ragione sappiamo che:
		\[
			p = -a/2, \quad q = \frac{\sqrt{4b - a^2}}{2}
		\]
		per la formula di Eulero sappiamo che $e^{\lambda_1 t} = e^{pt + iqt} = e^{pt} \cdot \cos(qt) + i \cdot e^{pt} \cdot \sin(qt)$ e abbiamo che $y_1 = e^{pt}\cos(qt)$, $y_2 = e^{pt}\sin(qt)$, e dunque:
		\[
			y(t) = e^{pt} [c_1 \cdot \cos(qt) + c_2 \cdot \sin(qt)]
		\]
		
	\end{description}
	
\end{document}